\documentclass[12pt, a4paper]{book}

% الحزم الأساسية
\usepackage{polyglossia}
\setmainlanguage[numerals=maghrib]{arabic}
\setotherlanguage{english}

% إعدادات الخطوط (يفضل خط أميري)
\newfontfamily\arabicfont[Script=Arabic]{Amiri}
\newfontfamily\quranfont[Script=Arabic]{Scheherazade} % لخط القرآن إن وجد

% حزم التنسيق
\usepackage[margin=2.5cm]{geometry}
\usepackage{titlesec}
\usepackage{fancyhdr}
\usepackage{tocloft}
\usepackage{amsmath}
\usepackage{setspace} % لضبط المسافات بين السطور
\usepackage{parskip}  % لضبط المسافات بين الفقرات
\usepackage{imakeidx} % حزمة الفهارس العلمية

% إعدادات الفهارس
\makeindex[name=quran, title=فهرس الآيات القرآنية, columns=2]
\makeindex[name=hadith, title=فهرس الأحاديث النبوية, columns=1]
\makeindex[name=general, title=فهرس الأعلام, columns=2]

% أوامر مخصصة للتحقيق والتنسيق العلمي
\newcommand{\matn}[1]{{\arabicfont\bfseries «#1»}} % لتمييز متن الرسالة
\newcommand{\aya}[2]{﴿#1﴾\footnote{سورة #2.}\index[quran]{#2}} % للآيات
\newcommand{\hadith}[2]{«#1»\footnote{#2}\index[hadith]{#1}} % للأحاديث
\newcommand{\tarjama}[2]{#1\footnote{ترجمة: #2}\index[general]{#1}} % للأعلام

\usepackage{hyperref}

% إخفاء الأرقام التلقائية من العناوين (الفصول والأقسام) والاعتماد على المحتوى فقط
\setcounter{secnumdepth}{-1}
% تحديد مستوى العناوين التي تظهر في الفهرس
\setcounter{tocdepth}{2}

% إعدادات الروابط
\hypersetup{
    colorlinks=true,
    linkcolor=blue,
    filecolor=magenta,      
    urlcolor=cyan,
    pdftitle={شرح كتاب الرسالة للإمام الشافعي},
    pdfpagemode=FullScreen,
}

\begin{document}

% تفعيل مسافة ونصف بين السطور لراحة العين في القراءة
\onehalfspacing

% صفحة الغلاف
\begin{titlepage}
    \centering
    \vspace*{5cm}
    {\Huge \textbf{شرح كتاب الرسالة للإمام الشافعي}}\\[2cm]
    {\LARGE \textbf{لفضيلة الشيخ: سامي}}\\[3cm]
    {\Large \textit{تفريغ وتنقيح وتنسيق}}\\[1cm]
    {\large فريق العمل}\\[2cm]
    \vfill
    {\large \today}
\end{titlepage}

\frontmatter
\tableofcontents

\chapter*{كلمة المحرر}
\addcontentsline{toc}{chapter}{كلمة المحرر}

الحمد لله الذي بنعمته تتم الصالحات، والصلاة والسلام على المبعوث رحمة للعالمين، نبينا محمد وعلى آله وصحبه أجمعين، أما بعد؛

فإن كتاب «الرسالة» للإمام الجليل أبي عبد الله محمد بن إدريس الشافعي (المتوفى سنة 204 هـ) يمثل قطب الرحى في علوم الشريعة، واللبنة الأولى التي أُسس عليها علم أصول الفقه، ولا يستغني عنه طالب علم ولا فقيه. وقد يسر الله تعالى لفضيلة الشيخ سامي -حفظه الله- شرح هذا السفر العظيم وتوضيح مقاصده، وفك عباراته في مجالس علمية عامرة بالفوائد والدرر.

وحرصاً منا على تعميم النفع بهذه المجالس، وحفظاً لهذا الشرح النفيس من الضياع، قمنا بالاعتناء بتفريغات هذه الدروس وإخراجها في حلة مطبوعة تليق بمكانة الكتاب وشرحه، وقد اتبعنا في إعداد هذا الكتاب المنهجية التالية:

\begin{enumerate}
    \item \textbf{التنقيح اللغوي وإعادة الصياغة:} قمنا بتحويل لغة الشرح الشفهية المفرغة إلى لغة عربية فصحى مقروءة وسلسة، مع إزالة اللزمات الكلامية، وحذف الاستطرادات العابرة التي تخرج عن مقصود الشرح، وذلك حفاظاً على التسلسل المنطقي للأفكار، مع الحرص التام على عدم الإخلال بالمعنى العلمي الذي قصده الشيخ.
    \item \textbf{التنسيق والتبويب:} تم تقسيم مادة الكتاب إلى فصول ومباحث ذات عناوين واضحة ودقيقة تعكس مضمون كل قسم، مما يسهل على القارئ فهم الهيكلة العامة والعودة إلى المسائل بيسر.
    \item \textbf{العناية بالنصوص الشرعية:} تم تمييز الآيات القرآنية بالأقواس المخصصة لها، وكتابتها بالرسم المناسب، وتنسيق الأحاديث النبوية ووضعها بين علامات التنصيص لتبرز عن كلام الشارح.
    \item \textbf{استخلاص الفوائد والملاحق:} لم نكتفِ بمجرد تنسيق الشرح، بل قمنا بعد كل درس باستخلاص أهم الفوائد والدرر التي نثرها الشيخ طيات شرحه، وقمنا بتبويبها وتصنيفها موضوعياً (عقدياً، ومنهجياً، وأصولياً، ولغوياً) وإدراجها في ملاحق خاصة في نهاية الكتاب، لتكون زاداً سريعاً للقارئ ومراجعة مركزة لمقاصد الشرح.
\end{enumerate}

وقد اعتمدنا في تنسيق الكتاب وإخراجه فنياً على نظام التنضيد الاحترافي (\LaTeX) لضمان أعلى معايير الجودة الطباعية.

وختاماً، نسأل الله العلي القدير أن يجعل هذا العمل خالصاً لوجهه الكريم، وأن يثيب فضيلة الشيخ على ما قدم، وأن ينفع بهذا الكتاب كل من قرأه أو ساهم في إخراجه، إنه سميع مجيب.

\vspace{1cm}
\noindent
\textbf{المحرر} \\
\textbf{فريق العمل}

\mainmatter

% تضمين الدروس تباعاً
\chapter{مقدمة في شرح الرسالة وخطبة الإمام الشافعي}

\section{مقدمة الدرس وخطبة الحاجة}
إن الحمد لله نحمده ونستعينه ونستغفره، ونعوذ بالله تعالى من شرور أنفسنا ومن سيئات أعمالنا، من يهده الله فلا مضل له، ومن يضلل فلا هادي له. وأشهد أن لا إله إلا الله وحده لا شريك له، وأشهد أن محمداً عبده ورسوله.

﴿يَا أَيُّهَا الَّذِينَ آمَنُوا اتَّقُوا اللَّهَ حَقَّ تُقَاتِهِ وَلَا تَمُوتُنَّ إِلَّا وَأَنتُم مُّسْلِمُونَ﴾، ﴿يَا أَيُّهَا النَّاسُ اتَّقُوا رَبَّكُمُ الَّذِي خَلَقَكُم مِّن نَّفْسٍ وَاحِدَةٍ وَخَلَقَ مِنْهَا زَوْجَهَا وَبَثَّ مِنْهُمَا رِجَالًا كَثِيرًا وَنِسَاءً وَاتَّقُوا اللَّهَ الَّذِي تَسَاءَلُونَ بِهِ وَالْأَرْحَامَ إِنَّ اللَّهَ كَانَ عَلَيْكُمْ رَقِيبًا﴾، ﴿يَا أَيُّهَا الَّذِينَ آمَنُوا اتَّقُوا اللَّهَ وَقُولُوا قَوْلًا سَدِيدًا * يُصْلِحْ لَكُمْ أَعْمَالَكُمْ وَيَغْفِرْ لَكُمْ ذُنُوبَكُمْ وَمَن يُطِعِ اللَّهَ وَرَسُولَهُ فَقَدْ فَازَ فَوْزًا عَظِيمًا﴾.

أما بعد؛ فإن أصدق الحديث كتاب الله تعالى، وإن خير الهدي هدي محمد عليه الصلاة والسلام، وإن شر الأمور محدثاتها، وإن كل محدثة بدعة، وكل بدعة ضلالة، وكل ضلالة في النار.

هذا هو الدرس الأول في القراءة في كتاب «الرسالة» للإمام أبي عبد الله الشافعي رحمه الله تعالى. وقد اهتم الشافعي رحمه الله في كتابه بأمرين عظيمين: الأول: تحرير القواعد الأصولية، وإقامة الأدلة عليها من الكتاب والسنة، وتأييد كلامه بشواهد من اللغة العربية. الثاني: الإكثار من الأمثلة لزيادة الإيضاح والتطبيق لقضايا أصول الشريعة وفروعها، مع مناقشة المخالفين بأسلوب سهل جزل. وهو أول ما دُوّن في أصول الفقه، ويكفينا شرفاً أن نتلقى عن الإمام الشافعي، نقرأ له ونتعلم منه، ونتأدب بأدبه، ونحفظ شيئاً من القواعد التي أصّلها.

\section{أهمية كتاب الرسالة وتحقيقاته}
ألّف الشافعي «الرسالة» مرتين؛ الرسالة القديمة التي أرسلها لعبد الرحمن بن مهدي حين سأله عن دلالات النصوص والإجماع والقياس، ثم أملى الرسالة الجديدة على تلميذه الربيع بن سليمان المرادي.

يقول أبو القاسم عثمان بن سعيد الأنماطي عن الإمام المزني رحمه الله: «أنا أنظر في كتاب الرسالة عن الشافعي منذ خمسين سنة، ما أعلم أني نظرت فيه مرة إلا وأنا أستفيد منه شيئاً لم أكن عرفته». فظل خمسين سنة يراجع ويذاكر هذا الكتاب الجليل.

وقد حُقق هذا الكتاب تحقيقاً جليلاً على يد شيخ شيوخنا أبي الأشبال أحمد بن محمد شاكر رحمه الله، على نسخة الربيع بن سليمان، وهي من أقدم المخطوطات على وجه الأرض (كُتبت سنة 265هـ). كما حققها الدكتور رفعت فوزي عبد المطلب عند تحقيقه لموسوعة «الأم».

\section{قراءة إسناد الكتاب ونسب الإمام الشافعي}
قال الربيع بن سليمان: بسم الله الرحمن الرحيم. أخبرنا أبو عبد الله محمد بن إدريس بن العباس بن عثمان بن شافع بن السائب بن عبيد بن عبد يزيد بن هاشم بن المطلب بن عبد مناف.

والمطلب بن عبد مناف كان أخاً لهاشم بن عبد مناف، والمطلب هو الذي ذهب إلى المدينة ليأتي بابن أخيه هاشم (شيبة الحمد)، فأردفه خلفه ودخل مكة، فقال القرشيون: «من هذا؟»، قالوا: «هذا المطلب»، وقالوا عمن معه: «عبد المطلب» يظنونه عبداً اشتراه، فصار «عبد المطلب» لقباً اشتهر به جد النبي صلى الله عليه وسلم. والشافعي قرشي صليبة من بني المطلب، أبناء عمومة رسول الله صلى الله عليه وسلم.

\section{شرح خطبة الإمام الشافعي في الرسالة}
ابتدأ الشافعي كتابه بالبسملة والحمد، اقتداءً بالكتاب الكريم. قال رحمه الله:
«الحمد لله الذي خلق السماوات والأرض وجعل الظلمات والنور ثم الذين كفروا بربهم يعدلون. والحمد لله الذي لا يُؤدّى شُكر نعمة من نعمه إلا بنعمة منه تُوجب على مُؤدّي ماضي نعمه بأدائها نعمة حادثة يجب عليه شكره بها...»

\subsection{الفرق بين الحمد والشكر}
الحمد يكون باللسان والقلب على جهة التعظيم، ويكون على النعمة وغيرها. أما الشكر فلا يكون إلا على نعمة أو جميل يُقدّم، ويكون بالقلب واللسان والجوارح. فالحمد أعم سبباً (لأنه على النعمة وغيرها)، والشكر أعم متعلقاً (لأنه بالجوارح والقلب واللسان). ونحن نحمد الله على كمال صفاته وعلى عظيم إنعامه.

\subsection{فقه القلوب عند الشافعي}
في قول الشافعي: «لا يُؤدّى شكر نعمة من نعمه إلا بنعمة منه...» درس تربوي عظيم؛ فالعبد لا يستطيع أداء شكر الله إلا بإعانة الله وتوفيقه، وهذا التوفيق نعمة جديدة توجب شكراً جديداً. فالمؤمن دائم التقلب في نعم الله، ودائم الشكر لها، وكلما شكر زادته النعم، ﴿لَئِن شَكَرْتُمْ لَأَزِيدَنَّكُمْ﴾.

ثم قال: «ولا يبلغ الواصفون كُنه عظمته، الذي هو كما وصف نفسه وفوق ما يصفه به خلقه». وهذا أصل عقدي، فلا يجوز وصف الله بغير ما وصف به نفسه، ولا يجوز التكييف أو التمثيل، لأن الخلق لن يبلغوا حقيقة عظمته سبحانه.

\section{حال الخلق عند بعثة النبي صلى الله عليه وسلم}
ذكر الشافعي انقسام الناس عند مبعث النبي صلى الله عليه وسلم إلى قسمين:
الأول: أهل الكتاب (اليهود والنصارى) الذين بدلوا أحكام كتبهم وحرّفوها.
الثاني: المشركون عبدة الأصنام والأوثان من العرب، وعبدة النيران والنجوم من الأعاجم. فبعث الله محمداً صلى الله عليه وسلم لإنقاذ البشرية.

وفي قوله: «فلما بلغ الكتاب أجله فحقّ قضاء الله بإظهار دينه الذي اصطفى بعد استعلاء معصيته التي لم يرضَ...» لمحة تربوية؛ فالتمكين وإظهار الدين يكون بقضاء الله، والواجب علينا الأخذ بالأسباب الشرعية وعدم الركون لمعصية الله طمعاً في التمكين.

\section{خصائص النبي محمد صلى الله عليه وسلم ومقاصد رسالته}
اصطفى الله محمداً صلى الله عليه وسلم من خيار خيار العرب، وختم به النبوة، وعمّ برسالته الناس كافة. وخصّ قومه (العرب) بالنذارة أولاً ﴿وَأَنذِرْ عَشِيرَتَكَ الْأَقْرَبِينَ﴾ ليكونوا حملة الرسالة إلى غيرهم.

\subsection{الاعتزاز باللغة العربية ومقاصده}
يؤكد الشافعي أن القرآن نزل بلسان عربي مبين خالص لا تشوبه عُجمة. والاعتزاز بالعربية هو اعتزاز بالدين؛ فمتابعة الهدي الظاهر في اللغة والسمت تورث الاهتداء في الباطن. وقد نقل الشيخ أحمد شاكر عن والده محمود شاكر تحذيره من دعوات الاستعمار لترجمة القرآن وجعله أعجمياً في مستعمراتهم، لإذابة الشخصية الإسلامية وقطع الرابطة العربية بين المسلمين.

\section{آداب طالب العلم ومسلك النجاة}
ختم الشافعي هذه المقدمة ببيان طبقات الناس في العلم، وأن العلم الحقيقي هو علم الكتاب والسنة. وأوصى طلبة العلم بأمور:
1. بلوغ غاية الجهد في الاستكثار من العلم.
2. الصبر على العوارض (الفقر، التعب، الانشغال).
3. إخلاص النية لله عز وجل.
4. الرغبة إلى الله والتذلل له لطلب العون والفهم.

فمن أدرك أحكام الله نصاً واستدلالاً، ووفقه الله للقول والعمل، فاز بالفضيلة في دينه ودنياه، ونوّرت في قلبه الحكمة.
\chapter{القرآن كتاب هداية والاعتصام بالسنة}

\section{القرآن كتاب هداية}
القرآن تبيان لكل شيء، ولكن تبيان في ماذا؟ في أمر الهداية؛ فهو كتاب هداية، وليس كتاب جبر وهندسة وكيمياء وفيزياء. حتى لا يُقال: نحن نريد حلاً للمسألة الفلانية، فنقول: ﴿فَاسْأَلُوا أَهْلَ الذِّكْرِ إِن كُنتُمْ لَا تَعْلَمُونَ﴾، فيُسأل أهلها. وهذا أيضاً من الأصول العظيمة في القرآن. وقال تعالى: ﴿وَكَذَٰلِكَ أَوْحَيْنَا إِلَيْكَ رُوحًا مِّنْ أَمْرِنَا مَا كُنتَ تَدْرِي مَا الْكِتَابُ وَلَا الْإِيمَانُ وَلَٰكِن جَعَلْنَاهُ نُورًا نَّهْدِي بِهِ مَن نَّشَاءُ مِنْ عِبَادِنَا ۚ وَإِنَّكَ لَتَهْدِي إِلَىٰ صِرَاطٍ مُّسْتَقِيمٍ﴾.

\section{الوصية بالتمسك بالسنة}
فيا طالب الصراط المستقيم، إذا كنت تريد الصراط المستقيم، فعليك بسنة النبي الكريم محمد صلى الله عليه وسلم؛ استمسك بها، وعضّ عليها بالنواجذ، وادعُ إليها، واحذر من خلافها، تنل وتفز بجنة ربك وبرحمته.

نقف عند باب "كيف البيان". أسأل الله بأسمائه الحسنى وصفاته العلى أن ينفعني وإياكم بما قلنا وبما سمعنا.
\chapter{باب كيف البيان وأقسامه}

\section{مقدمة الدرس وتعريف البيان والمجمل}
الحمد لله رب العالمين، والصلاة والسلام على رسول الله، أما بعد؛ فهذا هو الدرس الثاني من دروس شرح كتاب «الرسالة» للإمام أبي عبد الله الشافعي رحمه الله تعالى. ونشرع اليوم في «باب كيف البيان».

يقول الشافعي رحمه الله: «والبيان اسم جامع لمعانٍ مجتمعة الأصول متشعبة الفروع».
لفهم هذا الباب لابد من معرفة المجمل والمبين.
المجمل في اللغة: المجموع. وفي الاصطلاح: ما احتمل معنيين أو أكثر من غير ترجيح لواحد منهما.
أما البيان في اللغة: فهو الإيضاح والإظهار لما كان فيه خفاء. 

وقد اشترط بعض الأصوليين سبق الخفاء ليسمى بياناً، بينما اعتبر آخرون أن كل ما يزيل الإشكال ويوضح المعنى فهو بيان. والشافعي يرى أن البيان اسم يجمع تحته أصولاً مجتمعة تتشعب إلى فروع، يدركها ويفهمها أهل اللغة العربية، لأن القرآن نزل بلسانهم. فأقل ما في هذه المعاني أنها بيان لمن خوطب بها، متقاربة الاستواء في الفهم عندهم، وإن كان بعضها أشد تأكيداً في البيان من بعض. ومن يجهل لسان العرب قد يظن فيها الاختلاف والتشعب. وفي هذا حث لطلبة العلم على تعلم اللغة العربية والتبحر فيها لفهم مراد الله ورسوله.

\section{وجوه البيان عند الإمام الشافعي}
يقول الشافعي رحمه الله: «فجماع ما أبان الله لخلقه في كتابه مما تعبدهم به لما مضى من حكمه جل ثناؤه من وجوه». وذكر وجوه البيان، وهي إجمالاً:
\begin{enumerate}
    \item ما أبانه لخلقه نصاً في كتابه.
    \item ما أحكم فرضه بكتابه وبينه على لسان نبيه صلى الله عليه وسلم.
    \item ما سنه رسول الله صلى الله عليه وسلم مما ليس لله فيه نص حكم في القرآن.
    \item ما فرض الله على خلقه الاجتهاد في طلبه.
\end{enumerate}

\subsection{الاجتهاد في طلب الحق وإعمال النصوص}
قبل التفصيل في أقسام البيان، أشار الشافعي إلى النوع المتعلق بالاجتهاد. فقد ابتلانا الله بطاعته في الفرائض، وابتلانا كذلك بطاعته في الاجتهاد للوصول إلى الحق في بعض المسائل. ومن ذلك الاجتهاد في استقبال القبلة لمن غاب عن عين المسجد الحرام؛ قال تعالى: ﴿وَحَيْثُ مَا كُنتُمْ فَوَلُّوا وُجُوهَكُمْ شَطْرَهُ﴾. ففرض علينا الاجتهاد بالعقول التي ركبها فينا، والعلامات التي نصبها لنا كالشمس والقمر والنجوم والجبال والرياح، لنهتدي بها للقبلة. 

وكذلك الاجتهاد في معرفة «العدل» لقبول شهادته أو حكمه؛ قال تعالى: ﴿وَأَشْهِدُوا ذَوَيْ عَدْلٍ مِّنكُمْ﴾، والعدل هو المسلم البالغ العاقل الذي سلم من أسباب الفسق وخوارم المروءة، فنجتهد في تطبيق هذه الشروط على الأشخاص.
وكذلك الاجتهاد في «جزاء الصيد»؛ كأن يصيد المحرم صيداً، فيجتهد الحكمان العدلان في تحديد المثل من النعم الذي يقابل هذا الصيد ليكون فدية. وهذا النوع من الاجتهاد يسميه الشافعي «القياس»، وهو القياس الصحيح الذي يعتمد على إعمال النصوص والقرائن، بخلاف القول بالاستحسان المجرد عن الدليل. ومن هذا يتبين أنه لا يجوز لأحد (غير رسول الله صلى الله عليه وسلم) أن يحلل أو يحرم إلا بالاستدلال من الكتاب أو السنة أو الإجماع أو القياس الصحيح.

\section{تفصيل أبواب البيان}

\subsection{البيان الأول: ما أبانه الله لخلقه نصاً في كتابه}
وهو ما جاء نصاً واضحاً لا يحتمل تأويلاً، كفرضية الصيام، والصلاة، والزكاة، والحج، وكتحريم الزنا والخمر والميتة والدم. وضرب الشافعي لذلك أمثلة تتضمن زيادة في البيان اللفظي والتأكيد، كقوله تعالى: ﴿فَمَن لَّمْ يَجِدْ فَصِيَامُ ثَلَاثَةِ أَيَّامٍ فِي الْحَجِّ وَسَبْعَةٍ إِذَا رَجَعْتُمْ ۗ تِلْكَ عَشَرَةٌ كَامِلَةٌ﴾، وقوله: ﴿وَوَاعَدْنَا مُوسَىٰ ثَلَاثِينَ لَيْلَةً وَأَتْمَمْنَاهَا بِعَشْرٍ فَتَمَّ مِيقَاتُ رَبِّهِ أَرْبَعِينَ لَيْلَةً﴾. فهذا من باب زيادة التبيين لجملة العدد لمن خوطبوا بالآية.

\subsection{البيان الثاني: ما نص القرآن على فرضه وبينت السنة بعض أحكامه}
وهو ما أتى فرضه في القرآن، ولكن بعض جزئياته أو شروطه احتاجت إلى تقييد أو تخصيص أو بيان من السنة. وضرب له أمثلة:
\begin{itemize}
    \item \textbf{الوضوء والغسل:} قال تعالى: ﴿إِذَا قُمْتُمْ إِلَى الصَّلَاةِ فَاغْسِلُوا وُجُوهَكُمْ وَأَيْدِيَكُمْ إِلَى الْمَرَافِقِ﴾ الآية. فأتى القرآن بالفرض، لكن السنة بينت الكيفية، وأن العدد يجزئ فيه المرة، والكمال في الثلاث، وبينت موجبات الوضوء والغسل، ووضحت أن غسل الأعقاب واجب في الوضوء بقوله صلى الله عليه وسلم: «ويل للأعقاب من النار».
    \item \textbf{المواريث والوصية:} فصل القرآن أنصبة المواريث تفصيلاً دقيقاً، لكن اشترط أن يكون التقسيم ﴿مِن بَعْدِ وَصِيَّةٍ يُوصِي بِهَا أَوْ دَيْنٍ﴾. فجاءت السنة المطهرة في حديث سعد بن أبي وقاص رضي الله عنه لتقيد الوصية بألا تتجاوز الثلث، فقال صلى الله عليه وسلم: «الثلث، والثلث كثير».
\end{itemize}

\subsection{البيان الثالث: ما أحكم الله فرضه مجملاً وبينت السنة كيفيته}
وهذا النوع يختلف عن سابقه في أن الفرض يأتي في القرآن مجملاً إجمالاً تاماً، وتتولى السنة تفصيل كل أحكامه وكيفياته. كقوله تعالى: ﴿إِنَّ الصَّلَاةَ كَانَتْ عَلَى الْمُؤْمِنِينَ كِتَابًا مَّوْقُوتًا﴾، وقوله: ﴿وَأَقِيمُوا الصَّلَاةَ وَآتُوا الزَّكَاةَ﴾، و﴿وَلِلَّهِ عَلَى النَّاسِ حِجُّ الْبَيْتِ﴾. فبين النبي صلى الله عليه وسلم عدد الصلوات ومواقيتها وسننها بقوله: «صلوا كما رأيتموني أصلي»، وبين أنصبة الزكاة، وبين مناسك الحج بقوله: «خذوا عني مناسككم». 

\subsection{البيان الرابع: ما سنه النبي صلى الله عليه وسلم استقلالاً}
وهو كل ما سنه رسول الله صلى الله عليه وسلم مما ليس فيه نص كتاب مباشر. كتحريم الجمع بين المرأة وعمتها، والمرأة وخالتها، وغيرها من التشريعات. وهذا البيان حجة وفرض يجب قبوله، لأن الله تعالى فرض في كتابه طاعة رسوله طاعة مطلقة، قال تعالى: ﴿وَمَا آتَاكُمُ الرَّسُولُ فَخُذُوهُ وَمَا نَهَاكُمْ عَنْهُ فَانتَهُوا﴾. 

فمن ادعى أنه يتبع القرآن فقط وترك السنة، فقد كذب القرآن الذي أمره باتباع السنة. ومن قَبِلَ عن رسول الله صلى الله عليه وسلم سنته، فقد قَبِلَ عن الله عز وجل أمره وفرضه. ومن فرّق بينهما، فقد ردّ على الله الذي فرض طاعة نبيه. فكل شيء شرعه النبي صلى الله عليه وسلم فهو بيان من الحكمة التي آتاه الله إياها، ويجب التسليم له والتصديق به.
\chapter{البيان الخامس: الاجتهاد والقياس وشروط المفتي}

\section{مقدمة في البيان الخامس (الاجتهاد)}
الحمد لله رب العالمين، والصلاة والسلام على نبينا محمد وعلى آله وصحبه أجمعين، أما بعد؛ فهذا هو الدرس الرابع من دروس شرح كتاب «الرسالة» للإمام أبي عبد الله الشافعي رحمه الله تعالى.

تحدثنا سابقاً عن وجوه البيان، ونصل اليوم إلى «البيان الخامس»، وهو بيان الاجتهاد؛ أي أن تجتهد للوصول إلى الحق عن طريق إعمال النص. وكعادة الإمام الشافعي في تأصيل القواعد، فإنه يكثر من ضرب الأمثلة لتتضح القاعدة، فذكر للاجتهاد أمثلة تطبيقية:

\section{أمثلة تطبيقية على الاجتهاد والقياس}

\subsection{الاجتهاد في تحديد جهة القبلة}
قال الله تبارك وتعالى: ﴿وَمِنْ حَيْثُ خَرَجْتَ فَوَلِّ وَجْهَكَ شَطْرَ الْمَسْجِدِ الْحَرَامِ﴾. الفرض على المسلم أينما كان أن يولي وجهه شطر المسجد الحرام. فإن كان معاينًا للكعبة (يراها أمامه)، فالفرض أن يتوجه لعينها يقيناً ولا مجال للاجتهاد. أما إن كان غائباً عنها، فالفرض أن يجتهد في التوجه إلى جهتها.

و«الشطر» في لغة العرب يأتي بمعنى النصف، ويأتي بمعنى الجهة، وهو المراد هنا. وقد استدل الشافعي بأشعار العرب على أن الشطر بمعنى الجهة (التلقاء والقصد)، كقول خفاف بن ندبة: «تُغني الرسالةُ شطرَ عمروٍ»، أي جهته. 
واجتهاد الغائب عن القبلة يكون بالدلائل والعلامات التي نصبها الله، قال تعالى: ﴿وَهُوَ الَّذِي جَعَلَ لَكُمُ النُّجُومَ لِتَهْتَدُوا بِهَا فِي ظُلُمَاتِ الْبَرِّ وَالْبَحْرِ﴾، وقال: ﴿وَعَلَامَاتٍ وَبِالنَّجْمِ هُمْ يَهْتَدُونَ﴾. فالمسلم يعمل عقله في هذه العلامات للوصول إلى جهة الكعبة، وهذا الإعمال للعلامات لمعرفة القبلة هو من باب الاجتهاد الصحيح.

\subsection{الاجتهاد في معرفة العدل}
ضرب الشافعي مثالاً آخر بالاجتهاد في تحديد العدول الذين تقبل شهادتهم وحكمهم. قال تعالى: ﴿وَأَشْهِدُوا ذَوَىْ عَدْلٍ مِّنكُمْ﴾. والعدل هو: من له مَلَكَةٌ تحمله على ملازمة التقوى والمروءة.
\begin{itemize}
    \item \textbf{التقوى:} اجتناب الشرك، والفسق، والبدعة، والإصرار على الصغيرة. وقال طلق بن حبيب: «التقوى أن تعمل بطاعة الله على نور من الله ترجو ثواب الله، وأن تترك معصية الله على نور من الله تخاف عقاب الله».
    \item \textbf{المروءة:} أن يلتزم الإنسان بما يليق بأمثاله قولا وفعلاً وسلوكاً، ويجتنب خوارم المروءة (كأن يأكل العالِم في الأسواق مثلاً).
\end{itemize}
فنحن نجتهد في تطبيق هذه الشروط على الأشخاص لنحكم بعدالتهم.

\subsection{الاجتهاد في جزاء الصيد (المثلية)}
المثال الثالث هو الاجتهاد في جزاء من قتل صيداً وهو محرم؛ قال تعالى: ﴿فَجَزَاءٌ مِّثْلُ مَا قَتَلَ مِنَ النَّعَمِ يَحْكُمُ بِهِ ذَوَا عَدْلٍ مِّنكُمْ﴾. بعد أن نجتهد في اختيار الحَكَمَين العدلين، يجتهد الحكمان في تحديد «المثل» من بهيمة الأنعام ليقابل الصيد البري المقتول. فإذا اصطاد حماراً وحشياً، حكما عليه ببقرة مققاربة له في الحجم، وإذا اصطاد ضبعاً حكما بكبش، وإذا اصطاد أرنباً حكما بعناق أو ما شابهه. ولا يعدل إلى «القيمة» المالية إلا إذا عُدم «المثل». وهذا الإلحاق بين الصيد والنعم هو عين «القياس الصحيح».

\section{مصادر التشريع وأنواع القياس عند الشافعي}
يقرر الشافعي أصلاً عظيماً: لا يجوز لأحد أبدًا أن يقول في شيء حلال أو حرام إلا من جهة العلم. وجهة العلم هي: الكتاب، أو السنة، أو الإجماع، أو القياس.

وبيّن أن القياس الصحيح يكون موافقاً للخبر، وأن الله ورسوله يحلان ويحرمان لأحد أمرين:
\begin{enumerate}
    \item \textbf{أمر تعبدي محض:} نص غير معقول المعنى (كعدد ركعات الصلاة، وتحريم الجمع بين المرأة وعمتها)، وهذا يوقف عنده ولا يقاس عليه.
    \item \textbf{معنى معقول (علّة):} أن يُحرم الشيء لعلة واضحة. فإذا وجدت هذه العلة في فرع لم ينص عليه، ألحقناه بالأصل. كالخمر حُرّمت لعلة الإسكار، فإذا وجد الإسكار في أي مشروب آخر (شعير، ذرة، حشيش) فهو حرام. وهذا يسمى \textbf{(قياس المعنى) أو (قياس العلة)}.
\end{enumerate}
وقد نجد فرعاً يشبه أصلين، فنلحقه بأولى الأشياء شبهاً به، وهذا يسمى \textbf{(قياس الشبه)}.

\section{بين الإجماع والاختلاف}
ينبه الشافعي إلى أدب عظيم لطالب العلم: ما أجمعت عليه الأمة لا يجوز الخلاف فيه (كحرمة الطواف بالقباب والأضرحة والذبح لها). أما ما اختلفت فيه الأمة اختلافاً سائغاً، فيُنظر في الدليل، ويُعمل بالأقرب للصواب. وقد قال الشافعي لمناظره يونس بن عبد الأعلى: «ألا يستقيم أن نكون إخواناً وإن لم نتفق في مسألة؟».
ولكن هذا في اختلاف الأفهام وتنوع الاستدلال، وليس في الاختلاف الذي يصادم النصوص الصريحة كمن يحلل ما حرم الله، أو يسوي بين المسلم والكافر، أو يدعو لمبادئ كفرية كالعلمانية والديمقراطية المستوردة المصادمة لحاكمية الشريعة بدعوى الخلاف، فهذا اختلاف تضاد لا يقبل.

\section{شروط المفتي والمجتهد}
بيّن الشافعي أن المجتهد يجب أن تتوفر فيه علوم أساسية، أهمها:

\subsection{العلم بلغة العرب}
فالقرآن نزل بلسان عربي مبين، ولغة العرب هي لغة التشريع. والصحابة كانوا يفهمون دلالات الألفاظ بالسليقة. فعلى المفتي أن يعلم العام، والخاص، والمطلق، والمقيد، والعام الذي يراد به الخصوص. والاعتزاز بالعربية هو جزء من الاعتزاز بالهوية الإسلامية. وقد زعم بعض المقلدين وجود كلمات أعجمية في القرآن، ورد الشافعي ذلك بقوة، مقرراً أن كل ما في القرآن فهو عربي خالص، وما وافق لغات أخرى فهو إما من توارد اللغات، أو كلمات استعملتها العرب وعرّبتها قبل نزول القرآن فصارت من لسانها.

\subsection{العلم بالناسخ والمنسوخ ودلالات الأوامر}
يجب أن يعلم المفتي الناسخ والمنسوخ (ومنه ما نُسخ حكماً وبقي تلاوة، أو نُسخ تلاوة وبقي حكماً كآية الرجم). كما يجب أن يعلم دلالة الأوامر والنواهي؛ فمن الأوامر ما يكون للوجوب، ومنها ما يكون للندب، ومنها ما يكون للاباحة.

\section{التحذير من القول بلا علم ومن التقليد الأعمى}
ختم الشافعي هذا الباب بكلمة تكتب بماء الذهب: \textbf{«فالواجب على العالمين ألّا يقولوا إلا من حيث علموا. وقد تكلم في العلم من لو أمسك عن بعض ما تكلم فيه منه، لكان الإمساك أولى به، وأقرب من السلامة له إن شاء الله»}.
فلا يجوز للعالم أن يتكلم إلا بعلم راسخ من الكتاب والسنة. كما حذر من التقليد الأعمى، وعبر عنه بقوله: «وبالتقليد أغفل من أغفل منهم، والله يغفر لنا ولهم». فالمقلد الأعمى يتخبط في الظلمات، والواجب على طالب العلم إذا بلغه قول أن يسأل عن دليله الشرعي.
\chapter{إحاطة الأمة بلغة العرب وسنة النبي}

\section{سعة لغة العرب وحفظ السنة}
لسان العرب هو أوسع الألسنة مذهباً وأكثرها ألفاظاً، ولا يوجد رجل أحاط بلغة العرب إحاطة تامة إلا أن يكون نبياً. وكذلك لا يوجد عالم أحاط بسنة النبي صلى الله عليه وسلم إحاطة تامة بمفرده. ولكن كما أن علم اللسان لا يغيب عن مجموع العرب والعلماء، فإن علم السنة لا يغيب عن مجموع العلماء قاطبة.

وهذا مقصد عظيم جداً؛ أن يعلم العلماء أن السنة والعلم لا يضيعان في مجموعهم، بشرط أن يتقوا الله سبحانه وتعالى وأن يكونوا على حرص على الخير. 

\section{عصمة الأمة والاجتهاد الجماعي}
من المشاهد أن إحاطة الفرد الواحد بالسنة من الصعوبة بمكان، فالعالم قد ينسى وقد يجهل بعض السنن، كما حدث لأمير المؤمنين عمر بن الخطاب رضي الله عنه، وهو من هو في العلم والفضل، حين خفي عليه حديث الاستئذان ثلاثاً، ونسي مسألة التيمم للمجنب. فالعالم الفرد ينسى ويجهل، لكن مجموع الأمة معصومة عن الخطأ والضلال، ومنهج السلف الصالح معصوم عن الخطأ والضلال. 

والمقصد العظيم الذي يهدف إليه الشافعي رحمه الله تعالى هو أن الاجتهاد يكون جماعياً، بشرط أن يكون العلماء أتقياء أوفياء لدينهم، لا تدخلهم عصبية، ولا يدخلهم تقليد، ولا اتباع شهوة ولا هوى. بل يكون قصدهم نصرة سنة النبي صلى الله عليه وسلم، متبعين للشرع غير متأثرين بواقع منحرف، ولا خائفين من أحد إلا من الله عز وجل.

\section{كيف يحيط طالب العلم بالسنة؟}
من أراد أن يكاد يحيط علماً بالسنة، فعليه بالكتب الستة (صحيح البخاري، وصحيح مسلم، وسنن أبي داود، وسنن النسائي، وسنن الترمذي، وسنن ابن ماجه). ويضم إليها: مسند الإمام أحمد، ومسند أبي يعلى، ومسند البزار، ومعاجم الطبراني، وسنن الدارمي، وصحيحي ابن خزيمة وابن حبان، ومستدرك الحاكم، وسنن البيهقي.

من أحاط بهذه الدواوين يكاد أن يحيط بالسنة، وهيهات ألا ينسى أو يجهل أو يغفل، لكنه في النهاية يكون قد وصل إلى علم راسخ في سنة رسول الله صلى الله عليه وسلم. فرحم الله الشافعي وطيب ثراه، ونسأل الله أن يؤدبنا بأدبه ويفقهنا بفقهه ويرزقنا الإخلاص في القول والعمل.

\section{وصية للعوام بالرجوع للعلماء والصدق مع الله}
أما الكلام السابق فهو موجه لطالب العلم، أما العامي فنقول له: الزم كبار العلماء، فإذا ألمت بك ملمة فاسأل الكبار. 
ومع ذلك، فهذا العامي لو صدق مع الله لكان خيراً له، فالله عز وجل يقول: ﴿وَلَقَدْ يَسَّرْنَا الْقُرْآنَ لِلذِّكْرِ فَهَلْ مِن مُّدَّكِرٍ﴾. ولو أنه استفتى قلبه بصدق، كما قال النبي صلى الله عليه وسلم: «استفت قلبك وإن أفتاك المفتون وأفتوك»، وكان وقوفه وسؤاله بصدق كأنه واقف بين يدي الله عز وجل؛ هداه الله، كما وعد ربنا: ﴿وَالَّذِينَ جَاهَدُوا فِينَا لَنَهْدِيَنَّهُمْ سُبُلَنَا﴾.

\section{خاتمة ودعاء}
اللهم اقسم لنا من خشيتك ما تحول به بيننا وبين معاصيك، ومن طاعتك ما تبلغنا بها جنتك، ومن اليقين ما تهون به علينا مصائب الدنيا. اللهم متعنا بأسماعنا وأبصارنا وقوتنا ما أحييتنا، واجعله الوارث منا، واجعل ثأرنا على من ظلمنا، وانصرنا على من عادانا، اللهم لا تجعل مصيبتنا في ديننا، ولا تجعل الدنيا أكبر همنا ولا مبلغ علمنا، ولا النار مصيرنا يا أرحم الراحمين.

وصلى الله وسلم وبارك على سيد الأولين والآخرين وعلى آله وصحبه وسلم.
\chapter{القول بعلم وذم التقليد وعربية القرآن}

\section{وجوب القول بعلم وذم التقليد الأعمى}
ابتدأ الإمام الشافعي رحمه الله بتأصيل عظيم، حيث قال: «فالواجب على العالِمين ألا يقولوا إلا من حيث علموا. وقد تكلم في العلم من لو أمسك عن بعض ما تكلم فيه منه، لكان الإمساك أولى به وأقرب من السلامة له إن شاء الله». فلا يجوز للعالم أن يتكلم إلا بعلم راسخ، ولا يجوز لأحد أن يُحلل أو يُحرم إلا بدليل شرعي.
ثم حذر رحمه الله من التقليد الأعمى قائلاً: «وبالتقليد أغفل من أغفل منهم، والله يغفر لنا ولهم». فلا ينبغي لطالب العلم أن يقلد تقليداً أعمى، بل يجب أن يعرف الدليل وفقه السلف لهذا الدليل؛ فإن التقليد بغير دليل شرعي ضياع، وكما قيل: «لا يقلد إلا عصبي أو غبي».

\section{مكانة لغة العرب وسعتها}
القرآن الكريم نزل بلسان عربي مبين، والعالم الذي يتصدر للفتوى والاجتهاد لا بد أن يكون عالماً بلسان العرب؛ ليعلم متى يكون اللفظ عاماً أو خاصاً، ومطلقاً أو مقيداً. والعلم بلسان العرب يستوي مع العلم بالسنن عند أهل الفقه؛ فكما أنه لا يوجد فرد جمع السنن كلها، كذلك لا يوجد من أحاط بلغة العرب كاملة إلا نبي، ولكنها لا تغيب عن مجموع الأمة. وإذا جُمع علم عامة أهل العلم، لم يفتهم شيء من السنن ولا من لغة العرب.

\section{عربية القرآن الخالصة والرد على الشبهات}
أصّل الشافعي رحمه الله لحقيقة هامة، وهي أنه لا يوجد في القرآن أي حرف غير عربي، وأن القرآن كله عربي خالص، مستدلاً بقوله تعالى: ﴿وَمَا أَرْسَلْنَا مِن رَّسُولٍ إِلَّا بِلِسَانِ قَوْمِهِ﴾. 
ورد على من زعم وجود كلمات أعجمية (مثل: استبرق، وسندس) بأن العرب قد تكلمت بها وعرّبتها قبل نزول القرآن فصارت من لغتها. والاعتزاز بالعربية هو جزء من الاعتزاز بالدين والشخصية الإسلامية، ولا يجوز لأهل لسان النبي صلى الله عليه وسلم أن يكونوا أتباعاً لأهل لسان غيره، بل يجب أن يكون لسانهم هو المتبوع.

\section{النصيحة للمسلمين ووجوه البيان في لغة العرب}
بيّن الشافعي أن توضيح لغة العرب للناس هو من باب النصيحة، والنصيحة للمسلمين فرض ودين، كما قال النبي صلى الله عليه وسلم: «الدين النصيحة».
وقد خاطب الله العرب بكتابه على ما تعرف من معانيها، ومن وجوه ذلك في لسان العرب:
\begin{itemize}
    \item \textbf{العام الظاهر:} الذي يراد به العموم.
    \item \textbf{العام المراد به الخصوص:} لفظ عام لكن المراد به خاص، ويُستدل على ذلك بسياق الكلام.
    \item \textbf{تعدد الأسماء للمسمى الواحد:} كالأسد (ليث، حفص، حيدرة).
    \item \textbf{تعدد المعاني للاسم الواحد:} كالعين (عين البصر، عين الماء، الجاسوس، الذهب).
\end{itemize}
وهذه الوجوه تعرفها العرب بالسليقة، ومَنْ جَهِلَ لسانهم قد يستنكرها، ولا عبرة بإنكار من جهل لغة العرب التي نزل بها القرآن.

\section{خطورة التكلف والقول بغير علم}
ختم الشافعي هذا المبحث بقاعدة منهجية عظيمة، فقال: «ومن تكلف ما جهل، وما لم تثبته معرفته، كانت موافقته للصواب -إن وافقه من حيث لا يعرفه- غير محمودة، وكان بخطئه غير معذور».
فمن تكلم في الدين بغير علم ولا التزام بمنهج السلف، فإن وافق الصواب صدفة فإنه لا يُحمد على ذلك لأنه لم يسلك الطريق الصحيح، وإن أخطأ فلا يُعذر بخطئه. أما المجتهد الذي يسير على منهج السلف، فإنه إن أصاب فله أجران، وإن أخطأ فله أجر وهو معذور.
\chapter{بيان ما نزل من الكتاب عاماً وأقسامه}

\section{مقدمة الدرس وتعريف العام والخاص}
الحمد لله رب العالمين، والصلاة والسلام على نبينا محمد وعلى آله وصحبه أجمعين، أما بعد؛ فهذا الدرس يتابع تقريرات الإمام الشافعي رحمه الله في كتاب «الرسالة» حول لغة العرب التي نزل بها القرآن، وأهمية فهم دلالات الألفاظ فيها كالعام والخاص والمطلق والمقيد.

وقد عرّف الإمام الشوكاني «العام» بأنه: اللفظ المستغرق لجميع ما وضع له بحسب وضع واحد دفعة (مثل: يا أيها الناس). أما «التخصيص» فهو: إخراج بعض ما كان داخلاً تحت العموم. وعرفه ابن الحاجب بأنه: قصر العام على بعض مسمياته.

\section{أقسام العام في القرآن الكريم}
قسّم الإمام الشافعي ما نزل من القرآن الكريم من حيث العموم والخصوص إلى عدة أقسام:

\subsection{القسم الأول: العام الذي يراد به العموم ولا يدخله التخصيص}
وهو اللفظ العام الذي يستغرق جميع أفراده ولا يخرج منه شيء. ومثّل له الشافعي بقوله تعالى: ﴿اللَّهُ خَالِقُ كُلِّ شَيْءٍ وَهُوَ عَلَىٰ كُلِّ شَيْءٍ وَكِيلٌ﴾، وقوله: ﴿خَلَقَ السَّمَاوَاتِ وَالْأَرْضَ﴾، وقوله: ﴿وَمَا مِن دَابَّةٍ فِي الْأَرْضِ إِلَّا عَلَى اللَّهِ رِزْقُهَا﴾. فكل شيء في السماوات والأرض خلقه الله، وكل دابة فرزقها على الله، فهذا عام باقٍ على عمومه ولا مخصص له أبداً.

\subsection{القسم الثاني: العام الذي يجمع بين العموم والخصوص}
وهو أن تأتي الآية بلفظ عام، لكن يراد ببعض أجزائها العموم وببعضها الآخر الخصوص. ومثّل لذلك بقوله تعالى: ﴿مَا كَانَ لِأَهْلِ الْمَدِينَةِ وَمَنْ حَوْلَهُم مِّنَ الْأَعْرَابِ أَن يَتَخَلَّفُوا عَن رَّسُولِ اللَّهِ وَلَا يَرْغَبُوا بِأَنفُسِهِمْ عَن نَّفْسِهِ﴾. 
فقوله: ﴿أَن يَتَخَلَّفُوا﴾ عام يراد به الخصوص، إذ لا يشمل العجزة والنساء والأطفال والمرضى، بل يخص من يطيق الجهاد. أما قوله: ﴿وَلَا يَرْغَبُوا بِأَنفُسِهِمْ عَن نَّفْسِهِ﴾ فهو عام يراد به العموم؛ فلا يجوز لأي مؤمن (سواء كان صحيحاً أو مريضا، رجلاً أو امرأة) أن يرغب بنفسه عن نفس رسول الله صلى الله عليه وسلم.

ومن ذلك أيضاً قوله تعالى: ﴿إِنَّا خَلَقْنَاكُم مِّن ذَكَرٍ وَأُنثَىٰ وَجَعَلْنَاكُمْ شُعُوبًا وَقَبَائِلَ لِتَعَارَفُوا﴾ فهذا عام يشمل كل البشر، ثم قال: ﴿إِنَّ أَكْرَمَكُمْ عِندَ اللَّهِ أَتْقَاكُمْ﴾ وهذا عام يراد به الخصوص؛ لأن التقوى لا تكون إلا لمن عقلها وكان مكلفاً بالغاً، ويخرج منها غير المكلفين كالصبي والمجنون، لقوله صلى الله عليه وسلم: «رفع القلم عن ثلاثة: عن النائم حتى يستيقظ، وعن الصبي حتى يبلغ، وعن المجنون حتى يفيق».

\subsection{القسم الثالث: العام الظاهر الذي يراد به الخصوص}
وهو لفظ مخرجه عام على كل الناس، ولكن يُراد به فئة معينة يُستدل عليها بالسياق أو بأدلة أخرى، ومن أمثلته:
\begin{itemize}
    \item قوله تعالى: ﴿الَّذِينَ قَالَ لَهُمُ النَّاسُ إِنَّ النَّاسَ قَدْ جَمَعُوا لَكُمْ فَاخْشَوْهُمْ﴾. لفظ (الناس) عام، لكنه هنا أُطلق وأُريد به الخصوص في المواضع الثلاثة: فالذين قالوا هم بضعة أشخاص (نفر قليل)، والمقول لهم هم النبي صلى الله عليه وسلم وصحابته، والذين جمعوا هم جيش أبي سفيان المنصرف من أحد. وبقية أهل الأرض لم يقولوا ولم يُقَل لهم ولم يجمعوا.
    \item قوله تعالى: ﴿يَا أَيُّهَا النَّاسُ ضُرِبَ مَثَلٌ فَاسْتَمِعُوا لَهُ ۚ إِنَّ الَّذِينَ تَدْعُونَ مِن دُونِ اللَّهِ لَن يَخْلُقُوا ذُبَابًا﴾. الخطاب هنا بلفظ (الناس) عام، ولكنه مخصوص بعبدة الأصنام، ويخرج منه المؤمنون الموحدون.
    \item قوله تعالى: ﴿ثُمَّ أَفِيضُوا مِنْ حَيْثُ أَفَاضَ النَّاسُ﴾. المراد بالناس هنا الحجاج الحاضرون في ذلك الموقف، وليس كل الناس على وجه الأرض.
    \item قوله تعالى: ﴿وَقُودُهَا النَّاسُ وَالْحِجَارَةُ﴾. استثنى القرآن الكريم من هذا العموم أهل الإيمان والأنبياء بقوله: ﴿إِنَّ الَّذِينَ سَبَقَتْ لَهُم مِّنَّا الْحُسْنَىٰ أُولَٰئِكَ عَنْهَا مُبْعَدُونَ﴾.
\end{itemize}

\section{أثر السياق والتركيب والإضافة في بيان المعنى}
دلالات الألفاظ في لغة العرب تختلف باختلاف ثلاثة أمور: السياق، والتركيب، والإضافة.

\textbf{أثر التركيب:} 
يغير المعنى تماماً مع بقاء نفس الألفاظ؛ فجملة «ما في البيت إلا محمد» تفيد الحصر المطلق بأن البيت خالٍ إلا منه، بينما جملة «ما محمد إلا في البيت» (وهي مكونة من نفس الألفاظ الخمسة) تفيد حصر مكانه هو، وقد يكون معه في البيت غيره.

\textbf{أثر الإضافة:} 
«يد النملة» تختلف عن «يد الفيل»، ولله المَثَل الأعلى، فإضافة الصفة للخالق سبحانه تقتضي الكمال والجلال الذي لا يشابهه فيه شيء من خلقه.

\textbf{أثر السياق:} 
وقد ضرب الشافعي له مثلاً بكلمة (القرية)، فالعرب تطلق القرية وتريد بها الأبنية، وتطلقها وتريد بها السُّكان، والسياق هو المُحَدِّد. قال تعالى: ﴿وَاسْأَلْهُمْ عَنِ الْقَرْيَةِ الَّتِي كَانَتْ حَاضِرَةَ الْبَحْرِ إِذْ يَعْدُونَ فِي السَّبْتِ﴾. السياق يبيّن أن المعتدي في السبت ليس الجدران والأبنية، بل سكان القرية. وكذلك قوله: ﴿رَبَّنَا أَخْرِجْنَا مِنْ هَٰذِهِ الْقَرْيَةِ الظَّالِمِ أَهْلُهَا﴾، فليس كل أهل القرية ظَلَمة لوجود مستضعفين مؤمنين فيها، لكن الأغلبية أو المتسلطين كانوا هم الظلمة. فهذه الوجوه تفهمها العرب بسليقتها، ومن جهل لغة العرب استشكل عليه فهم كتاب الله.
\chapter{تخصيص السنة لعموم القرآن ووجوب طاعة الرسول}

\section{مقدمة في تخصيص السنة للقرآن}
الحمد لله والصلاة والسلام على رسول الله، أما بعد؛ فقد تناولنا في الدرس السابق ما نزل من القرآن عاماً وأقسامه. ونشرع في هذا الدرس في بيان قاعدة أصولية جليلة قررها الإمام الشافعي رحمه الله، وهي: «ما نزل من القرآن عاماً ودلت السنة على أنه يراد به الخاص». فالسنة المطهرة تأتي لتخصص عموم القرآن الكريم وتُقيد مطلقه، وقد ضرب الشافعي لذلك ستة أمثلة فقهية تطبيقية:

\section{الأمثلة التطبيقية على تخصيص السنة للقرآن}

\subsection{المثال الأول: آيات المواريث والوصية}
قال تعالى: ﴿وَلِأَبَوَيْهِ لِكُلِّ وَاحِدٍ مِّنْهُمَا السُّدُسُ مِمَّا تَرَكَ إِن كَانَ لَهُ وَلَدٌ﴾، وقال: ﴿وَلَكُمْ نِصْفُ مَا تَرَكَ أَزْوَاجُكُمْ﴾. هذه الآيات جاءت بلفظ عام يشمل كل والد وولد وكل زوج وزوجة. لكن السنة المطهرة خصصت هذا العموم، فبينت أنه لا يتوارث أهل ملتين (المسلم والكافر)، وأن القاتل لا يرث.

كذلك في الوصية، قال تعالى: ﴿مِن بَعْدِ وَصِيَّةٍ يُوصَىٰ بِهَا أَوْ دَيْنٍ﴾، فالآية عامة في أي مقدار للوصية، لكن السنة خصصتها بحديث سعد بن أبي وقاص رضي الله عنه: «الثلث والثلث كثير»، فلا تنفذ الوصية فيما زاد عن الثلث. كما بينت السنة وأجمع المسلمون على أن سداد الدين مُقَدَّم على الوصية، رغم أن الآية ذكرت الوصية أولاً.

\subsection{المثال الثاني: آية الوضوء وغسل القدمين}
قال تعالى: ﴿إِذَا قُمْتُمْ إِلَى الصَّلَاةِ فَاغْسِلُوا وُجُوهَكُمْ وَأَيْدِيَكُمْ إِلَى الْمَرَافِقِ وَامْسَحُوا بِرُءُوسِكُمْ وَأَرْجُلَكُمْ إِلَى الْكَعْبَيْنِ﴾. قصد الله عز وجل غسل القدمين كغسل الوجه واليدين، وظاهر الآية أنه لا يجزئ إلا الغسل لجميع المتوضئين. لكن سنة النبي صلى الله عليه وسلم بتواتر مسحه على الخفين، دلت على أن المراد بغسل القدمين بعض المتوضئين دون بعض، فمن لبس الخفين على طهارة خُصَّ بجواز المسح، ومن لم يلبسهما وجب عليه الغسل.

\subsection{المثال الثالث: آية حد السرقة}
قال تعالى: ﴿وَالسَّارِقُ وَالسَّارِقَةُ فَاقْطَعُوا أَيْدِيَهُمَا﴾. لفظ «السارق» عام، ولفظ «اليد» عام (يشمل من أطراف الأصابع إلى الإبط). فجاءت السنة المطهرة لتخصص السارق بمن سرق مالاً محرزاً يبلغ ربع دينار فصاعداً (فلا قطع في ثمر معلق ولا في أقل من ربع دينار)، وخصصت السنة اليد بـ «الكف» ليُقطع من المفصل (الرسغ) وليس من المرفق أو العضد.

\subsection{المثال الرابع: آية حد الزنا}
قال تعالى: ﴿الزَّانِيَةُ وَالزَّانِي فَاجْلِدُوا كُلَّ وَاحِدٍ مِّنْهُمَا مِائَةَ جَلْدَةٍ﴾، وهذا عام. لكن السنة دلت على أن المراد بجلد المائة هما «الحران البكران» فقط، لأن النبي صلى الله عليه وسلم رجم الثيب المحصن ولم يجلده (كما في قصة ماعز والغامدية). أما الإماء والعبيد، فقد بين القرآن أن عليهم نصف ما على المحصنات (أي خمسين جلدة)، فخُصصت الآية الكريمة بالسنة.

\subsection{المثال الخامس: سهم ذوي القربى في الغنائم}
قال تعالى: ﴿وَاعْلَمُوا أَنَّمَا غَنِمْتُم مِّن شَيْءٍ فَأَنَّ لِلَّهِ خُمُسَهُ وَلِلرَّسُولِ وَلِذِي الْقُرْبَىٰ﴾. «ذو القربى» لفظ عام قد يشمل كل قريش، أو كل بني عبد مناف. لكن السنة خصصته بطائفتين فقط هما: بنو هاشم وبنو المطلب. ولما سأل عثمان بن عفان وجبير بن مطعم (وهما من بني عبد شمس وبني نوفل أبناء عبد مناف) عن سبب إعطاء بني المطلب ومنعهم مع أن قرابتهم واحدة، شبّك النبي صلى الله عليه وسلم أصابعه وقال: «إنما بنو هاشم وبنو المطلب شيء واحد في الجاهلية والإسلام». فدلت السنة على تخصيص ذوي القربى بهاتين البطنتين.

\subsection{المثال السادس: استثناء السَّلَب من عموم الغنائم}
الآية السابقة جعلت الغنائم مقسومة (يُؤخذ خمسها ويوزع الباقي)، لكن السنة خصت «السَّلَب» (وهو ما يكون على المقتول من سلاح وعتاد ومال) من عموم الغنيمة، فجعلته للقاتل خالصاً إذا قتل مشركاً في إقبال وحالة قتال. ودليل ذلك قضاء النبي صلى الله عليه وسلم لـ أبي قتادة رضي الله عنه يوم حنين بسلب المشرك الذي قتله بعد أن أقام البينة، وقال: «من قتل قتيلاً له عليه بينة فله سلبه». فهذا السلب لا يُخمّس ولا يدخل في عموم الغنيمة.

\section{فرضية طاعة رسول الله صلى الله عليه وسلم}
بعد أن بيّن الشافعي أمثلة تخصيص السنة للقرآن، أراد أن يُؤصل مسألة في غاية الأهمية، وهي أن اتباع النبي صلى الله عليه وسلم فرض حتمي، وأن سنته مستقلة بالتشريع.

\subsection{اقتران الإيمان بالله بالإيمان بالرسول}
جعل الله الإيمان برسوله صلى الله عليه وسلم كمالاً وشرطاً للإيمان به جل وعلا. قال تعالى: ﴿إِنَّمَا الْمُؤْمِنُونَ الَّذِينَ آمَنُوا بِاللَّهِ وَرَسُولِهِ﴾، فلو آمن عبد بالله ولم يؤمن برسوله لم يقع عليه اسم الإيمان أبداً.

\textbf{تنبيه منهجي:} أورد الإمام الشافعي رحمه الله في هذا الموضع استدلالاً بقوله تعالى: ﴿فَآمِنُوا بِاللَّهِ وَرَسُولِهِ﴾، مع أن الآية في سورة النساء جاءت بلفظ الجمع: ﴿فَآمِنُوا بِاللَّهِ وَرُسُلِهِ﴾، والآيات الدالة على الإفراد كثيرة في سور أخرى (كالأعراف والتغابن). وقد أشار العلامة أحمد شاكر رحمه الله إلى أن هذا وهم وسبق لسان من الإمام الشافعي، وأن بقاءه في المخطوطات قروناً ينسخه العلماء دون تنبه يدل على خطورة «التقليد» حتى لكبار الأئمة. وهذا تطبيق عملي لقول الشافعي نفسه: «وبالتقليد أغفل من أغفل منهم، والله يغفر لنا ولهم»، فالعصمة للوحيين فقط.

\subsection{السنة هي «الحكمة» في القرآن}
قال تعالى: ﴿لَقَدْ مَنَّ اللَّهُ عَلَى الْمُؤْمِنِينَ إِذْ بَعَثَ فِيهِمْ رَسُولًا مِّنْ أَنفُسِهِمْ يَتْلُو عَلَيْهِمْ آيَاتِهِ وَيُزَكِّيهِمْ وَيُعَلِّمُمُ الْكِتَابَ وَالْحِكْمَةَ﴾، وقال موجهاً الخطاب لنساء النبي: ﴿وَاذْكُرْنَ مَا يُتْلَىٰ فِي بُيُوتِكُنَّ مِنْ آيَاتِ اللَّهِ وَالْحِكْمَةِ﴾.
قال الشافعي: «سمعت من أرضى من أهل العلم بالقرآن يقول: الحكمة سُنَّة رسول الله صلى الله عليه وسلم». فإن الله منّ على خلقه بتعليمهم الكتاب (وهو القرآن) والحكمة (وهي السنة)، ولا يجوز أن يقال الحكمة هنا إلا سنة رسول الله صلى الله عليه وسلم لأنها قُرنت بكتاب الله.

\subsection{الرد على منكري السنة (القرآنيين)}
إن من يدعي الاكتفاء بالقرآن ليرد سنة النبي صلى الله عليه وسلم هو في الحقيقة مكذب للقرآن. فالقرآن هو الذي قال: ﴿وَمَا آتَاكُمُ الرَّسُولُ فَخُذُوهُ وَمَا نَهَاكُمْ عَنْهُ فَانتَهُوا﴾.
فمن قَبِلَ عن رسول الله صلى الله عليه وسلم سنته، فقد قَبِلَ عن الله فَرْضَه. ومن فرَّق بينهما وقَبِلَ واحداً وترك الآخر، فقد ردَّ على الله أمره. وسنة رسول الله صلى الله عليه وسلم مُبَيِّنةٌ عن اللهِ مرادَه، دليلةٌ على خاصِّه وعامِّه، ولا يجوز لأحد أن يُخالفها كائناً من كان.
\chapter{فرض طاعة رسول الله ومكانة السنة من القرآن}

\section{طاعة الرسول صلى الله عليه وسلم مطلقة ومستقلة}
الحمد لله والصلاة والسلام على رسول الله، أما بعد؛ فقد عقد الإمام الشافعي رحمه الله تعالى في كتاب «الرسالة» باباً في بيان فرض طاعة رسول الله صلى الله عليه وسلم، سواء أكانت مقرونة بطاعة الله عز وجل أم كانت مذكورة وحدها.

فالله عز وجل فرض طاعة نبيه صلى الله عليه وسلم في كل حال، وسواء كان النبي صلى الله عليه وسلم مبيناً لما جاء عن ربه، أو أتى بحكم مستقل ليس منصوصاً عليه في كتاب الله، فالفرض علينا أن نطيعه. قال تعالى: ﴿وَمَا كَانَ لِمُؤْمِنٍ وَلَا مُؤْمِنَةٍ إِذَا قَضَى اللَّهُ وَرَسُولُهُ أَمْرًا أَن يَكُونَ لَهُمُ الْخِيَرَةُ مِنْ أَمْرِهِمْ ۗ وَمَن يَعْصِ اللَّهَ وَرَسُولَهُ فَقَدْ ضَلَّ ضَلَالًا مُّبِينًا﴾. فلا خيار لأهل الإيمان في طاعة رسول الله صلى الله عليه وسلم متى ما صح عنه الأمر.

\section{طاعة أولي الأمر مقيدة ومستثناة}
قال تعالى: ﴿يَا أَيُّهَا الَّذِينَ آمَنُوا أَطِيعُوا اللَّهَ وَأَطِيعُوا الرَّسُولَ وَأُولِي الْأَمْرِ مِنكُمْ﴾. يُلاحظ أن الله عز وجل كرر لفظ "أطيعوا" مع الرسول صلى الله عليه وسلم للدلالة على استقلال طاعته، بينما لم يكرره مع "أولي الأمر"؛ لأن طاعة أولي الأمر (من الأمراء والعلماء) مقيدة بطاعة الله وطاعة رسوله.

فلا طاعة لمخلوق في معصية الخالق، وقد قال النبي صلى الله عليه وسلم: «إنما الطاعة في المعروف». فطاعة أولي الأمر ليست طاعة مطلقة في كل شيء، بل هي طاعة مستثناة؛ فإذا أمروا بمعصية فلا سمع لهم ولا طاعة في تلك المعصية. ولهم علينا حق السمع والطاعة وعدم الخروج عليهم ما أقاموا شرع الله، وعليهم إقامة الدين وحفظ الثغور وأعراض ودماء المسلمين. 

\section{الرد عند التنازع وحجية القياس}
قال تعالى: ﴿فَإِن تَنَازَعْتُمْ فِي شَيْءٍ فَرُدُّوهُ إِلَى اللَّهِ وَالرَّسُولِ إِن كُنتُمْ تُؤْمِنُونَ بِاللَّهِ وَالْيَوْمِ الْآخِرِ﴾. الاختلاف واقع بين الأمة وولاة أمرها وعلمائها، ولكن المرجعية عند التنازع يجب أن تكون لكتاب الله وسنة رسوله صلى الله عليه وسلم على فهم سلف الأمة، وليس إلى آراء البشر أو القوانين الوضعية والمناهج المنحرفة كالديمقراطية أو العلمانية.

فإن لم نجد نصاً ظاهراً في مسألة التنازع، رددناه قياساً على كتاب الله وسنة رسوله. والقياس نوعان:
\begin{itemize}
    \item \textbf{قياس صحيح:} وهو القياس الذي فيه إعمال للنصوص الشرعية، كإلحاق فرع بأصل في حكم لعلة جامعة بينهما (مثل القياس في تحديد القبلة أو معرفة العدل).
    \item \textbf{قياس باطل فاسد الاعتبار:} وهو القياس المصادم للنصوص، كقياس إبليس عندما أُمر بالسجود لآدم فقال: ﴿أَنَا خَيْرٌ مِّنْهُ﴾، فقاس مع وجود النص الصريح بالسجود فكان قياساً باطلاً.
\end{itemize}

\section{وجوه السنة مع القرآن الكريم}
السنة النبوية مع كتاب الله عز وجل تأتي على وجوه:
\begin{enumerate}
    \item \textbf{نص موافق:} أن تنزل السنة مؤكدة وموافقة لما جاء به نص القرآن الكريم (كفرضية الصلاة والزكاة).
    \item \textbf{بيان للمجمل:} أن يأتي الفرض مجملاً في القرآن، فتبين السنة كيفيته وتفاصيله عن الله عز وجل، مثل مواقيت الصلوات وعدد ركعاتها، وكيفية الزكاة وأنصبتها.
    \item \textbf{تشريع مستقل:} أن يسن النبي صلى الله عليه وسلم ما ليس فيه نص في كتاب الله (كتحريم الجمع بين المرأة وعمتها أو خالتها، وتحريم كل ذي ناب من السباع).
\end{enumerate}
وهذا الوجه الثالث اختلف فيه العلماء في تأصيله: هل شرعه النبي صلى الله عليه وسلم بما أعطاه الله من الحكمة وبما فرض من طاعته؟ أم أن لكل سنة أصلاً في القرآن يعلمه من يعلمه ويجهله من يجهله؟ أم أنها وحي ألقي في روعه صلى الله عليه وسلم؟ 
والخلاصة: أياً كان التوجيه، فإن طاعة رسول الله صلى الله عليه وسلم فرض محتم، وكل ما جاء به وجب قبوله، ولا عذر لأحد في مخالفة أمر عَرَفَ أنه من أمر رسول الله صلى الله عليه وسلم.

\section{الرد على القرآنيين منكري السنة}
إن من يدعي الاكتفاء بالقرآن ليرد سنة النبي صلى الله عليه وسلم هو في الحقيقة مكذب للقرآن ومتبع لهواه؛ لأن القرآن نفسه أمر بطاعة الرسول وحذر من مخالفته. قال الشافعي: استدللنا بآية ﴿وَأَنزَلَ اللَّهُ عَلَيْكَ الْكِتَابَ وَالْحِكْمَةَ﴾ على أن الحكمة هي سنة رسول الله صلى الله عليه وسلم؛ لأنها قُرنت بالكتاب.
فمن قَبِلَ عن الله فرائضه، يجب أن يقبل عن رسول الله سننه؛ لأن الله هو الذي افترض طاعته. ومن فرّق بينهما فقد ردّ على الله جل وعلا، وسنة رسول الله صلى الله عليه وسلم مبيّنة عن الله مراده.

\section{فتاوى فقهية منثورة في ختام الدرس}
أجاب الشيخ حفظه الله في نهاية الدرس عن جملة من الأسئلة الفقهية، نلخصها فيما يلي:
\begin{itemize}
    \item \textbf{قصر وقت الليل والجمع بين الصلوات:} سُئل عمن يعيش في بلاد لا يتجاوز فيها الليل أربع ساعات، ويشق عليه انتظار العشاء لعمله بالنهار. الجواب: المشقة تجلب التيسير، ويجوز له جمع العشاء مع المغرب جمع تقديم إذا شق عليه ذلك، والأولى أن ينظم وقت نومه ليدرك الصلوات في وقتها خروجاً من الخلاف.
    \item \textbf{التوبة من المال الحرام المستثمر:} سُئل عمن كسب مالاً من تجارة الخمر واشترى به داراً، ثم تاب وأراد استثمار الدار. الجواب: كان الأولى التخلص منه بالكلية إن أمكن، وإذا استثمر هذا المال وجب عليه أن يقدر رأس المال الحرام ويجعله ديناً في ذمته، فيتخلص مما يقابله من أرباح بالصدقة بنية التخلص، وله أن ينفق على نفسه وعياله من الحلال المتبقي قدر الحاجة.
    \item \textbf{الجمع للمسجون المستضعف:} سُئل عمن سُجن في بلاد الكفر ويشق عليه أداء كل صلاة في وقتها بسبب التضييق. الجواب: يجوز له الجمع للمشقة، فـ "المشقة تجلب التيسير" والله لا يريد بعباده العسر.
\end{itemize}
\chapter{مبحث الناسخ والمنسوخ ومكانة السنة}

\section{مقدمة في تعريف النسخ وحكمته}
الحمد لله والصلاة والسلام على رسول الله، أما بعد؛ فقد شرع الإمام الشافعي رحمه الله تعالى في بيان مبحث عظيم من مباحث الأصول، وهو "الناسخ والمنسوخ".

\textbf{النسخ في اللغة:} يأتي بمعنى الإزالة والنقل، فيقال: نسخت الشمس الظل (أي أزالته)، ونسخت الكتاب (أي نقلته). 
\textbf{وفي اصطلاح الأصوليين:} هو رفع حكم متقدم بدليل متراخٍ عنه. والأصل فيه قول الله جل وعلا: ﴿مَا نَنَسَخْ مِنْ آيَةٍ أَوْ نُنسِهَا نَأْتِ بِخَيْرٍ مِّنْهَا أَوْ مِثْلِهَا﴾.

وقد مهّد الشافعي رحمه الله ببيان طلاقة القدرة الإلهية، وأن الله يفعل ما يشاء ويحكم ما يريد لا معقب لحكمه. ومن رحمته بعباده التدرج في التشريع؛ كتحريم الخمر الذي جاء متدرجاً، وكذلك الصيام. فنسخ بعض الأحكام وتغييرها هو من تمام رحمة الله بالعباد للتخفيف عنهم والتوسعة عليهم، كنسخ خمسين صلاة لتصبح خمساً في العمل وخمسين في الأجر. فالحمد لله على نعمه في إثبات الأحكام ونسخها.

\section{مذهب الشافعي في النسخ (السنة لا تنسخ القرآن)}
قرر الشافعي أصلاً عظيماً، وهو أن القرآن لا ينسخه إلا قرآن، والسنة لا ينسخها إلا سنة. ومذهبه أن السنة لا تنسخ القرآن أبداً، ولا يوجد مثال صحيح صريح سالم من المعارضة يثبت أن القرآن نُسخ بسنة.
وإنما السنة تابعة للكتاب، مفسرة ومبينة له. قال تعالى: ﴿قُلْ مَا يَكُونُ لِي أَنْ أُبَدِّلَهُ مِن تِلْقَاءِ نَفْسِي ۖ إِنْ أَتَّبِعُ إِلَّا مَا يُوحَىٰ إِلَيَّ﴾، فالنبي صلى الله عليه وسلم لا ينسخ حكم القرآن ولا يبدله من تلقاء نفسه، بل التشريع ونسخ القرآن من خصائص رب العالمين وحده.

وإذا نُسخت سنة بسنة، فإن النبي صلى الله عليه وسلم يوضح ذلك بسنة أخرى حتى تقوم الحجة، ولا يترك الأمة في لبس. وقد حذر الشافعي من دعوى أن السنة تُنسخ بالقرآن دون وجود سنة توضح ذلك، لئلا يتذرع أصحاب الأهواء (كالقرآنيين أو بعض المقلدين) برد السنن الصحيحة بحجة أنها تخالف القرآن أو أنها منسوخة، كما زعموا في أحكام البيوع المحرمة وحدود الزنا والمسح على الخفين.

\section{علاقة السنة بالقرآن والرد على المتعصبين}
السنة لا تكون أبداً إلا موافقة لكتاب الله عز وجل؛ لأن الكل وحي من الله. ومن ادعى التناقض بينهما فهو مشكك أو جاهل. واتباع السنة هو في الحقيقة اتباع للقرآن؛ لأن القرآن هو الذي أمر باتباعها، فمن قبل عن رسول الله فقد قبل عن الله أمره.

وقد نقل العلامة أحمد شاكر رحمه الله تحذيراً شديداً للمقلدين المتعصبين للمذاهب، الذين يردون الأحاديث الصحاح بدعوى احتمال نسخها نصرة لمذاهبهم (كقول بعضهم: كل حديث خالف مذهبنا فهو منسوخ أو مؤول). وهذا الصنيع يخرج عامة السنن من أيدي الناس، وهو الذي مهد الطريق لاحقاً للقوانين الوضعية والعلمانية لتحل محل الشريعة بدعوى المصلحة. ولا يسع مسلماً تبلغه سنة صحيحة إلا أن يتبعها.

\section{هل النسخ يكون إلى غير بدل؟}
قرر الشافعي، ووافقه جمهور المحققين كالعلامة الشنقيطي، أن النسخ لا يكون إلا إلى بدل، لقوله تعالى: ﴿نَأْتِ بِخَيْرٍ مِّنْهَا أَوْ مِثْلِهَا﴾. فليس ينسخ فرض أبداً إلا أثبت مكانه فرض، كنسخ التوجه لبيت المقدس وإثبات التوجه للكعبة مكانه. 

\section{أمثلة تطبيقية على الناسخ والمنسوخ}

\subsection{نسخ قيام الليل بالصلوات الخمس}
فرض الله قيام الليل أولاً بقوله: ﴿قُمِ اللَّيْلَ إِلَّا قَلِيلًا﴾، ثم نسخه وخففه بقوله: ﴿فَاقْرَءُوا مَا تَيَسَّرَ مِنْهُ﴾. ثم نُسخ ذلك كله وأصبح نافلة بقوله: ﴿وَمِنَ اللَّيْلِ فَتَهَجَّدْ بِهِ نَافِلَةً لَّكَ﴾. وبينت السنة المتواترة (كحديث الأعرابي) أن الله لم يفرض على العباد سوى الصلوات الخمس في اليوم والليلة، وأن ما عداها تطوع، ولا يجوز لأحد أن يقول بوجوب قيام الليل بعد هذا النسخ، وإن كان يستحب استحباباً شديداً.

\subsection{نسخ المصابرة في القتال}
قال تعالى: ﴿إِن يَكُن مِّنكُمْ عِشْرُونَ صَابِرُونَ يَغْلِبُوا مِائَتَيْنِ﴾، فكان الفرض ألا يفر الواحد من المسلمين أمام عشرة من المشركين. ثم خفف الله عنهم لعلمه بضعفهم، فنسخها بقوله: ﴿الْآنَ خَفَّفَ اللَّهُ عَنكُمْ وَعَلِمَ أَنَّ فِيكُمْ ضَعْفًا ۚ فَإِن يَكُن مِّنكُم مِّائَةٌ صَابِرَةٌ يَغْلِبُوا مِائَتَيْنِ﴾، فأصبح الفرض ألا يفر الواحد من المسلمين أمام اثنين من المشركين، والآية شارحة لنفسها ولا تحتاج لتفسير.

\subsection{نسخ حبس الزانيات}
كان حد الزنى في أول الأمر الحبس للمرأة والأذى للرجل: ﴿فَأَمْسِكُوهُنَّ فِي الْبُيُوتِ حَتَّىٰ يَتَوَفَّاهُنَّ الْمَوْتُ أَوْ يَجْعَلَ اللَّهُ لَهُنَّ سَبِيلًا﴾. فبين النبي صلى الله عليه وسلم هذا "السبيل" في حديث عبادة بن الصامت: «خذوا عني، قد جعل الله لهن سبيلا، البكر بالبكر جلد مائة وتغريب عام، والثيب بالثيب جلد مائة والرجم». ثم نُسخ الجلد عن الثيب بالرجم فقط كما فعل النبي صلى الله عليه وسلم مع ماعز والغامدية.

\section{أحكام فقهية مستنبطة (الحائض، المغمى عليه، السكران)}

\subsection{قضاء الصلاة للمرأة الحائض}
الحيض أمر جبلّي خلقه الله في المرأة، فهي ليست عاصية بوجوده. وقد أسقط الله عنها فرض الصلاة أيام حيضها. ولما كانت معذورة بغير جناية منها، أجمع العلماء على أنها لا تقضي الصلاة، ولكنها تقضي الصيام (لأن الصيام خفيف يقع شهراً في السنة، بخلاف الصلاة التي تتكرر يومياً فيشق قضاؤها).

\subsection{صلاة المغمى عليه والسكران}
المغمى عليه (أو المغلوب على عقله بمرض وعارض سماوي) لا يقضي الصلاة إذا خرج وقتها وهو فاقد للوعي؛ قياساً على الحائض، لأنه فقد مناط التكليف (العقل) بغير جناية منه ولا فعل.
أما "السكران"، فإذا فاتته الصلاة وجب عليه قضاؤها إذا أفاق، ولا يُقاس على المغمى عليه؛ لأن السكران أدخل المانع على نفسه متعمداً، فهو عاصٍ بفعله، والجهة غير منفكة، ولذا تبطل صلاته حال سكره لقوله تعالى: ﴿لَا تَقْرَبُوا الصَّلَاةَ وَأَنتُمْ سُكَارَىٰ﴾، ويلزمه القضاء بعد إفاقته عقوبة له.

\section{فقه إقامة الحدود}
ساق الشافعي حديث العسيف (الأجير) الذي زنى بامرأة مستأجره، فافتداه أبوه بمائة شاة وجارية ظناً منه أنه يُرجم، ثم سأل العلماء فاخبروه أن على ابنه جلد مائة وتغريب عام. فارتفعوا للنبي صلى الله عليه وسلم، فأبطل الفدية ورد الشياه والجارية، وأمر بجلد الشاب وتغريبه. ثم أرسل أُنَيْسًا الأسلمي للمرأة: «فإن اعترفت فارجمها».
وفي هذا دلالة عظيمة على كمال الشريعة؛ فإقامة الحدود لا تُبنى على الشبهات، بل النبي صلى الله عليه وسلم لم يكتفِ بإقرار الشاب والزوج، بل أرسل للمرأة لتُقر بنفسها، فلما اعترفت رُجمت، ولو أنكرت لَدُرئ عنها الحد. وهذا يرد على شبهات العلمانيين الذين يصورون الشريعة كأنها متعطشة للدماء، بل هي غاية في الرحمة والاحتياط للنفوس والأعراض.
\chapter{الناسخ والمنسوخ بدلالة السنة والإجماع وأحكام اللعان}

\section{الناسخ والمنسوخ الذي تدل عليه السنة والإجماع}
الحمد لله والصلاة والسلام على رسول الله، أما بعد؛ فقد شرع الإمام الشافعي رحمه الله تعالى في بيان قسم آخر وهو: "الناسخ والمنسوخ الذي تدل عليه السنة والإجماع". ومقصوده أن هناك آيات وأحاديث دلت السنة المطهرة والإجماع على أنها منسوخة.
ويجب التنبيه هنا إلى قاعدة أصولية هامة: أن الإجماع في ذاته لا يَنْسَخ ولا يُنْسَخ؛ لأنه لا إجماع إلا بعد موت رسول الله صلى الله عليه وسلم وانقطاع الوحي. وإنما الإجماع يَدُل على النسخ، أي يُجمع العلماء على أن هذا النص قد نُسخ بنص آخر كان مستقراً في عهد النبوة.

\subsection{نسخ الوصية للوارث بآيات المواريث}
قال الله تبارك وتعالى: ﴿كُتِبَ عَلَيْكُمْ إِذَا حَضَرَ أَحَدَكُمُ الْمَوْتُ إِن تَرَكَ خَيْرًا الْوَصِيَّةُ لِلْوَالِدَيْنِ وَالْأَقْرَبِينَ﴾، وقال في حق الزوجة: ﴿وَالَّذِينَ يُتَوَفَّوْنَ مِنكُمْ وَيَذَرُونَ أَزْوَاجًا وَصِيَّةً لِّأَزْوَاجِهِم﴾. فهاتان الآيتان في سورة البقرة تثبتان الوصية للوالدين وللزوجة.
ثم أنزل الله عز وجل آيات المواريث في سورة النساء مفصلة لأنصبة الورثة. فاحتمل الأمر أن تكون الوصية ثابتة مع الميراث، واحتمل أن تكون آيات المواريث ناسخة لآيات الوصية.
فطلب أهل العلم الدلالة من كتاب الله فلم يجدوا نصاً يقطع بالمراد، فطلبوه في سنة رسول الله صلى الله عليه وسلم، فوجدوا أن أهل المغازي والفتيا نقلوا عامة عن عامة أن النبي صلى الله عليه وسلم قال في خطبته عام الفتح: «لا وصية لوارث». وأجمع أهل العلم قاطبة على ذلك. فدل هذا الحديث مع الإجماع المذكور على أن آيات النساء ناسخة لآيات البقرة، وأن الميراث ناسخ للوصية بالنسبة للوارث.

\subsection{حكم الوصية لغير الوارث ولغير الأقارب}
إذا بطلت الوصية للوارث، فهل تجوز لغير الوارث من الأقارب وغيرهم؟ ذهب الإمام طاوس رحمه الله إلى القول بأن الوصية نسخت للوالدين وثبتت للقرابة غير الوارثين فقط، فمن أوصى لغير قرابة لم تجز وصيته.
فرد الشافعي هذا القول بالدليل، واستدل بحديث عمران بن حصين رضي الله عنه: أن رجلاً من الأنصار أعتق ستة مماليك له عند موته ليس له مال غيرهم، فجزأهم النبي صلى الله عليه وسلم أثلاثاً، فأقرع بينهم، فأعتق اثنين وأرق أربعة.
ووجه الدلالة من الحديث: أن إعتاقه لهم عند الموت وصية، وهم مماليك ليسوا بأقارب له، فدل على جواز الوصية لغير الأقارب، ودل أيضاً على إبطال الاستسعاء في حقهم، وأن الوصية مقتصرة شرعاً على حدود "الثلث" فقط (لأنه أعتق الثلث وأرق الباقي). 

\section{الفرائض التي أنزلها الله نصاً (القذف واللعان)}
شرع الشافعي بعد ذلك في بيان باب الفرائض التي نص عليها القرآن ولا تحتاج إلى بيان إضافي من السنة لكفايتها، ومثّل لذلك بحد القذف. قال تعالى: ﴿وَالَّذِينَ يَرْمُونَ الْمُحْصَنَاتِ ثُمَّ لَمْ يَأْتُوا بِأَرْبَعَةِ شُهَدَاءَ فَاجْلِدُوهُمْ ثَمَانِينَ جَلْدَةً﴾. والمحصنات هنا هن: الحرائر البوالغ العفائف.
فهذا حكم رمي الأجنبي للمرأة بالزنا، أنه يُجلد ثمانين جلدة إن لم يأت بأربعة شهداء، ولا يقبل منه غير ذلك.

\subsection{أحكام اللعان وما زادته السنة على القرآن}
أما إن كان القاذف هو الزوج، فقد فرق القرآن بين حكمه وحكم الأجنبي، فجعل للزوج مخرجاً باللعان: ﴿وَالَّذِينَ يَرْمُونَ أَزْوَاجَهُمْ وَلَمْ يَكُن لَّهُمْ شُهَدَاءُ إِلَّا أَنفُسُهُمْ فَشَهَادَةُ أَحَدِهِمْ أَرْبَعُ شَهَادَاتٍ بِاللَّهِ إِنَّهُ لَمِنَ الصَّادِقِينَ﴾ الآيات.
فإذا تلاعنا، درأ ذلك عنهما الحد. وقد بيّن الشافعي أن القرآن اكتفى ببيان صفة اللعان وعدد الأيمان (أربع شهادات والخامسة اللعنة أو الغضب)، فلم يحتج الصحابة لنقل لفظ النبي صلى الله عليه وسلم في تلقينهم الأيمان لكونها مفصلة وواضحة في التنزيل. 
لكنهم نقلوا في المقابل الأحكام التي زادتها السنة ولم تُذكر في القرآن، وهي: التفريق بين المتلاعنين، ونفي الولد عن الزوج وإلحاقه بالأم. دلالة على أن السنة تستقل بتشريع الأحكام.

\section{الحكم بالظاهر لا بالبواطن والسرائر}
من أعظم القواعد القضائية والمنهجية التي أرساها النبي صلى الله عليه وسلم في قصة المتلاعنين (عويمر العجلاني وزوجته): الحكم بالظاهر. 
فقد قال صلى الله عليه وسلم في شأنهما: «إن أحدكما كاذب»، وذكر علامات خِلْقية ظاهرة في الولد، فقال: «إن جاءت به أُحَيْمِرَ كأنه وحَرَةٌ... فلا أراه إلا قد صدق عليها»، فلما جاءت به على النعت المكروه (الذي يثبت التهمة)، قال صلى الله عليه وسلم: «لولا الأيمان لكان لي ولها شأن»، أي لولا أن الله حكم بدرء الحد عنها باللعان لأقام عليها الحد بسبب وجود الشبه بالرجل المتهَم.

فلم يحكم النبي صلى الله عليه وسلم بالعلامة الظاهرة ولا بما يعلمه من القرائن، بل أمضى حكم الله النافذ باللعان، مغلباً للظاهر الشرعي.
وهذا يوجب على الحكام والقضاة ألا يقضوا إلا بالظاهر (بالبينة أو الإقرار)، ولا يقضوا بعلمهم الشخصي ولا بالقرائن المجردة عن البينة. كما قال صلى الله عليه وسلم في الحديث الصحيح: «إنما أنا بشر، وإنكم تختصمون إلي، ولعل بعضكم أن يكون ألحن بحجته من بعض، فأقضي له على نحو ما أسمع، فمن قضيت له من حق أخيه شيئاً فلا يأخذه، فإنما أقطع له قطعة من النار». 
فالشريعة جاءت لتبني الأحكام على الظواهر المنضبطة، وتَكِل السرائر إلى الخبير بها سبحانه وتعالى.
\chapter{الفرائض المنصوصة وتفصيل السنة لمجمل القرآن}

\section{الفرائض المنصوصة التي سن رسول الله معها}
الحمد لله رب العالمين، والصلاة والسلام على رسول الله، أما بعد؛
يواصل الإمام الشافعي رحمه الله بيان موقع السنة من القرآن، ويذكر هنا أمثلة للفرائض التي نزل أصلها في القرآن، وجاءت السنة المطهرة لتوضحها وتبين كيفياتها وشروطها؛ إذ لو اكتفى الناس بتدبر القرآن لعرفوا الفرض، ولكن السنة أتت لتفصل وتبيّن.

\subsection{أحكام الوضوء والغسل (التفريق بين الفرض والسنة)}
قال الله تعالى: ﴿إِذَا قُمْتُمْ إِلَى الصَّلَاةِ فَاغْسِلُوا وُجُوهَكُمْ وَأَيْدِيَكُمْ إِلَى الْمَرَافِقِ وَامْسَحُوا بِرُءُوسِكُمْ وَأَرْجُلَكُمْ إِلَى الْكَعْبَيْنِ ۚ وَإِن كُنتُمْ جُنُبًا فَاطَّهَّرُوا﴾.
أبان القرآن أن طهارة الجنب تكون بالغسل، وأن طهارة الحدث الأصغر تكون بالوضوء. وسنّ رسول الله صلى الله عليه وسلم في غسل الجنابة أن يتوضأ المغتسل وضوء الصلاة أولاً، فهذا من سنته مع الغسل المنصوص.

وفي الوضوء، احتمل قوله تعالى: ﴿فَاغْسِلُوا وُجُوهَكُمْ﴾ غسلة واحدة أو أكثر، فجاءت السنة لتبيّن أن النبي صلى الله عليه وسلم توضأ (مرة مرة)، وتوضأ (مرتين مرتين)، وتوضأ (ثلاثاً ثلاثاً).
فدلت سنته حين اكتفى بالمرة الواحدة على أن الفرض الذي يجزئ في الوضوء هو المرة الواحدة، وأن ما زاد عليها (الثنتان والثلاث) هو نافلة وسنة مستحبة لطلب الفضل، وليس بفرض يفسد الوضوء بتركه.

كما بيّنت السنة دخول المرفقين والكعبين في الغسل، وأن الغسل يعمها؛ فقد كانت الآية تحتمل الغسل إليهما أو غسلهما، فبيّن فعله صلى الله عليه وسلم دخولهما.

\textbf{فائدة لغوية:} الوَضوء (بفتح الواو) هو الماء الذي يُتوضأ به، والوُضوء (بضم الواو) هو الفعل وحركة غسل الأعضاء.

\subsection{موانع الإرث (تخصيص السنة لعموم المواريث)}
قال تعالى: ﴿يُوصِيكُمُ اللَّهُ فِي أَوْلَادِكُمْ ۖ لِلذَّكَرِ مِثْلُ حَظِّ الْأُنثَيَيْنِ﴾، وقال في الكلالة (وهو من مات وليس له أصل ولا فرع وارث): ﴿وَإِن كَانَ رَجُلٌ يُورَثُ كَلَالَةً أَوِ امْرَأَةٌ وَلَهُ أَخٌ أَوْ أُخْتٌ فَلِكُلِّ وَاحِدٍ مِّنْهُمَا السُّدُسُ﴾.

هذه الآيات وردت بلفظ عام يشمل كل والد وولد وكل زوج وزوجة. فهل يرث القاتل عمداً? وهل يتوارث المسلم والكافر? وهل يرث العبد?
أتت السنة المطهرة لتخصص هذا العموم بضوابط وموانع، وبيّنت أن الميراث لا يثبت إلا بشروط:
\begin{itemize}
    \item \textbf{اتحاد الدين:} قال صلى الله عليه وسلم: «لا يَرِثُ الْمُسْلِمُ الْكَافِرَ وَلَا الْكَافِرُ الْمُسْلِمَ». فلا يتوارث أهل ملتين.
    \item \textbf{الحرية:} لأن العبد وما يملك هو ملك لسيده، فهو يباع ويوهب، فلا يرث ولا يُورث، ولو أورثناه لأعطينا المال لسيده الذي لا فرض له في كتاب الله.
    \item \textbf{البراءة من القتل العمد:} قال صلى الله عليه وسلم: «ليس لقاتل شيء»، فمن استعجل الشيء قبل أوانه عوقب بحرمانه، ويُمنع القاتل عمداً من الميراث عقوبة له.
\end{itemize}
فدلت السنة على أن الله تعالى إنما أراد بالآيات الخاصة بالميراث أناساً مخصوصين تتوفر فيهم هذه الشروط، دون غيرهم ممن شملهم اللفظ العام.

\subsection{تخصيص البيوع المحرمة من عموم (أحل الله البيع)}
قال تعالى: ﴿وَأَحَلَّ اللَّهُ الْبَيْعَ وَحَرَّمَ الرِّبَا﴾. هل كل البيوع حلال? 
لقد نهى رسول الله صلى الله عليه وسلم عن بيوع تراضى بها المتبايعان، ولم يكن فيها غرر مجهول، ولكنها تضمنت الربا المحرم. كأن يبيع الرجل الذهب القديم بذهب جديد ويدفع الفرق، فهذا بيع بالتراضي ولكنه محرم؛ لقوله صلى الله عليه وسلم: «الذهب بالذهب مثلاً بمثل ويداً بيد».

فدلت السنة الكريمة على أن الله تعالى إنما أحل البيوع التي لم يحرمها على لسان نبيه، فخصصت السنة عموم الآية، وبيّنت أن من البيوع ما هو باطل وإن تراضى عليه الناس.

\section{بيان السنة لمجمل الفرائض}
أحكم الله فرضه في كتابه مجملاً، كالصلاة والزكاة والحج، ثم بيّن على لسان نبيه صلى الله عليه وسلم كيفيات هذا الفرض:

\subsection{تفصيل أحكام الصلاة ومواقيتها}
قال تعالى: ﴿إِنَّ الصَّلَاةَ كَانَتْ عَلَى الْمُؤْمِنِينَ كِتَابًا مَّوْقُوتًا﴾. فأجمل الفرض. 
ثم بيّن النبي صلى الله عليه وسلم مواقيت الصلاة، وعددها بأنها خمس صلوات في اليوم والليلة لا يزاد عليها في الفرض. وبيّن أعداد ركعاتها، وأن الصبح ركعتان يجهر فيهما، والظهر والعصر أربع يسر فيهما، والمغرب ثلاث يجهر في اثنتين، والعشاء أربع يجهر في اثنتين. وبيّن أن الدخول فيها بالتكبير والخروج منها بالتسليم، وسنّ قصر الرباعية في السفر.

\subsection{صلاة الخسوف وصلاة الخوف}
وسنّ النبي صلى الله عليه وسلم في صلاة الكسوف ركعتين في كل ركعة ركوعين وسجودين، كما روى ذلك ابن عباس وعائشة رضي الله عنهما باختلاف يسير في الألفاظ واتفاق في الهيئة.

وكذلك في صلاة الخوف؛ فرغم أن الآية ذكرت صلاة الخوف مجملة: ﴿وَإِذَا كُنتَ فِيهِمْ فَأَقَمْتَ لَهُمُ الصَّلَاةَ فَلْتَقُمْ طَائِفَةٌ مِّنْهُم مَّعَكَ﴾، فقد صلاها النبي صلى الله عليه وسلم بهيئات متعددة تناسب حال الحرب وموقع العدو. ففي غزوة الخندق أخّر الصلاة حتى كُفي القتال فصلاها جمعاً، فلما نزلت صلاة الخوف نُسخ التأخير، وصار الواجب أداء الصلاة في وقتها بحسب الاستطاعة (رجالاً أو ركباناً، مستقبلي القبلة أو غير مستقبليها) في حال المسايفة واشتداد الخوف، ولا تُترك الصلاة أبداً لتُقضى بعد خروج وقتها إلا للنائم والناسي.

\section{خاتمة: مكانة السنة والرد على منكريها}
هذا الباب من أجلّ ما يُرد به على الطائفة المنحرفة التي تسمي نفسها بـ "القرآنيين"، وهم في الحقيقة مكذبون للقرآن؛ إذ لو اتبعوا القرآن حقاً لاتبعوا من أمره القرآن بطاعته. فالقرآن أتى بأصول الفرائض مجملة، ولولا سنة رسول الله صلى الله عليه وسلم ما عرف المسلمون كيف يصلون، ولا متى يزكون، ولا كيف يحجون، فمن ردّ السنة فقد هدم الدين كله وانسلخ منه.
\chapter{جمل الفرائض: الزكاة والحج وأحكام النساء والطعام}

\section{جمل الفرائض في الزكاة وتفصيل السنة لها}
الحمد لله والصلاة والسلام على رسول الله، أما بعد؛ فقد شرع الإمام الشافعي رحمه الله تعالى في بيان جمل الفرائض التي نزل أصلها في القرآن الكريم وأتت تفصيلاتها وتفريعاتها في السنة النبوية المطهرة. وابتدأ بالزكاة بعد الصلاة.

الزكاة في اللغة: الطهارة والنماء. وفي الشرع: أخذ مال مخصوص من أشخاص مخصوصين في أوقات معلومة لأصناف معلومة.
قال الله جل وعلا: ﴿وَأَقِيمُوا الصَّلَاةَ وَآتُوا الزَّكَاةَ﴾، وقال في ذم الكفار: ﴿الَّذِينَ هُمْ يُرَاءُونَ * وَيَمْنَعُونَ الْمَاعُونَ﴾، فسر بعض السلف (كعلي وابن عباس رضي الله عنهم) الماعون هنا بالزكاة المفروضة. 

وقال تعالى: ﴿خُذْ مِنْ أَمْوَالِهِمْ صَدَقَةً تُطَهِّرُهُمْ وَتُزَكِّيهِمْ بِهَا﴾. مخرج هذه الآية عام، فيحتمل أن تؤخذ الزكاة من جميع الأموال، فجاءت السنة المطهرة فدلت على أن الزكاة في بعض الأموال دون بعض، وحددت الأنصبة والمقادير:

\subsection{زكاة الماشية}
لما كانت الماشية أصنافاً، أخذ رسول الله صلى الله عليه وسلم الزكاة من الإبل والبقر والغنم، ولم يأخذ من الخيل والبغال والحمير، وقال: «ليس على المسلم في عبده ولا فرسه صدقة». وحدد أنصابها؛ فجعل زكاة الإبل تبتدئ من خمس، والبقر من ثلاثين، والغنم من أربعين، فاستدللنا بسنته على وجوب الزكاة فيما أخذ منه وسقوطها عما ترك.

\subsection{زكاة الزروع والثمار}
أخذ رسول الله صلى الله عليه وسلم الزكاة من النخل والعنب بالخَرْص (وهو تقدير الخبير لما على الشجر من ثمر)، وجعل فيها العشر إن سقيت بماء السماء أو العيون، ونصف العشر إن سقيت بالدلو والمشقة. ولم يأخذ من باقي الشجر (كالتين والجوز واللوز).
كما أخذ الزكاة من الحنطة (القمح) والشعير والذرة وكل ما يقتات ويدخر، ولم يأخذ مما لا يقتات ويدخر (كحب الرشاد والكزبرة).

\subsection{زكاة النقدين والركاز}
وفرض صلى الله عليه وسلم الزكاة في الورِق (الفضة) والذهب لأنهما نقد الناس وثمن مبيعاتهم، ولم يأخذ من النحاس والحديد والرصاص شيئاً. وبين أن زكاة النقدين والماشية تؤخذ مرة كل حَول (سنة)، بينما زكاة الزروع تؤخذ يوم الحصاد لقوله تعالى: ﴿وَآتُوا حَقَّهُ يَوْمَ حَصَادِهِ﴾. وسنّ في الركاز (ما استخرج من دفائن الأرض) الخُمس في يوم خروجه.

\section{جمل الفرائض في الحج ومناسك السنة}
قال تعالى: ﴿وَلِلَّهِ عَلَى النَّاسِ حِجُّ الْبَيْتِ مَنِ اسْتَطَاعَ إِلَيْهِ سَبِيلًا﴾. فذكر الفرض مجملاً.
فبين رسول الله صلى الله عليه وسلم أن "السبيل" هو الزاد والراحلة، وعلمنا كيف نلبي، وما يلبس المحرم وما يجتنبه (كالطيب والثياب المخيطة)، وعلمنا أعمال عرفة ومزدلفة والرمي والحلق والطواف؛ مردداً: «خذوا عني مناسككم».
فلو أن امرأً لم يعلم من السنة إلا ما فصله النبي صلى الله عليه وسلم لمجمل القرآن لكانت الحجة قائمة عليه بوجوب طاعة رسول الله صلى الله عليه وسلم، ووجب عليه أن يعلم أن سنة رسول الله لا تخالف كتاب الله أبداً.

\section{أحكام العدد والمحرمات من النساء}

\subsection{عدة المتوفى عنها زوجها الحامل}
قال تعالى: ﴿وَالَّذِينَ يُتَوَفَّوْنَ مِنْكُمْ وَيَذَرُونَ أَزْوَاجًا يَتَرَبَّصْنَ بِأَنْفُسِهِنَّ أَرْبَعَةَ أَشْهُرٍ وَعَشْرًا﴾، وقال تعالى: ﴿وَأُولَاتُ الْأَحْمَالِ أَجَلُهُنَّ أَنْ يَضَعْنَ حَمْلَهُنَّ﴾.
فظن بعض أهل العلم أن الحامل المتوفى عنها زوجها تعتد بأبعد الأجلين (تجمع بين العدتين). فجاءت السنة لتبين الأمر؛ فلما وضعت سُبيعة الأسلمية رضي الله عنها حملها بعد وفاة زوجها بأيام قليلة، قال لها رسول الله صلى الله عليه وسلم: «قد حللت فتزوجي». فدلت السنة على أن وضع الحمل ينهي العدة مطلقاً وإن كان بعد سويعات من الوفاة.

\subsection{المحرمات من النساء (الجمع بين المرأة وعمتها)}
قال تعالى معدداً المحرمات في النكاح: ﴿حُرِّمَتْ عَلَيْكُمْ أُمَّهَاتُكُمْ وَبَنَاتُكُمْ...﴾ ثم قال: ﴿وَأُحِلَّ لَكُمْ مَا وَرَاءَ ذَلِكُمْ﴾.
فكان ظاهر الآية جواز نكاح أي امرأة لم تذكر في الآية. لكن السنة المطهرة بينت أن قوله ﴿وَأُحِلَّ لَكُمْ مَا وَرَاءَ ذَلِكُمْ﴾ مقيد بما لم تحرمه السنة أيضاً. فقد نهى النبي صلى الله عليه وسلم أن يجمع الرجل بين المرأة وعمتها، وبين المرأة وخالتها. والتحريم هنا ليس تحريماً مؤبداً لذات العمة أو الخالة كتحريم الأم والأخت، بل هو تحريم "الجمع" بينهما في عصمة واحدة فقط.

\subsection{حتى تنكح زوجاً غيره (اشتراط الوطء)}
قال تعالى في المطلقة ثلاثاً: ﴿فَإِنْ طَلَّقَهَا فَلَا تَحِلُّ لَهُ مِنْ بَعْدُ حَتَّى تَنْكِحَ زَوْجًا غَيْرَهُ﴾. احتمل لفظ "النكاح" العقد المجرد، واحتمل العقد مع الإصابة (الجماع). فجاءت السنة موضحة لمراد الله؛ حينما جاءت امرأة رفاعة القرظي تشتكي أن زوجها الثاني لم يصبها، فقال لها النبي صلى الله عليه وسلم: «أتريدين أن ترجعي إلى رفاعة؟ لا، حتى تذوقي عُسَيْلَتَهُ ويذوق عُسَيْلَتَكِ». فبينت السنة أن التحليل لا يتم إلا بنكاح صحيح يتبعه وطء.

\section{ما أحل وحرم من الطعام}
قال تعالى: ﴿قُلْ لَا أَجِدُ فِي مَا أُوحِيَ إِلَيَّ مُحَرَّمًا عَلَى طَاعِمٍ يَطْعَمُهُ إِلَّا أَنْ يَكُونَ مَيْتَةً أَوْ دَمًا مَسْفُوحًا أَوْ لَحْمَ خِنْزِيرٍ...﴾.
ظاهر الآية حصر المحرمات في هذه المذكورات، ولكن السنة المطهرة جاءت بتحريم "كل ذي ناب من السباع، وكل ذي مخلب من الطير". فدلت السنة إما على أن الحصر في الآية كان لما سألوا عنه حينها، أو أنه مخصوص، فما حرمه رسول الله فهو حرام كأنما نزل في القرآن.

\section{ما تجتنبه المعتدة للوفاة}
قال تعالى في المعتدة للوفاة: ﴿فَإِذَا بَلَغْنَ أَجَلَهُنَّ فَلَا جُنَاحَ عَلَيْكُمْ فِيمَا فَعَلْنَ فِي أَنْفُسِهِنَّ بِالْمَعْرُوفِ﴾.
القرآن فرض عليها أن تتربص أربعة أشهر وعشراً، ولم يذكر الأشياء التي تجتنبها المعتدة، فجاءت السنة النبوية فألزمتها بأمرين زائدين على مجرد التربص عن الأزواج: 
1. الامتناع عن الزينة والطيب وما يدعو للرغبة فيها.
2. لزوم السكن والمكث في بيت زوجها الذي مات وهي فيه، فلا تخرج منه إلا لحاجة ضرورية.

\section{خاتمة: العذر بالجهل لمن خالف السنة من العلماء}
يقرر الشافعي أصلاً منهجياً عظيماً: كل عالم رُوي عنه قول يخالف سنة صحيحة لرسول الله صلى الله عليه وسلم، فإنه معذور مأجور؛ لأنه إنما خالفها لجهله بها وعدم بلوغها إياه، ولو علمها لرجع إليها وترك قوله. ولا يجوز لمسلم إذا بلغته السنة أن يتركها لقول عالم مهما علا شأنه؛ فالحجة في قول رسول الله صلى الله عليه وسلم، ومن أطاع الرسول فقد أطاع الله.
\chapter{العلل في الأحاديث واختلاف الروايات}

\section{مقدمة في اختلاف الأحاديث وموافقة السنة للقرآن}
الحمد لله والصلاة والسلام على رسول الله، أما بعد؛ فقد عقد الإمام الشافعي رحمه الله تعالى باباً عظيماً أسماه «باب العلل في الأحاديث»، والعلة في اصطلاح المحدثين: سبب خفي قادح يقدح في صحة الحديث مع أن ظاهره السلامة، والتقييم هنا أعم من العلة الاصطلاحية، فهو يقصد أسباب الاختلاف الظاهري بين الأحاديث.

وقد سأل سائل الإمام الشافعي عن أحاديث ظاهرها الاختلاف والتضاد؛ منها ما يبين القرآن، ومنها ما يزيد عليه، ومنها ما ينهى عنه النبي صلى الله عليه وسلم فيعده العلماء تحريماً وفي موضع آخر يعدونه كراهة، ومنها ما يؤخذ به ويترك ما يعارضه. فما هي الحجة والقاعدة في هذا الأخذ والترك والترجيح؟

أجاب الشافعي بتأصيل عظيم: «كل ما سن رسول الله صلى الله عليه وسلم مع كتاب الله، فهو موافق لكتاب الله». فإما أن تأتي السنة بنص يوافق نص القرآن، وإما أن تبين مجمل القرآن وتفصله، وإما أن تأتي بسنة ليس فيها نص في القرآن. وفي كل هذه الحالات، يجب التسليم لأن الله هو الذي فرض علينا طاعة رسوله صلى الله عليه وسلم.

\section{الناسخ والمنسوخ في السنة}
كما أن القرآن ينسخ بعضه بعضاً (كنسخ المصابرة من واحد لعشرة إلى واحد لاثنين)، كذلك سنة رسول الله صلى الله عليه وسلم ينسخ بعضها بعضاً. ومثاله قوله صلى الله عليه وسلم: «كنت نهيتكم عن زيارة القبور فزوروها»، وقوله: «كنت نهيتكم عن ادخار لحوم الأضاحي بعد ثلاث فكلوا وادخروا». فالسنة ينسخ بعضها بعضاً لأن الكل وحي من عند الله عز وجل.

\section{أسباب الاختلاف الظاهري بين الأحاديث}
لا يوجد في أحاديث رسول الله صلى الله عليه وسلم الصحيحة اختلاف تضاد أبداً، وإنما يقع الاختلاف في فهم السامع لعدة أسباب بينها الشافعي:

\subsection{العام المراد به الخصوص}
رسول الله صلى الله عليه وسلم عربي اللسان والدار، وقد يقول القول عاماً ويريد به العام، وقد يطلق اللفظ عاماً ويريد به الخاص، فمن لم يفهم لغة العرب توهم التعارض.

\subsection{سماع الجواب دون السؤال}
قد يُسأل النبي صلى الله عليه وسلم عن شيء فيجيب على قدر المسألة، فيحضر الراوي فيسمع الجواب ولا يدرك السؤال ومناطه، فيحدث بالجواب مطلقاً. كمن روى قوله: «هو الطهور ماؤه الحل ميتته» دون أن يذكر أنه كان جواباً لمن سأل عن الوضوء بماء البحر لتوفير الماء العذب للشرب.

\subsection{اختلاف الحالين}
قد يسن النبي صلى الله عليه وسلم في شيء حكماً، وفيما يخالفه حكماً آخر، فلا يفرق بعض السامعين بين اختلاف الحالين اللتين سن فيهما. 
ومثاله بيوع العرايا؛ فقد نهى النبي صلى الله عليه وسلم عن الغرر والجهالة والمزابنة، ولكنه رخص لبعض الفقراء الذين لا يملكون نقوداً أن يشتروا الرطب على النخل بخرصه من التمر ليأكلوه (بيوع العرايا). فمن لم يفرق بين الحالين ظن التعارض، وإنما العرايا حالة رخصة مخصوصة لحاجة، والنهي العام باقٍ على التحريم.

\subsection{حفظ المنسوخ دون الناسخ}
قد يحفظ بعض الصحابة الحديث المنسوخ ويغيب عنه الناسخ، فيحدث بما حفظ. ويحفظ صحابي آخر الناسخ فيحدث به. وهذا لا يذهب على عامة الصحابة، فلا تضيع سنة ناسخة عن مجموع الأمة أبداً، فإذا جمعت طرق الحديث تبين الناسخ من المنسوخ.

\section{وجوب التسليم للسنة وترك الاعتراض}
يقرر الشافعي قاعدة إيمانية جليلة: «لا يجوز أن تقول لسنة رسول الله ﷺ: لِمَ ولا كيف؟». 
فلا يجوز لمسلم أن يفرق بين ما جمعه النبي صلى الله عليه وسلم، ولا أن يجمع بين ما فرقه. ومن اعترض على حكم رسول الله صلى الله عليه وسلم متسائلاً على وجه الإنكار "لماذا فرق بين هذا وهذا؟" فهو إما جاهل، أو مرتاب شاك، والشك والارتياب في كلام النبي صلى الله عليه وسلم شر من الجهل وقد يؤدي للكفر.
فالمؤمن يقول: «سمعنا وأطعنا»، ويسلم لحكم الله ورسوله دون حرج. أما ادعاء بعض المنتكسين في أيامنا أنهم "سلفيون عقلانيون" يقدمون العقل على النقل، فهو طعن في السلفية وبوابة للعلمانية، فالعقل الصريح لا يعارض النقل الصحيح أبداً.

\section{طرق دفع التعارض (الترجيح)}
إذا وجدنا حديثين ظاهرهما الاختلاف، فإننا نجمع بينهما بالوجوه السابقة، فإن لم يمكن، صرنا إلى الترجيح بناءً على أسس منها:
\begin{itemize}
    \item أن يكون أحد الحديثين له دلالة ومواكبة من كتاب الله أو سنة نبيه تدعمه.
    \item أن يكون أحد الحديثين يوافق القواعد الكلية والأصول الثابتة.
    \item أن يكون رواته أحفظ وأكثر ضبطاً.
\end{itemize}
والأصل أن ما نهى عنه رسول الله صلى الله عليه وسلم فهو على التحريم، حتى تأتي دلالة تصرفه إلى الكراهة أو التنزيه.
\chapter{السنن الناسخة والمنسوخة وتخصيص الأحكام}

\section{نسخ النهي عن ادخار لحوم الأضاحي}
الحمد لله والصلاة والسلام على رسول الله، أما بعد؛ فقد استكمل الإمام الشافعي رحمه الله تعالى الحديث عن موقع السنة من القرآن، وبدأ يفصل في "السنن الناسخة والمنسوخة". فكما أن القرآن ينسخ بعضه بعضاً، فكذلك سنة رسول الله صلى الله عليه وسلم ينسخ بعضها بعضاً.

ومن ذلك نهي النبي صلى الله عليه وسلم عن ادخار لحوم الأضاحي بعد ثلاث. فقد روى ابن عمر وعلي بن أبي طالب رضي الله عنهما أن رسول الله صلى الله عليه وسلم نهى عن أكل لحوم الضحايا بعد ثلاث. وقد حفظا المنسوخ وحدثا به، ولم تبلغهما الرخصة الناسخة.
بينما حفظت أم المؤمنين عائشة رضي الله عنها المنسوخ والناسخ وسبب التشريع فقالت: «دَفَّ ناسٌ من أهل البادية حضرة الأضحى في زمان النبي صلى الله عليه وسلم (والدافة: هم القوم يسيرون جماعة ونزلوا ضيوفاً)، فقال النبي صلى الله عليه وسلم: ادخروا لثلاث وتصدقوا بما بقي. فلما كان بعد ذلك قيل يا رسول الله لقد كان الناس ينتفعون بضحاياهم يُجملون منها الودك (الشحم) فقال: إنما نهيتكم من أجل الدافة التي دفت حضرة الأضحى، فكلوا وتصدقوا وادخروا».

فحديث عائشة رضي الله عنها من أبين ما يوجد في الناسخ والمنسوخ، وهو يوضح أن بعض الرواة قد يحفظ المنسوخ ولا يبلغه الناسخ، ولا تضيع السنة بتمامها عن مجموع الأمة. 
وقد رجّح المحققون (كالعلامة أحمد شاكر) أن نهي النبي صلى الله عليه وسلم هنا لم يكن تشريعاً عاماً منسوخاً، بل كان تصرفاً منه على سبيل نظر الإمام والحاكم لمصلحة الرعية لوجود المجاعة (الدافة)، فلما زالت العلة زال النهي.

\section{صلاة الخوف ونسخ تأخير الصلاة عن وقتها}
من أمثلة الناسخ والمنسوخ في السنة: تأخير الصلاة عن وقتها في شدة الخوف.
فقد روى أبو سعيد الخدري رضي الله عنه أن المشركين شغلوا النبي صلى الله عليه وسلم يوم الخندق عن صلاة الظهر والعصر والمغرب حتى دخل الليل، فأمر بلالاً فأقام وصلاها تباعاً. وكان هذا قبل نزول شرعية (صلاة الخوف).

فلما نزل قوله تعالى: ﴿وَإِذَا كُنتَ فِيهِمْ فَأَقَمْتَ لَهُمُ الصَّلَاةَ...﴾ شرع النبي صلى الله عليه وسلم صلاة الخوف بكيفيات متعددة (كصلاة ذات الرقاع وعسفان)، فصارت الصلاة تؤدى في وقتها على أي هيئة أمكنت (رجالاً أو ركباناً، مستقبلي القبلة أو غير مستقبليها في حال المسايفة). وبذلك نُسخت رخصة تأخير الصلاة عن وقتها لعذر الخوف، ولا يجوز تأخيرها بحال بل تُصلى في وقتها قدر الاستطاعة.

\section{نسخ الحبس في الزنا وتفصيل إقامة الحدود}
كان عقاب الزاني في أول الإسلام الحبس والأذى، كما قال تعالى: ﴿وَاللَّاتِي يَأْتِينَ الْفَاحِشَةَ مِن نِّسَائِكُمْ فَاسْتَشْهِدُوا عَلَيْهِنَّ أَرْبَعَةً مِّنكُمْ ۖ فَإِن شَهِدُوا فَأَمْسِكُوهُنَّ فِي الْبُيُوتِ حَتَّىٰ يَتَوَفَّاهُنَّ الْمَوْتُ أَوْ يَجْعَلَ اللَّهُ لَهُنَّ سَبِيلًا﴾. 
ثم نسخ الله هذا الحبس بآية سورة النور: ﴿الزَّانِيَةُ وَالزَّانِي فَاجْلِدُوا كُلَّ وَاحِدٍ مِّنْهُمَا مِائَةَ جَلْدَةٍ﴾. 
وبينت السنة المطهرة في حديث عبادة بن الصامت هذا "السبيل"، فقال صلى الله عليه وسلم: «خذوا عني، قد جعل الله لهن سبيلاً؛ البكر بالبكر جلد مائة وتغريب عام، والثيب بالثيب جلد مائة والرجم». ثم نُسخ الجلد عن الثيب واكتفي بالرجم فقط كما فعل النبي صلى الله عليه وسلم في رجم ماعز والغامدية.

\subsection{تنصيف الحد للإماء}
قال تعالى في الإماء: ﴿فَإِذَا أُحْصِنَّ فَإِنْ أَتَيْنَ بِفَاحِشَةٍ فَعَلَيْهِنَّ نِصْفُ مَا عَلَى الْمُحْصَنَاتِ مِنَ الْعَذَابِ﴾. 
وقد بين الشافعي أن التنصيف لا يكون إلا فيما يقبل العدد وهو (الجلد)، فتجلد الأمة خمسين جلدة. أما الحدود المتلفة للنفوس كـ (الرجم) فلا يمكن تنصيفها؛ لأن المرجوم قد يموت بحجر واحد أو بألف، وما لا يقبل العدد لا يقبل التنصيف، ولذلك أجمع العلماء على أن الأمة إذا زنت لا تُرجم.

\section{دقة الشريعة البالغة في الدماء والأعراض (قصة العسيف)}
لإثبات عظمة الشريعة واحتياطها في الدماء والأعراض، أورد الشافعي قصة (العسيف)، وهو الأجير الذي زنى بامرأة مستأجره، فافتداه أبوه بمائة شاة وجارية ظناً منه أن ابنه سيُرجم، فلما ارتفعوا للنبي صلى الله عليه وسلم، أبطل هذه الفدية لعدم شرعيتها، وأمر بجلد الشاب مائة وتغريبه عاماً (لأنه بكر غير محصن). 

وأما المرأة المشتكى عليها، فقد قال النبي صلى الله عليه وسلم لأُنيس الأسلمي: «واغدُ يا أُنيس إلى امرأة هذا، فإن اعترفت فارجمها». وفي هذا رد مفحم على أعداء الشريعة؛ فرغم اعتراف الشريك (الشاب)، وإقرار الزوج، وافتداء أب الشاب بالمال، لم يقم النبي صلى الله عليه وسلم الحد على المرأة فوراً! بل اشترط إما أن تعترف هي بنفسها، أو توجد البينة القطعية بـ (أربعة شهداء)، فاعترفت فرجمها، ولو أنكرت لدُرئ عنها الحد.

\subsection{خضوع الذميين لقانون العقوبات الإسلامي}
روى الشافعي عن ابن عمر رضي الله عنهما أن النبي صلى الله عليه وسلم رجم يهوديين زنيا. وهذا دليل ساطع على أن الشريعة الإسلامية هي الحاكمة والنافذة في ديار الإسلام على جميع المقيمين فيها والمواطنين، حتى من غير المسلمين (كأهل الذمة)، إذا ترافعوا إلينا أو أعلنوا منكرهم، ولا حصانة لقانون كفري يخالف شرع الله داخل دولة الإسلام.

\section{صلاة المأموم خلف الإمام القاعد (بين النسخ والإحكام)}
من أسباب الاختلاف بين العلماء في تصحيح الأحاديث دعوى النسخ في صلاة المأموم خلف الإمام إذا صلى قاعداً لعذر.
فقد ثبت في حديث أنس أن النبي صلى الله عليه وسلم سقط من فرس فجُحش شقه (أي جُرح)، فصلى قاعداً وصلى الصحابة وراءه قعوداً، وقال: «وإذا صلى جالساً فصلوا جلوساً أجمعون».
ثم ثبت في حديث عائشة في مرض موته صلى الله عليه وسلم، أنه خرج فصلى قاعداً إلى جنب أبي بكر، وكان الناس يصلون قياماً خلف أبي بكر.

وقد ذهب الإمام الشافعي إلى أن فعل النبي صلى الله عليه وسلم في مرض موته ناسخ لأمره الأول بجلوس المأمومين، وأن من أطاق القيام لا يجوز له أن يجلس خلف الإمام المريض.
بينما ذهب فقهاء آخرون (كالإمام أحمد، ورجحه المحقق أحمد شاكر) إلى عدم النسخ، لأن الأمر بالقعود جاء بأعلى صيغ التأكيد والنهي عن التشبه بالأعاجم، وأن صلاة أبي بكر بالناس قياماً في مرض الموت تحتمل أنه كان هو الإمام المُسمِع للناس، وهي حالة خاصة. والخلاف في ذلك معتبر ومبني على دقة توجيه النصوص.
\chapter{أوجه اختلاف الأحاديث (الجمع والترجيح)}

\section{اختلاف التنوع (صيغ التشهد والقراءات)}
يتابع الإمام الشافعي رحمه الله مناقشة الأحاديث التي ظاهرها الاختلاف، ويبين كيف يتعامل معها الفقيه. ومن ذلك اختلاف الصحابة في ألفاظ "التشهد". فقد روى ابن مسعود، وابن عباس، وعمر بن الخطاب، وعائشة، وجابر، وابن عمر رضي الله عنهم صيغاً للتشهد تختلف في بعض الألفاظ بالزيادة والنقصان والتقديم والتأخير.

والحقيقة أن هذا من "اختلاف التنوع" وليس اختلاف تضاد. فالنبي صلى الله عليه وسلم نوّع في العبادة تيسيراً وتوسعة على الأمة، فكل هذه الصيغ صحيحة ومجزئة، وكلها تعظيم لله عز وجل. وقد اختار الشافعي تشهد ابن عباس رضي الله عنهما (التحيات المباركات الصلوات الطيبات لله...) لأنه أوسع وأكثر لفظاً، مع إقراره بصحة الباقي وعدم تعنيفه لمن أخذ بغيره مما ثبت عن النبي صلى الله عليه وسلم. 

\subsection{القرآن أُنزل على سبعة أحرف}
ومما يشبه اختلاف التنوع، اختلاف القراءات القرآنية. فقد استدل الشافعي بحديث عمر بن الخطاب حين سمع هشام بن حكيم يقرأ سورة الفرقان على غير ما أقرأه النبي صلى الله عليه وسلم، فلما احتكما إليه قال: «إن هذا القرآن أُنزل على سبعة أحرف فاقرؤوا ما تيسر منه». فهذا تنويع وتوسعة ورأفة من الله بالأمة لاختلاف لهجاتهم، ما لم يكن في الاختلاف إحالة وتغيير للمعنى. وإذا جاز هذا الاختلاف في لفظ القرآن للتوسعة، فجوازه في السنن والأذكار (كالتشهد) من باب أولى.

\section{الجمع والترجيح في أحاديث الربا}
ذكر الشافعي اختلاف الرواية على وجه آخر يحتاج إلى جمع أو ترجيح. فالأحاديث المتواترة (كحديث أبي سعيد الخدري، وأبي هريرة، وابن عمر، وعبادة بن الصامت) تحرم "ربا الفضل"، وتنص على أنه لا يجوز بيع الذهب بالذهب، والفضة بالفضة إلا مثلاً بمثل ويداً بيد دون زيادة.

ولكن جاء في المقابل حديث أسامة بن زيد رضي الله عنهما أن النبي صلى الله عليه وسلم قال: «إنما الربا في النسيئة» (أي ربا الأجل). وظاهره حصر الربا في النسيئة فقط وإباحة ربا الفضل.

\textbf{طرق توجيه حديث أسامة:}
\begin{enumerate}
    \item \textbf{مسلك الجمع:} يُحمل حديث أسامة على اختلاف الأصناف (كبيع الذهب بالفضة، أو التمر بالقمح)، فهنا يجوز التفاضل، ويكون الربا المحرم هو النسيئة (التأخير) فقط. أو أن أسامة رضي الله عنه أدرك الجواب ولم يسمع السؤال الذي قيد هذه الإجابة بمسألة معينة.
    \item \textbf{مسلك الترجيح:} إن تعذر الجمع، صرنا إلى الترجيح؛ فالذين رووا تحريم ربا الفضل جماعة كبار أحفظ وأقدم صحبة (كعثمان وعبادة وأبي هريرة وابن عمر)، وحديث الجماعة أولى بالقبول من حديث الواحد.
\end{enumerate}

\section{اختلاف الأحاديث في وقت صلاة الفجر}
ومن الاختلاف الظاهري ما ورد في وقت صلاة الفجر؛ ففي حديث عائشة وزيد بن ثابت وسهل بن سعد: أن النبي صلى الله عليه وسلم كان يصلي الصبح بغلس (أي في الظلمة). وفي المقابل حديث رافع بن خديج: «أسفروا بالفجر فإنه أعظم للأجر». 

\textbf{توجيه الشافعي للتعارض:}
بيّن الشافعي أننا إذا احتجنا للترجيح، نرجح حديث التغليس والمبادرة للصلاة لأسباب:
\begin{itemize}
    \item أنه أشبه وأقرب لظاهر القرآن في المبادرة للخيرات ﴿حَافِظُوا عَلَى الصَّلَوَاتِ وَالصَّلَاةِ الْوُسْطَىٰ﴾، فمن قدمها في أول وقتها كان أولى بالمحافظة.
    \item أنه مروي عن عدد أكبر وأحفظ من الصحابة.
    \item أن الخلفاء الراشدين (أبا بكر وعمر وعثمان وعلي) كانوا يدخلون فيها بغلس.
\end{itemize}
أما حديث الإسفار، فيُحمل على التأكد من طلوع الفجر الصادق (المعترض في الأفق) وعدم التعجل قبل دخول الوقت، أو أن الإسفار يكون في الانصراف منها بسبب إطالة القراءة فيها.

\section{اختلاف الحال في النهي عن استقبال القبلة للغائط والبول}
ومن الاختلاف الذي يرجع لتغير الحال: حديث أبي أيوب الأنصاري في النهي المطلق: «لا تستقبلوا القبلة ولا تستدبرها لغائط أو بول». يقابله حديث ابن عمر: «رأيت رسول الله صلى الله عليه وسلم على لبنتين مستقبلاً بيت المقدس لحاجته» (أي مستدبراً الكعبة).

\textbf{طريقة الجمع بينهما:} 
حمل الشافعي حديث النهي على قضاء الحاجة في "الصحراء" والفضاء المفتوح؛ لسعتها ولئلا تُقذر القبلة أو تُكشف العورات تجاه المصلين. وحمل رخصة الاستدبار (أو الاستقبال) على "البنيان والمنازل"؛ لضيقها ووجود الستر فيها. فابن عمر حدث بما رأى في البنيان، وأبو أيوب حدث بما سمع من النهي العام، والفقيه يجمع بينهما بالتفريق بين حال الصحراء وحال البنيان.

\section{خاتمة في أدب الاختلاف}
إن الخلاف إما أن يكون اختلاف أفهام وتنوع، كصيغ التشهد وقراءة القرآن، فهذا يُعذر فيه المخالف ولا يُعنف. وإما أن يكون اختلاف تضاد ومصادمة للنصوص، كمن يحلل ما حرم الله، أو من يعارض الدليل الواضح بالهوى والآراء المحدثة؛ فهذا لا يُقبل، والكثرة لا تبرر الباطل ﴿وَإِن تُطِعْ أَكْثَرَ مَن فِي الْأَرْضِ يُضِلُّوكَ عَن سَبِيلِ اللَّهِ﴾. 
وقد ضرب الشافعي أروع الأمثلة في أدب الخلاف، حين رد على شيخه سفيان بن عيينة في مسألة نسخ أحاديث، فرد عليه بالأدلة والتاريخ دون تنقيص، ليعلمنا أن العلم مبني على الدليل الشرعي وليس على التقليد الأعمى.
\chapter{أوجه الاختلاف بين الأحاديث: الجمع والتخصيص}

\section{حكم قتل النساء والصبيان في البيات (الغارات الليلية)}
يواصل الإمام الشافعي رحمه الله مناقشة أوجه الاختلاف الظاهري بين الأحاديث، ومن ذلك:
حديث الصعب بن جثامة رضي الله عنه: أن النبي ﷺ سُئل عن أهل الدار من المشركين يُبَيَّتون (أي يُهجم عليهم ليلاً) فيُصاب من نسائهم وذراريهم، فقال: «هم منهم» وفي رواية «هم من آبائهم». وظاهره إباحة إصابتهم.
ويقابله حديث ابن كعب بن مالك: أن النبي ﷺ لما بعث إلى ابن أبي الحُقَيق نهى عن قتل النساء والولدان. وظاهره التحريم.

\textbf{طريقة الجمع أو الترجيح:}
ذهب سفيان بن عيينة إلى أن حديث النهي ناسخ لحديث الإباحة (هم منهم). لكن الشافعي لم يقنع بهذا، وناقشه نقاشاً علمياً وتاريخياً؛ فبين أن حديث الصعب كان في عمرة النبي ﷺ المتأخرة، ويقيناً كان بعد مقتل ابن أبي الحقيق، فلا يمكن أن يكون المتقدم ناسخاً للمتأخر. 
ثم جمع الشافعي بينهما بفقه دقيق: أن النهي عن قتل النساء والولدان يُحمل على (التعمد والقصد) وهم متميزون بعيدون عن المقاتلة؛ لأنهم لم يبلغوا كفراً يُقاتلون عليه، ولا معنى لقتالهم. أما حديث «هم منهم» فيُحمل على البيات (الغارة الليلية) أو اختلاطهم بالمقاتلين بحيث لا يمكن تمييزهم، فإذا أُصيبوا خطأً فلا دية فيهم ولا كفارة ولا قود، لأنهم كفار في دار كفر، وليس لهم حكم الإيمان الذي يعصم الدم.

\subsection{دلالة من القرآن على سقوط الدية والكفارة في دار الكفر}
استدل الشافعي على سقوط الكفارة والدية عنهم بآيات قتل الخطأ في سورة النساء؛ فالمؤمن في دار الإيمان (أو المعاهد) فيه الكفارة والدية ﴿وَدِيَةٌ مُّسَلَّمَةٌ إِلَىٰ أَهْلِهِ وَتَحْرِيرُ رَقَبَةٍ﴾. أما المؤمن الذي في دار الكفر ففيه الكفارة دون الدية ﴿فَإِن كَانَ مِن قَوْمٍ عَدُوٍّ لَّكُمْ وَهُوَ مُؤْمِنٌ فَتَحْرِيرُ رَقَبَةٍ﴾. فإذا كان المؤمن تسقط ديته في دار الكفر، فمن باب أولى أن يسقط القود والدية والكفارة عن نساء المشركين وصبيانهم إذا أُصيبوا خطأً في البيات.

\section{غسل الجمعة بين الوجوب والاستحباب}
ومن الأحاديث التي اختلف في فقهها:
حديث أبي سعيد الخدري: «غسل الجمعة واجب على كل محتلم». وحديث ابن عمر: «من جاء منكم الجمعة فليغتسل». ظاهرها الوجوب الذي لا تصح الصلاة إلا به.

\textbf{توجيه الشافعي وتخصيصه:}
استدل الشافعي بحديث دخول عثمان بن عفان رضي الله عنه المسجد يوم الجمعة وعمر يخطب، فعاتبه عمر لعدم اغتساله واكتفائه بالوضوء، فلو كان الغسل واجباً شرطياً (كركن أو شرط للصلاة) لأمره عمر بالخروج ليغتسل، ولما واصل عثمان صلاته بدونه. فدل إقرار عمر وعثمان والصحابة على أن الغسل هنا (واجب في الاختيار والنظافة) وليس شرطاً في صحة الصلاة.
ويؤيده حديث: «من توضأ يوم الجمعة فبها ونعمت، ومن اغتسل فالغسل أفضل». فالأوامر تُحمل على الاستحباب وطلب الأفضلية، لا على الركنية.

\section{الخطبة على خطبة الأخ والبيع على بيعه}
قال ﷺ: «لا يخطب أحدكم على خطبة أخيه». ظاهره التحريم المطلق بمجرد الخطبة.
لكن الشافعي بين أن هذا النهي مصروف لمعنى خاص؛ واستدل بحديث فاطمة بنت قيس رضي الله عنها، حين طلقها زوجها، فجاءت تستشير النبي ﷺ وأخبرته أن معاوية وأبا جهم خطباها، فقال ﷺ: «أما أبو جهم فلا يضع عصاه عن عاتقه، وأما معاوية فصعلوك لا مال له، انكحي أسامة بن زيد».
فالنبي ﷺ خطبها لأسامة بعد أن علم بخطبة معاوية وأبي جهم، فدل ذلك على أن النهي (لا يخطب أحدكم على خطبة أخيه) إنما يكون بعد (الركون والقبول) من المرأة ووليها. أما قبل الموافقة والقبول، فالحال واحدة، ويجوز للرجل أن يتقدم لخطبتها كما فعل النبي ﷺ لأسامة.

\subsection{لا يبيع الرجل على بيع أخيه}
وهو مقاس على ما قبله؛ لقوله ﷺ: «المتبايعان بالخيار ما لم يتفرقا»، فنهيه عن البيع على بيع أخيه إنما يكون قبل التفرق (في وقت الخيار)، فإذا جاء مشترٍ آخر ليفسد البيعة بزيادة السعر، فهذا هو المحرم. أما إذا تم البيع ولزم العقد وتفرقا، فلا يضره أن يبيع أو يشتري لأن البيع قد لزم ولا يملك فسخه.

\section{الصلاة في أوقات النهي (الفرائض وذوات الأسباب)}
نهى النبي ﷺ عن الصلاة بعد العصر حتى تغرب الشمس، وبعد الصبح حتى تطلع الشمس. 
واحتمل النهي أن يكون عاماً لكل الصلوات، واحتمل أن يكون مخصصاً للنوافل. فبحث الشافعي عن الدلائل، فوجد أن النبي ﷺ قال: «من أدرك ركعة من الصبح قبل أن تطلع الشمس فقد أدرك الصبح»، وهذا يقتضي أنه سيُكمل صلاته أثناء طلوع الشمس وهو وقت نهي. وقال: «من نسي صلاة فليصلها إذا ذكرها»، وجعل وقت الذكر هو وقتها ولو كان في وقت نهي.
وقال في الطواف: «يا بني عبد مناف لا تمنعوا أحداً طاف بهذا البيت وصلى أية ساعة شاء من ليل أو نهار».

فدل هذا كله على أن النهي مخصص بـ (النوافل المطلقة) التي لا سبب لها. أما الصلوات المفروضة، والصلوات التي لها سبب (كسنة الطواف، وركعتي الكسوف، والجنائز) فتُصلى في أي وقت، ولا تدخل في النهي. 
وقد طاف ابن عمر وابن عباس والحسن والحسين وصلوا ركعتي الطواف بعد العصر وبعد الصبح، وأما تأخير عمر بن الخطاب لركعتي الطواف حتى طلعت الشمس، فلأنه حفظ النهي المجرد ولم تبلغه الرخصة، ومن علم حجة على من لم يعلم.

\section{وجوب التسليم للسنة إذا ثبتت}
يختم الشافعي هذا الباب بكلمة عظيمة في منهج الاستدلال: «وإذا ثبت عن رسول الله ﷺ الشيء، فهو اللازم لجميع من عرفه، لا يقويه ولا يوهنه شيء غيره، بل الفرض الذي على الناس اتباعه... ولم يجعل الله لأحد معه أمراً يخالف أمره». 
فالسنة حاكمة بنفسها، ولا تحتاج لعمل الناس لتقويتها، ولا يضعفها ترك بعض الناس لها.
\chapter{الجمع بين الأحاديث، ودلالة النهي، وأقسام العلم}

\section{الجمع بين الأحاديث (المزابنة والعرايا)}
الحمد لله والصلاة والسلام على رسول الله، أما بعد؛ فيواصل الإمام الشافعي رحمه الله تعالى بيان الأحاديث التي ظاهرها الاختلاف وكيفية الجمع بينها.
ومن ذلك: ما روي عن ابن عمر رضي الله عنهما أن رسول الله ﷺ «نهى عن المزابنة». والمزابنة: بيع الثمر بالتمر كيلاً (أي بيع الرطب على النخل بالتمر الجاف). وفي حديث سعد بن أبي وقاص رضي الله عنه أنه سُئل النبي ﷺ عن بيع التمر بالرطب، فقال: «أينقص الرطب إذا يبس؟»، قالوا: نعم، فنهى عن ذلك. 
فهذه الأحاديث تدل على تحريم هذا البيع صراحة لوجود التفاضل والجهالة فيه.

وفي المقابل، روي عن زيد بن ثابت رضي الله عنه أن رسول الله ﷺ «رخص لصاحب العرية أن يبيعها بخرصها». والعرايا: أن يمتلك الفقير رطباً على النخل ويحتاج إلى تمر ليأكله، فرخص له النبي ﷺ أن يبيع هذا الرطب بتمر يخرصه (يقدره) أهل الخبرة، ما لم يبلغ خمسة أوسق.
ظاهر هذين الحديثين التعارض؛ فالأول ينهى مطلقاً، والثاني يبيح. والجمع بينهما عند الشافعي: أن النهي عام عن المزابنة لحرمة الجهالة والغرر والربا، وأن العرايا رخصة مستثناة من هذا النهي العام لحاجة الفقراء بحدود وضوابط (وهي الخرص وعدم تجاوز خمسة أوسق). فالعام باقٍ على عمومه، والعرايا مخصوصة ومستثناة للحاجة.

\section{الجمع بين بيع المجهول وإباحة السَّلَم}
ومن ذلك أيضاً: حديث حكيم بن حزام رضي الله عنه، حيث قال: «يا رسول الله يأتيني الرجل فيريد مني المبيع ليس عندي، أفأبتاعه له من السوق؟» فقال ﷺ: «لا تبع ما ليس عندك».
ويقابله حديث ابن عباس رضي الله عنهما: «قدم رسول الله ﷺ المدينة وهم يسلفون في التمر السنتين والثلاث، فقال: من أسلف في شيء فليُسلف في كيل معلوم، ووزن معلوم، إلى أجل معلوم». والسَّلَم (السَّلَف) هو بيع شيء ليس موجوداً عند العقد.

\textbf{وجه الجمع بينهما:}
النهي في حديث حكيم عن بيع «عين» ليس يملكها البائع وليست بحضرته وقد تهلك قبل تسليمها (كأن يبيع سيارة محددة لا يملكها). أما «السَّلَم» المباح، فهو بيع «موصوف في الذمة» مضمون على البائع (كأن يبيع طناً من القمح بصفة كذا يسلمه بعد سنة)، فهذا مباح تيسيراً على الناس. فلا تعارض بين النهي عن بيع المعيَّن الغائب، وبين إباحة بيع الموصوف المضمون في الذمة. 
فالجمع بين الأحاديث هو مسلك العلماء، ولا يصار إلى دعوى التعارض والنسخ إلا عند تعذر الجمع.

\section{صفة نهي الله ورسوله (بين التحريم والكراهة)}
قسم الشافعي رحمه الله ما نهى عنه الشرع إلى قسمين: نهي يقتضي التحريم والفساد، ونهي يقتضي الكراهة والتنزيه (أو الإرشاد).

\subsection{النهي الذي يقتضي التحريم والفساد}
إذا جاء النهي عن شيء في أصل العقود أو العبادات، فإنه يقتضي التحريم وإبطال العقد. كنهي الله ورسوله عن نكاح المرأة على عمتها أو خالتها، ونهيه عن الجمع بين الأختين، ونهيه عن نكاح المتعة، ونهيه عن نكاح الشغار (وهو أن يزوجه أخته على أن يزوجه الآخر أخته دون مهر).
فأي نكاح أو بيع وقع على هذه الوجوه المنهي عنها فهو مفسوخ وباطل.
ويقرر الشافعي قاعدة عظيمة هنا: أن «المعصية لا تُحِلُّ مُحَرَّماً». فالأصل في الأبضاع والأموال التحريم، ولا تستباح إلا بما شرعه الله. والغاية المشروعة لا تُنال بوسيلة محرمة؛ فالتمكين للدين لا ينال بالانحراف عن منهج السلف واللجوء للبدع أو النظم المستوردة، وما عند الله لا ينال بمعصيته.

\subsection{النهي الذي يقتضي الكراهة والتنزيه (الإرشاد)}
قد يأتي النهي عن شيء أصله مباح، ويكون النهي لعلة أخرى خارجة عن ذات الشيء، فهذا نهي أدب وإرشاد. ومن أمثلته:
\begin{itemize}
    \item \textbf{النهي عن اشتمال الصمَّاء والاحتباء:} الأصل في الثياب الإباحة، لكن نهى النبي ﷺ أن يلتف الرجل بثوب واحد لا يستطيع إخراج يده منه (الصمَّاء)، ونهى أن يحتبي (يضم ركبتيه إلى صدره) في ثوب واحد ليس تحته غيره. والنهي هنا ليس لكون الثوب محرماً، بل خشية انكشاف العورة، فهو نهي إرشاد.
    \item \textbf{النهي عن الأكل من وسط الصحفة:} الأصل إباحة الطعام، لكن نهى النبي ﷺ عن الأكل من أعلى الصحفة (وسطها) وأمر بالأكل من الحواف؛ لأن البركة تنزل في وسطها. فهذا نهي تأديب وليس تحريماً لذات الطعام.
    \item \textbf{النهي عن التعريس بالطريق:} التعريس هو النزول للراحة والنوم في السفر. وقد نهى النبي ﷺ عن التعريس في قارعة الطريق؛ لأنها ممر الهوام (كالحيات) ومسلك الدواب. فالنهي هنا للإرشاد والمحافظة على المسلم، وليس تحريماً لجلوسه على الأرض.
\end{itemize}
ومن فعل شيئاً من هذا عالماً بالنهي فهو عاصٍ (لمخالفته الأمر)، ولكنه ليس كمن فعل المحرمات القطعية المفسدة للعقود، لأن هذه تتعلق بأمور مباحة الأصل.

\section{أقسام العلم الشرعي}
عقد الشافعي باباً نفيساً في أقسام العلم، فبين أن العلم ينقسم إلى قسمين:

\subsection{علم العامة (فرض العين)}
وهو العلم الذي لا يسع مكلفاً بالغاً عاقلاً جهله. وهو نص كتاب الله وما تواتر من سنة رسول الله ﷺ في أصول الدين، مثل: فرضية الصلوات الخمس، وصيام رمضان، وحج البيت، وزكاة الأموال، وتحريم الزنا والقتل والسرقة والخمر. 
هذا العلم يعلمه العامة عن العامة، ولا يُقبل فيه التأويل، ولا يُعذر فيه أحد بالجهل، ولا يختلف فيه مسلمان.

\subsection{علم الخاصة (فرض الكفاية)}
وهو العلم بتفاصيل الأحكام، وفروع المسائل، ودقائق السنن والناسخ والمنسوخ، وما يستنبط بالقياس والاجتهاد.
هذا النوع من العلم ليس فرضاً على جميع الأمة، بل هو (فرض كفاية)؛ إذا قام به مَن فيهم غَناء وكفاية من طلبة العلم والعلماء، سقط الإثم عن الباقين، ونال القائمون به الفضل والدرجات العُلى.

واستدل الشافعي على فرض الكفاية بآيات الجهاد؛ فإن الله أوجب الجهاد، ولكنه رفع الحرج عن القاعدين (من غير أولي الضرر) ووعدهم الحسنى، وفضل المجاهدين عليهم درجات. فلو كان فرض عين لأثم القاعدون. وكذلك قال تعالى: ﴿وَمَا كَانَ الْمُؤْمِنُونَ لِيَنفِرُوا كَافَّةً ۚ فَلَوْلَا نَفَرَ مِن كُلِّ فِرْقَةٍ مِّنْهُمْ طَائِفَةٌ لِّيَتَفَقَّهُوا فِي الدِّينِ﴾، فجعل النفير لطلب العلم كفاية في طائفة تنذر قومها إذا رجعوا إليهم.
ومن أمثلة فرض الكفاية أيضاً: صلاة الجنازة وتغسيل الميت ودفنه، ورد السلام؛ إذا قام به واحد من القوم أجزأ عن الباقين.
\chapter{حجية خبر الآحاد وشروط الحديث الصحيح}

\section{مقدمة في خبر الخاصة (خبر الآحاد)}
الحمد لله والصلاة والسلام على رسول الله، أما بعد؛ فقد شرع الإمام الشافعي رحمه الله تعالى في بيان باب عظيم من أبواب أصول الفقه والحديث، وهو "باب خبر الواحد".
وقد سأله سائل: حدد لي أقل ما تقوم به الحجة على أهل العلم حتى يثبت عليهم خبر الخاصة؟ 
و(خبر الخاصة) عند الشافعي يقابله (خبر العامة)؛ فخبر العامة هو المتواتر الذي يعلمه عامة المسلمين (كفرضية الصلوات الخمس)، أما خبر الخاصة فهو ما ينقله ويعلمه أهل العلم وطلابه، وهو ما نطلق عليه اصطلاحاً (خبر الآحاد).
فأجاب الشافعي: أقل ما تقوم به الحجة هو خبر الواحد عن الواحد حتى يُنتهى به إلى النبي ﷺ، أو من يُنتهى به إليه دونه من الصحابة والتابعين.

\section{شروط الحديث الصحيح عند الإمام الشافعي}
قرر الشافعي رحمه الله أنه لا تقوم الحجة بخبر الخاصة (الآحاد) إلا باجتماع أمور وشروط في رواته، وهي التي استقر عليها علماء مصطلح الحديث في تعريف "الحديث الصحيح"، وهي:
\begin{enumerate}
    \item \textbf{العدالة:} أن يكون الراوي ثقة في دينه، معروفاً بالصدق في حديثه (العدل).
    \item \textbf{الضبط والفهم:} أن يكون عاقلاً فاهماً لما يحدّث به، عالماً بما يحيل معاني الحديث من الألفاظ (إذا كان يروي بالمعنى).
    \item \textbf{الأداء كما سمع:} أن يكون ممن يؤدي الحديث بحروفه كما سمع إذا لم يكن عالماً بلغة العرب وما يحيل المعاني؛ لئلا يحوّل الحلال إلى حرام بخطأ في اللفظ.
    \item \textbf{ضبط الصدر أو الكتاب:} أن يكون حافظاً إن حدّث من حفظه، وضابطاً لكتابه حافظاً له إن حدّث من كتابه.
    \item \textbf{عدم الشذوذ (موافقة الثقات):} إذا شَرِك أهل الحفظ في الحديث وافق حديثهم، ولا يخالف الثقات.
    \item \textbf{السلامة من التدليس:} أن يكون بريئاً من أن يكون مدلّساً؛ يحدّث عمن لقي ما لم يسمع منه.
    \item \textbf{اتصال السند:} أن يكون السند موصولاً إلى النبي ﷺ أو إلى من دونه، فكل راوٍ في السند يُشترط فيه هذه الشروط ليكون مثبتاً لمن حدّثه وعمن حدّث عنه.
\end{enumerate}

\section{الفرق بين الشهادة والرواية}
طلب السائل من الشافعي أن يضرب له مثلاً ويقيس خبر الآحاد على شيء من الشهادات التي يعرفها الناس. فبين الشافعي أن الرواية والشهادة يجتمعان في أشياء (كاشتراط العدالة والصدق)، لكنهما يفترقان في أمور جوهرية:
\begin{itemize}
    \item \textbf{في العدد والذكورة:} تُقبل في الرواية رواية الواحد ورواية المرأة الواحدة، بينما لا تُقبل في الشهادة شهادة الواحد وحده ولا المرأة وحدها.
    \item \textbf{في صيغة الأداء:} يُقبل في الرواية صيغة (عنعنة) الثقة غير المدلس كقوله: "حدثني فلان عن فلان"، بينما لا تُقبل في الشهادة إلا بالتصريح القطعي: "سمعتُ، أو رأيتُ، أو أشهدني".
    \item \textbf{في الترجيح:} تختلف الأحاديث فيُرجَّح ببعضها استدلالاً بكتاب أو سنة أو إجماع أو قياس، وهذا لا يوجد في باب الشهادات.
\end{itemize}

\subsection{لماذا يُحتاط للرواية أكثر من الشهادة؟}
بيّن الشافعي أن إحالة معنى الحديث أخفى من إحالة معنى الشهادة. فالشاهد يُنشئ كلاماً يعبر به عما رأى بلفظه، والخطأ في الشهادة يضر فرداً بعينه (في مال أو عقوبة). أما المحدّث فهو ناقل لحكم شرعي يحلل ويحرم للأمة جمعاء، والخطأ في لفظة قد يحيل الحلال حراماً؛ لذا كان الاحتياط في الرواية والتشديد في شروط رواتها (كالضبط والفهم والسلامة من التدليس) أشد من الاحتياط في الشهود.

\section{حكم رواية المدلّس وتعديل الرواة}
قرر الشافعي مذهبه الصارم في التدليس قائلاً: «ومن عَرَفناه دلَّس مَرَّة، فقد أبان لنا عورته في روايته». فمن دلّس (بأن روى عمن لقيه ما لم يسمع منه بصيغة موهمة كـ "عن")، لا يُقبل حديثه حتى يصرّح بالتحديث (سمعت أو حدثني). والتدليس ليس كذباً محضاً، إذ لو كان كذباً لردت روايته بالكلية، لكنه أظهر عورة في دقة نقله، فوجب التوقف في عنعنته.
كما بيّن أن رواية الثقة عن شيخ لا نعرفه لا تُعد توثيقاً له، بل لا بد من تعديل صريح للراوي.

\section{التحذير من الكذب على رسول الله ﷺ}
ختم الشافعي هذا المبحث ببيان ورع المحدثين وحرصهم على ألا يكذبوا على رسول الله ﷺ؛ مستدلاً بالأحاديث المتواترة في ذلك كقوله ﷺ: «من يقل عليّ ما لم أقل فليتبوأ مقعده من النار»، وقوله: «حدثوا عن بني إسرائيل ولا حرج، وحدثوا عني ولا تكذبوا عليّ».

فبيّن أن النبي ﷺ لم يُبح الكذب في أخبار بني إسرائيل، بل أباح التحديث عنهم على جهة البلاغ لتعذر الأسانيد إليهم ولطول المدة، ولكنه شدّد في النقل عنه ﷺ بألا يكون إلا بنقل الإسناد والتثبت الدقيق، والاحتراز التام من الكذب عليه، فكلامه دين وشرع إلى يوم القيامة، ولا يستدل على الصدق والكذب إلا بمعرفة الرواة ونقل الإسناد.
\chapter{حجية خبر الآحاد والأدلة على قبوله}

\section{نضارة وجه مبلغ السنة ودلالتها على الآحاد}
الحمد لله والصلاة والسلام على رسول الله، أما بعد؛ فقد بوب الإمام الشافعي رحمه الله تعالى باباً عظيماً لبيان الحجة في تثبيت خبر الواحد، ليرد به على المعتزلة وأهل البدع الذين يشككون في خبر الواحد ويردونه في العقائد أو الأحكام.
استدل الشافعي أولاً بحديث عبد الله بن مسعود رضي الله عنه أن النبي ﷺ قال: «نَضَّرَ الله عبداً -وفي رواية: امرأً- سَمِعَ مَقَالَتي فَحَفِظَهَا وَوَعَاهَا وَأَدَّاهَا، فَرُبَّ حَامِلِ فِقْهٍ غَيْرُ فَقِيهٍ، وَرُبَّ حَامِلِ فِقْهٍ إِلَى مَنْ هُوَ أَفْقَهُ مِنْهُ. ثَلَاثٌ لا يُغِلُّ عَلَيْهِنَّ قَلْبُ مُسْلِمٍ: إِخْلَاصُ العَمَلِ للهِ، وَالنَّصِيحَةُ لِلْمُسْلِمِينَ، وَلُزُومُ جَمَاعَتِهِمْ، فَإِنَّ دَعْوَتَهُمْ تُحِيطُ مِنْ وَرَائِهِمْ».
والنضارة هي حسن الوجه وبريقه. وقد استنبط الشافعي من هذا الحديث حجية خبر الواحد؛ لأن النبي ﷺ ندب إلى استماع مقالته وحفظها وأدائها، وخص بالذكر «امرأً» أو «عبداً» وهو (الواحد). فدل ذلك على أنه لا يأمر الواحد بأن يؤدي عنه إلا وخبر هذا الواحد تقوم به الحجة على من أُدي إليه؛ لأنه إنما يؤدي حلالاً وحراماً وأمراً ونهياً. 
كما دل الحديث على التفرقة بين الحفظ والفقه، فليس كل راوٍ وحافظ فقيهاً مستنبطاً، والخير كل الخير فيمن جمع بين الحفظ والفقه والأصول.

\section{لزوم جماعة المسلمين وحجية الإجماع}
في قوله ﷺ: «ولزم جماعتهم» حجة ظاهرة على اعتبار إجماع المسلمين حجة لازمة وملزمة، لا يجوز لمؤمن يؤمن بالله واليوم الآخر أن يخالفه. وما أجمع عليه الصحابة والتابعون سلفاً لا يجوز نقضه بادعاء إجماعات معاصرة مبنية على شريعة مستوردة كالديمقراطية أو غيرها، فمن ادعى الإجماع على باطل يخالف إجماع السلف الصالح المتقدم فهو كمن ادعى إجماع الرافضة على تكفير الصحابة؛ كلاهما إجماع باطل لا وزن له في دين الله.

\section{وجوب التسليم للسنة (حديث الأريكة)}
ساق الشافعي حديث أبي رافع رضي الله عنه أن النبي ﷺ قال: «لَا أُلْفِيَنَّ أَحَدَكُمْ مُتَّكِئًا عَلَى أَرِيكَتِهِ يَأْتِيهِ الأَمْرُ مِنْ أَمْرِي مِمَّا أَمَرْتُ بِهِ أَوْ نَهَيْتُ عَنْهُ، فَيَقُولُ: لَا نَدْرِي! مَا وَجَدْنَا فِي كِتَابِ اللهِ اتَّبَعْنَاهُ». 
وهذا تثبيت من النبي ﷺ لسنته وإعلام للأمة أنها لازمة لهم وإن لم يجدوا لها نص حكم في كتاب الله. فالأريكة هي السرير أو ما يتكئ عليه الرجل، وفيه وعيد شديد لمن يرد سنة رسول الله ﷺ بحجة الاكتفاء بالقرآن، فمن رد السنة فقد رد أمر الله جل وعلا.

\section{أمثلة وتطبيقات من العهد النبوي لقبول خبر الواحد}
ذكر الشافعي أمثلة عملية متواترة المعنى تدل على أن الصحابة كانوا يقبلون خبر الواحد ويقيمون به الحجة وينقضون به اليقين السابق:

\subsection{خبر أم سلمة في القبلة للصائم}
رجل من الأنصار قَبَّل امرأته وهو صائم، فوجد في نفسه من ذلك، فأرسل امرأته تسأل أم سلمة رضي الله عنها، فأخبرتها أن النبي ﷺ يقبل وهو صائم. فرجعت إلى زوجها فلم يقبل ظناً منه أنها خصوصية للنبي ﷺ. فلما استوثق وعلم، سَلَّم. والشاهد أن النبي ﷺ قال لأم سلمة: «أَلَا أَخْبَرْتِيهَا أَنِّي أَفْعَلُ ذَلِكَ؟» فلو لم تكن الحجة قائمة بخبر امرأة واحدة عن النبي ﷺ لم يأمرها بذلك.

\subsection{تحول أهل قباء عن قبلتهم بخبر واحد}
كان الصحابة يصلون الصبح في قباء متوجهين إلى بيت المقدس (وهي القبلة اليقينية التي فرضوها وسمعوها من النبي ﷺ)، فجاءهم آتٍ واحد، فأخبرهم أن القبلة قد حولت إلى الكعبة، فاستداروا وهم في صلاتهم. فلو لم يكن خبر الواحد الثقة حجة قاطعة، ما ترك الصحابة فرضاً متيقناً لخبر رجل واحد، وما استداروا دون أن يرسلوا ليتأكدوا من النبي ﷺ. فلما أقرهم النبي ﷺ على ذلك، ثبت أن خبر الواحد حجة.

\subsection{إراقة أبي طلحة للخمر بخبر المخبر}
كان أنس بن مالك رضي الله عنه يسقي أبا طلحة وأبا عبيدة بن الجراح وأبي بن كعب شراباً من فضيخ وتمر (وكان حلالاً حينها). فجاءهم رجل واحد فقال: إن الخمر قد حُرِّمت. فقال أبو طلحة: يا أنس، قم إلى هذه الجرار فاكسرها. 
ومعلوم أن إضاعة المال محرمة، وهؤلاء من كبار الصحابة وعلماء الأمة، فلو لم تقم الحجة عندهم بخبر هذا الواحد ما أراقوا هذا المال الذي كان حلالاً إلى وقت قريب، وما رضوا بإتلافه بمجرد خبر رجل واحد.

\subsection{إرسال الآحاد لتبليغ الأحكام والحدود والعقائد}
\begin{itemize}
    \item \textbf{في الحدود:} أرسل النبي ﷺ أنيساً الأسلمي وحده في قضية العسيف قائلاً: «وَاغْدُ يَا أُنَيْسُ إِلَى امْرَأَةِ هَذَا، فَإِنْ اعْتَرَفَتْ فَارْجُمْهَا». فلو لم تقم الحجة بالواحد ما أرسله وحده لإقامة حد الرجم الذي يزهق النفس.
    \item \textbf{في الأحكام ومناسك الحج:} أرسل علي بن أبي طالب رضي الله عنه وحده على جمل يصرخ في منى: «إنها أيام طعام وشراب فلا يصومن أحد». وأرسل ابن مربع الأنصاري وحده ليقول للناس في عرفة: «قِفُوا عَلَى مَشَاعِرِكُمْ فَإِنَّكُمْ عَلَى إِرْثٍ مِنْ إِرْثِ أَبِيكُمْ إِبْرَاهِيمَ». وبعث أبا بكر وعلياً رضي الله عنهما وحدهما لإقامة الحج ونبذ العهود للمشركين.
    \item \textbf{في العقائد والتوحيد:} بعث معاذ بن جبل وحده إلى اليمن ليدعوهم إلى التوحيد (شهادة أن لا إله إلا الله)، وإلى الصلاة والزكاة. فلو كان خبر الواحد لا يُقبل في العقائد (كما تزعم المبتدعة) لما أرسل النبي ﷺ معاذًا وحده لدعوة أهل اليمن للعقيدة والأحكام معاً.
\end{itemize}

\section{الفرق بين القاضي والشاهد في الأخبار}
أوضح الشافعي أن قضاء القاضي هو في حقيقته إخبار وإمضاء لبينة أو إقرار ثبت عنده. فالقاضي يلزمه أن ينفذ حكم الله بعلمه من خلال البينة، ويكون مخبراً بالحلال والحرام الذي لزمه. بينما الشاهد الذي يشهد على واقعة لا ينفذ الحكم بشهادته منفرداً، بل لا بد من نصاب الشهادة (كشاهدين). فدل هذا على أن مقام النقل والرواية (كالإخبار بالحكم الشرعي) يختلف عن مقام الشهادة على الوقائع والأعيان.
\chapter{قبول الصحابة لخبر الآحاد وشروط الحديث المرسل}

\section{رجوع الصحابة وخلفائهم لخبر الآحاد}
الحمد لله والصلاة والسلام على رسول الله، أما بعد؛ فيواصل الإمام الشافعي رحمه الله تعالى إقامة الحجج الدامغة على حجية خبر الواحد (الآحاد) الملزم في العقائد والأحكام، وذلك بسرد وقائع عملية متواترة لكبار الصحابة، رجعوا فيها عن فتاواهم وأقضيتهم واجتهاداتهم بمجرد بلوغهم خبراً لواحد من الصحابة عن رسول الله ﷺ.

\subsection{رجوع عمر بن الخطاب في دية الأصابع والميراث من الدية}
قضى عمر بن الخطاب رضي الله عنه باجتهاده في دية أصابع اليدين تفاضلاً؛ فقضى في الإبهام بخمس عشرة، وفي التي تليها بعشر، وفي الوسطى بعشر، وهكذا (لأن المنافع والجمال تختلف). حتى وُجد كتاب آل عمرو بن حزم الذي كتبه النبي ﷺ لعمرو حين بعثه لليمن، وفيه: «وفي كل إصبع مما هنالك عشر من الإبل»، فرجع عمر رضي الله عنه عن قضائه إلى نص حديث النبي ﷺ. 

وكذلك كان عمر رضي الله عنه يفتي بأن (الدية للعاقلة)، ولا ترث المرأة من دية زوجها المقتول شيئاً، حتى أخبره الضحاك بن سفيان رضي الله عنه أن رسول الله ﷺ كتب إليه: «أن يورّث امرأة أشيم الضبابي من دية زوجها»، فرجع عمر عن رأيه إلى سنة رسول الله ﷺ.

\subsection{رجوعه في دية الجنين وفي أخذ الجزية من المجوس}
ناشد عمر بن الخطاب رضي الله عنه الناس: من سمع من النبي ﷺ في إسقاط الجنين شيئاً؟ (وكاد أن يقضي فيه باجتهاده)، فقام حَمَلُ بن مالك بن النابغة رضي الله عنه فأخبره أنه كان متزوجاً بامرأتين (ضرتين)، فضربت إحداهما الأخرى بعمود خيمة فألقت جنيناً ميتاً، فقضى فيه النبي ﷺ بـ «غُرة» (أي عبد أو أمة). فقال عمر: «لو لم أسمع فيه لقضينا فيه برأينا»، ورجع لحكم النبي ﷺ.
وكذلك توقف عمر في أخذ الجزية من المجوس لأنهم ليسوا أهل كتاب، حتى شهد عبد الرحمن بن عوف رضي الله عنه أنه سمع النبي ﷺ يقول: «سُنّوا بهم سنة أهل الكتاب»، فأخذ منهم الجزية بخبر عبد الرحمن وحده.

\section{توجيه طلب عمر لخبر ثانٍ في بعض الأحاديث}
فإن قيل: قد طلب عمر شاهداً آخر مع أبي موسى الأشعري رضي الله عنه في حديث «الاستئذان ثلاثاً»، فهل هذا لعدم قبوله لخبر الواحد؟ 
يُجاب بأن عمر لم يفعل ذلك رداً لخبر الواحد، وإلا لكان متناقضاً يقبل الآحاد مرات ويرده أخرى! بل فعل ذلك من باب (الاحتياط والتثبت) وسد الباب أمام من قد تسول له نفسه الكذب على رسول الله ﷺ، ولذلك قال لأبي موسى: «أما إني لم أتهمك، ولكن خشيت أن يتقول الناس على رسول الله ﷺ». والاحتياط بطلب التأكيد لا ينفي حجية خبر الواحد القاطع.

\section{أمثلة أخرى لرجوع الصحابة والتابعين للآحاد}
\begin{itemize}
    \item \textbf{رجوع ابن عباس في طواف الوداع للحائض:} كان زيد بن ثابت يفتي بأن الحائض لا تنفر حتى تطوف بالبيت طواف الوداع قياساً على عموم الحجاج، فأحاله ابن عباس إلى امرأة أنصارية سألت النبي ﷺ فأمرها بالنفور وسقوط طواف الوداع عنها، فرجع زيد رضي الله عنه لخبر المرأة الواحدة.
    \item \textbf{رجوع ابن عباس في مسألة الخضر:} كذّب ابن عباس نَوْفاً البَكَالي حين زعم أن موسى صاحب الخضر ليس موسى بني إسرائيل، فلما حدّثه أُبي بن كعب بحديث النبي ﷺ في قصة موسى والخضر، سلم ابن عباس ورجع.
    \item \textbf{رجوع طاووس وابن عمر:} رجع طاووس عن صلاة ركعتين بعد العصر حين نهاه ابن عباس وأخبره بنهي النبي ﷺ عنهما. ورجع ابن عمر عن (المخابرة) وهي كراء الأرض الزراعية لما أخبره رافع بن خديج بنهي النبي ﷺ عنها، وترك ما كان يراه حلالاً ينتفع به لخبر رجل واحد.
    \item \textbf{رجوع القضاة والولاة:} قضى سعد بن إبراهيم بن عبد الرحمن بن عوف في مسألة برأي ربيعة الرأي، فأخبره ابن أبي ذئب بحديث عن النبي ﷺ يخالف قضاءه، فدعا سعد بكتاب القضية فشقه ونقض حكمه، وقال: «أنفذ قضاء رسول الله ﷺ وأرد قضاء سعد».
\end{itemize}
وقد أجمع علماء الأمصار في مكة والمدينة والشام واليمن (كطاووس ووهب بن منبه ومكحول) والكوفة (كابن أبي ليلى والأسود وعلقمة والشعبي) على تثبيت خبر الواحد والعمل به في كل زمان.

\section{شروط العمل بالحديث المرسل عند الشافعي}
الحديث المنقطع، ومنه "المرسل" (وهو ما رفعه التابعي للنبي ﷺ مسقطاً الصحابي)، هل تقوم به الحجة؟ 
أسس الإمام الشافعي قواعد صارمة ومحكمة لقبول المرسل، فبيّن أن مراسيل صغار التابعين لا تُقبل ألبتة لكثرة تساهلهم، أما مراسيل (كبار التابعين) كسعيد بن المسيب، فلا تُقبل أيضاً إلا باجتماع شروط تدل على صحة المخرج، وهي:
\begin{enumerate}
    \item \textbf{أن يسند من وجه آخر:} بأن يُروى نفس الحديث بإسناد متصل (مسند) يوافق معنى المرسل.
    \item \textbf{أن يوافق مرسلاً آخر:} بأن يرسله تابعي آخر كبير أخذ العلم عن شيوخ غير شيوخ التابعي الأول، لانتفاء التواطؤ.
    \item \textbf{موافقة قول صحابي:} أن يوجد من أقوال الصحابة ما يوافق دلالة هذا الحديث المرسل.
    \item \textbf{موافقة عوام أهل العلم:} أن يفتي عوام أهل العلم وعامتهم بمقتضى ما جاء في هذا المرسل.
    \item \textbf{شروط في الراوي المُرسِل نفسه:} أن يكون التابعي المرسِل لا يروي (إذا سمّى) إلا عن الثقات المأمونين، ولا يرغب عن الرواية عن مجهول، وإذا شرك الحُفّاظ المأمونين في حديث لم يخالفهم، بل يوافقهم أو ينقص عنهم ولا يزيد.
\end{enumerate}
ومع كل هذه الشروط، فإن الحديث (المسند المتصل) يظل أقوى من (المرسل)؛ لأن المرسل يحمل انقطاعاً ورواية عن مجهول، والاحتياط للدين يقتضي تقديم المتصل، لكن من توافرت في مرسلاته هذه الشروط من كبار التابعين أُلحقت بالحجة.

\section{خاتمة في خطورة رد السنن}
شنّع الشافعي رحمه الله على من يرد الأحاديث المسندة الصحيحة بالهوى أو العصبية، كمن يتمسك بالأحاديث المرسلة الضعيفة لنصرة مذهبه، ويرد الأحاديث المتصلة الصحيحة التي تلزمه الحجة بها، فهذا من اتباع الهوى والتقليد الأعمى المذموم الذي يودي بصاحبه للضلال.
\chapter{حجية الإجماع والقياس وإبطال الاستحسان}

\section{مقدمة في حجية الإجماع}
الحمد لله والصلاة والسلام على رسول الله، أما بعد؛ فقد عقد الإمام الشافعي رحمه الله تعالى باباً نفيساً في "الإجماع" ليبين حجية ما اجتمع عليه المسلمون مما ليس فيه نص حكم يُحكى عن النبي ﷺ.
فبين أن علماء الأمة إذا أجمعوا على مسألة وذكروا دليلها، فهذا هو الدليل. وأما إذا أجمعوا ولم يذكروا دليلاً نصياً، فإننا نتبعهم في إجماعهم؛ لأن سنة رسول الله ﷺ إن غابت عن بعضهم فلن تغيب عن جميعهم، ولأن عامة المسلمين وعلمائهم لا تجتمع على خلاف لسنة رسول الله ﷺ، ولا تجتمع على خطأ وضلالة أبداً.

\section{الأدلة على حجية الإجماع ولزوم الجماعة}
واستدل الشافعي رحمه الله على حجية الإجماع بحديث ابن مسعود رضي الله عنه أن النبي ﷺ قال: «نَضَّرَ الله عبداً سَمِعَ مَقَالَتي فَحَفِظَهَا وَوَعَاهَا وَأَدَّاهَا... ثَلَاثٌ لا يُغِلُّ عَلَيْهِنَّ قَلْبُ مُسْلِمٍ: إِخْلَاصُ العَمَلِ للهِ، وَالنَّصِيحَةُ لِلْمُسْلِمِينَ، وَلُزُومُ جَمَاعَتِهِمْ، فَإِنَّ دَعْوَتَهُمْ تُحِيطُ مِنْ وَرَائِهِمْ».
وكذلك بحديث عمر بن الخطاب رضي الله عنه بالجابية حين خطب الناس فقال: إن رسول الله ﷺ قام فينا مقامي فيكم، فقال: «...عَلَيْكُمْ بِالْجَمَاعَةِ، وَإِيَّاكُمْ وَالْفُرْقَةَ، فَإِنَّ الشَّيْطَانَ مَعَ الْوَاحِدِ وَهُوَ مِنَ الاثْنَيْنِ أَبْعَدُ، مَنْ أَرَادَ بُحْبُوحَةَ الْجَنَّةِ فَلْيَلْزَمِ الْجَمَاعَةَ».

\subsection{المعنى الصحيح للزوم الجماعة}
سُئل الشافعي: ما معنى أمر النبي ﷺ بلزوم جماعتهم؟
فقال: لا معنى له إلا لزوم إجماعهم؛ لأن أبدان المسلمين متفرقة في البلدان ولا يمكن لزومها، ولأن اجتماع الأبدان قد يضم المسلم والكافر والبر والفاجر. فلم يكن للزوم أبدانهم معنى، وإنما المعنى: لزوم ما هم عليه من التحليل والتحريم والطاعة. فمن قال بما تقول به جماعة المسلمين (في العقيدة والمنهج والأحكام) فقد لزم جماعتهم، ومن خالف ما أجمعوا عليه فقد خالف الجماعة.
وهذا الإجماع هو إجماع علماء السلف، ولا عبرة بدعاوى الإجماع المعاصرة التي تخالف ما كان عليه السلف الصالح، كمن يدعي الإجماع على النظم الوضعية والديمقراطية المستوردة المصادمة لحاكمية الشريعة.

\section{الاجتهاد والقياس: اسمان لمعنى واحد}
ثم انتقل الشافعي لبيان القياس، وذكر أن القياس والاجتهاد هما اسمان لمعنى واحد. وجماعهما: أن كل ما نزل بمسلم ففيه حكم لازم، إما أن يكون نصاً بعينه فيُتبع، وإما أن تطلب فيه الدلالة على سبيل الحق بالاجتهاد، وهذا الاجتهاد هو "القياس".

\subsection{التكليف بالظاهر دون الباطن}
وبيّن الشافعي أن العباد كُلفوا بإصابة الحق في الظاهر والباطن معاً فيما فيه نص، وكُلفوا بإصابة الحق في الظاهر (وهو ما ظهر لهم بالاجتهاد) فيما لا نص فيه.
وضرب لذلك أمثلة:
\begin{itemize}
    \item \textbf{استقبال القبلة:} من كان معاينًا للكعبة فهو يصيب الحق ظاهراً وباطناً بيقين. ومن غاب عنها فهو مأمور بالاجتهاد بالعلامات والنجوم، فإذا اجتهد وتوجه فقد أدى ما عليه، وهو مصيب في الظاهر، ولا يعلم الباطن (حقيقة الإصابة الدقيقة) إلا الله.
    \item \textbf{الحكم بعدالة الشهود:} نحن مأمورون بقبول شهادة العدل. ولا يُعرف العدل إلا بما يظهر لنا من صلاحه وملازمته للتقوى والمروءة، فنحكم بعدالته ظاهراً، وقد يكون في الباطن بخلاف ذلك، لكننا لم نُكلف بعلم الغيب.
    \item \textbf{الاجتهاد في جزاء الصيد:} في قوله تعالى ﴿فَجَزَاءٌ مِّثْلُ مَا قَتَلَ مِنَ النَّعَمِ﴾، يحكم العدلان بأقرب الأشياء شبهاً بالصيد المقتول من النعم، وهذا التقريب هو من باب الاجتهاد والقياس.
\end{itemize}

\subsection{أجر المجتهد}
واستدل بحديث عمرو بن العاص رضي الله عنه: «إذا حكم الحاكم فاجتهد ثم أصاب فله أجران، وإذا حكم فاجتهد ثم أخطأ فله أجر». فالمجتهد مأجور على بذل جهده لإعمال النصوص، فإن وافق الحق عند الله (الباطن) نال أجرين، وإن وافق الظاهر عنده وأخطأ الباطن نال أجراً واحداً.

\section{إبطال الاستحسان}
حذر الشافعي تحذيراً شديداً من القول بالاستحسان المجرد عن الدليل والقياس، وقرر أن أحداً لا يجوز له أن يفتي في دين الله إلا من جهة العلم (الكتاب، السنة، الإجماع، القياس).
فلا يجوز أن يُقال بالاستحسان (وهو القول بالرأي والتشهي دون مستند شرعي). وضرب لذلك مثلاً دنيوياً: إذا أتلف شخص سلعة لآخر، فلا يُلجأ في تقييمها إلا لخبير خابر بالسوق يقيسها بمثيلاتها، ولا يجوز له أن يلقي سعراً بالتعسف والاستحسان. فإذا كان هذا في أموال الدنيا، فكيف بحلال الله وحرامه؟
فالاستحسان المجرد ما هو إلا تلذذ واتباع للهوى، والواجب على العالِم ألا يقول إلا من جهة العلم، ولا يحل له أن يتخبط في دين الله برأيه المجرد.
\chapter{شروط المجتهد والمفتي وأنواع القياس}

\section{شروط الاجتهاد والفتوى}
الحمد لله والصلاة والسلام على رسول الله، أما بعد؛ فقد عقد الإمام الشافعي رحمه الله تعالى باباً نفيساً في "شروط المجتهد" ومن يحق له أن يتصدر للفتوى والقياس. وقد قال شيخ شيوخنا العلامة أحمد شاكر رحمه الله عن هذه الفِقرات: «هذه الدرر الغالية والحكم البالغة هي أحسن ما قرأت في شروط الاجتهاد».

قال الشافعي رحمه الله: «ولا ينبغي للمفتي أن يفتي أحداً إلا أن يكون عالماً بكتاب الله، وعلم ناسخه ومنسوخه، وخاصه وعامه، وأدبه. وعالماً بسنن رسول الله ﷺ، وأقاويل أهل العلم قديماً وحديثاً (مواطن الإجماع والاختلاف). وعالماً بلسان العرب، عاقلاً يميز بين المشتبه ويعقل القياس».

فإن عُدم واحدة من هذه الخصال، لم يحل له أن يقول بقياس أو يفتي. ومثّل لذلك بمن يفتقد أداة من أدوات الفهم؛ كمن كان عالماً بالأصول ولكنه لا يعقل القياس، أو كان عاقلاً للقياس ومضيعاً لعلم الأصول؛ فكلاهما لا يجوز له أن يقيس. 
كما لا يجوز أن يوصف الطريق لأعمى فيقال له: «اجعل كذا عن يمينك وكذا عن يسارك»، وهو لا يبصر. وكما لا يجوز لجاهل بأسعار السوق أن يُقيّم سلعة، فكذلك لا يقاس في حلال الله وحرامه إلا بآلة متكاملة من العلم.

\subsection{الإنصاف وأدب المجتهد في الخلاف}
من أعظم شروط المجتهد التي نص عليها الشافعي: بلوغ غاية الجهد في البحث عن الحق، والإنصاف من نفسه.
فيجب على طالب العلم ألا يمتنع من الاستماع لمن خالفه؛ لأنه قد يتنبه بالاستماع إلى ترك الغفلة، ويزداد تثبيتاً فيما اعتقده من الصواب. فالمخالف إما أن يزيدك يقيناً على ما أنت عليه إذا ظهر لك ضعف دليله، وإما أن يردك إلى الصواب إذا كنت مخطئاً. والواجب أن تنصف من نفسك فتتبع الحق حيث ظهر. ولا يكون العالم أعمى مقلداً يقلد بغير حجة.

\section{القياس ومراتبه}
بعد أن بيّن الشافعي شروط القياس والمقيس، شرع في بيان أقسام القياس ومراتبه قوةً ووضوحاً:

\subsection{أقوى القياس (قياس الأولى)}
يقرر الشافعي أن أقوى أنواع القياس هو: أن يحرّم الله في كتابه أو يحرّم رسوله ﷺ القليل من الشيء، فيُعلم يقيناً أن الكثير منه أشد تحريماً، بفضل الكثرة على القلة. 
وكذلك إذا حُمد العبد على يسير من الطاعة، كان ما هو أكثر منها أولى أن يُحمد عليه. وإذا أبيح الكثير من الشيء، كان الأقل منه أولى أن يكون مباحاً.

\textbf{أمثلة وتطبيقات:}
\begin{itemize}
    \item قال تعالى: ﴿فَمَن يَعْمَلْ مِثْقَالَ ذَرَّةٍ خَيْرًا يَرَهُ * وَمَن يَعْمَلْ مِثْقَالَ ذَرَّةٍ شَرًّا يَرَهُ﴾. فما هو أكثر من مثقال ذرة من الخير أعظم حمداً وأجراً، وما هو أكثر من ذرة من الشر أعظم مأثماً.
    \item قال ﷺ: «إن الله حرم من المؤمن دمه وماله وأن يُظن به إلا خيراً». فإذا حرم مجرد سوء الظن المخفي، فالتصريح بالقول الباطل فيه والإيذاء الفعلي أشد تحريماً من باب أولى.
    \item أباح الله لنا دماء أهل الكفر المحاربين غير المعاهدين وأموالهم بأسرها، فما كان من نيل شيء من أبدانهم دون الدماء أو أخذ بعض أموالهم فهو أولى بالإباحة.
\end{itemize}

\section{أمثلة فقهية دقيقة على القياس (العلة والشبه)}

\subsection{القياس في وجوب نفقة الأب العاجز}
قال تعالى: ﴿وَالْوَالِدَاتُ يُرْضِعْنَ أَوْلَادَهُنَّ حَوْلَيْنِ كَامِلَيْنِ ۖ لِمَنْ أَرَادَ أَن يُتِمَّ الرَّضَاعَةَ ۚ وَعَلَى الْمَوْلُودِ لَهُ رِزْقُهُنَّ وَكِسْوَتُهُنَّ بِالْمَعْرُوفِ﴾. 
دلت الآية والسنة (كحديث هند بنت عتبة: «خذي من ماله ما يكفيك وولدك بالمعروف») على أن نفقة الولد الصغير واجبة على الوالد، والعلة في ذلك أن الولد بضعة من أبيه، وأنه عاجز في حال صغره عن الكسب وإغناء نفسه. 

فقاس الشافعي على ذلك: أنه إذا بلغ الأب من الكِبَر والعجز مبلغا لا يغني فيه نفسه بكسب ولا مال، وولده غني قادر؛ وجبت النفقة لِلأب على ولده. لأن الولد من الوالد، فكما أُجبر الأب على نفقة ابنه الصغير لعجزه، يُجبر الابن على نفقة أبيه لعجزه؛ فلا يُضيّع أحدهما ما هو بضعة منه.

\subsection{القياس في قاعدة الخراج بالضمان}
قضى النبي ﷺ في رجل اشترى عبداً فاستغله (أي جعله يعمل ويُدِرُّ عليه مالاً)، ثم اكتشف فيه المشتري عيباً قديماً فرده على بائعه الأول، فطلب البائع غلة العبد (أجرته في تلك الأيام). فقضى النبي ﷺ بأن العبد يُرد بالعيب، وأن للمشتري أن يحبس الغلة لنفسه ولا يردها؛ لأنه كان «ضامناً» للعبد، فلو مات في تلك الفترة لمات من مال المشتري، فقوبل الضمان باستحقاق الخراج (الخراج بالضمان).

قاس الشافعي على هذه السنة الجليلة أشياء أخرى لم تُنص بعينها:
فمن اشترى نخلة فأثمرت، أو جارية فولدت، أو ماشية فحلبت، ثم ظهر عيب قديم فردها المشتري؛ فإن الثمر والولد واللبن من حق المشتري، لأن كل هذا حدث في فترة ملكه وضمانه، فلحق بالأصل المُقاس عليه وهو غلة العبد.

\subsection{القياس في الربويات}
نهى النبي ﷺ عن التفاضل في الذهب بالذهب، والفضة بالفضة، والتمر والبر والشعير. 
والعلة في المأكولات الثلاثة أنها مأكولة تُكال وتُدخر للاقتيات. فقاس الشافعي والجمهور عليها: العسل، والسمن، والزيت، والسكر، وغيرها مما يقاس بالوزن أو الكيل ويكون مأكولاً مدخراً. فهذه لا تُباع بجنسها إلا متماثلة ويداً بيد قياساً على المنصوص عليه. ومنع قياس المأكولات على الذهب والفضة (الدنانير والدراهم)؛ لاختلاف العلة والجنس، فالدنانير أثمان للأشياء في أنفسها.

\subsection{اشتراك العاقلة في دية الخطأ وما دونها}
أجمع أهل العلم وتواترت السنة أن دية قتل الخطأ للمسلم الحر (وهي مائة من الإبل) تتحملها عاقلة الجاني (أي عصبته وقرابته)، وتُقسّط عليهم في ثلاث سنين.
وأجمعوا على أن الجناية العمد تكون في مال الجاني خاصة لا تشاركه عاقلته.

ولكن اختلفوا فيما لو جنى جناية خطأً لا تبلغ دية كاملة، بل تبلغ أقل من الثلث (كالموضحة وهي الشجة التي توضح العظم، وديتها خمس من الإبل). هل تتحملها العاقلة أم تكون في مال الجاني المخطئ وحده؟
قاس الشافعي رحمه الله ذلك قائلاً: إذا كان رسول الله ﷺ قد ألزم العاقلة بتحمل الجناية الكبرى (وهي الدية الكاملة) مواساةً لصاحبهم في الخطأ، فمن باب أولى أو مساواةً أن يتحملوا عنه الجناية الصغرى (كالنصف والثلث وما دونهما) قياساً على الأصل؛ لأنه إذا لزمهم الأكثر لزمهم الأقل في نفس العلة (وهي الخطأ)، وهذا من تمام القياس المطرد.
\chapter{حدود الإجماع والقياس، وأدب المناظرة والاختلاف}

\section{التحذير من ادعاء الإجماع في مسائل الخلاف}
الحمد لله والصلاة والسلام على رسول الله، أما بعد؛ فقد استكمل الإمام الشافعي رحمه الله تعالى حديثه في مناقشة مخالفيه، محذراً من التساهل في ادعاء "الإجماع" في مسائل يوجد فيها خلاف معتبر بين السلف.

يقول الشافعي راداً على من يدعي الإجماع لنصرة مذهبه: «لست أقول -ولا أحد من أهل العلم- هذا مجمع عليه، إلا لما لا تلقى عالماً أبداً إلا قاله لك، وحكاه عمن قبله؛ كالظهر أربع، وكتحريم الخمر، وما أشبه هذا. وقد أجده يقول: "المجمع عليه"، وأجد من أهل العلم كثيراً يقولون بخلافه! وأجد عامة أهل البلدان على خلاف ما يقول: المجمع عليه!». 
فالإجماع اليقيني هو الإجماع على ما عُلم من الدين بالضرورة، وينبغي الحذر الشديد من ادعاء الإجماع الباطل لتمرير الآراء.

\section{قواعد في ضمان العاقلة لجناية الخطأ}
ناظر الشافعي رحمه الله أحد فقهاء الأحناف في مسألة جناية الخطأ وما تتحمله "العاقلة" (وهم عصبة الرجل وقرابته).
فالقاتل خطأً تتحمل عاقلته دية المقتول (مائة من الإبل)، لكن هل تتحمل العاقلة الجراحات الصغيرة التي لا تبلغ ثلث الدية (كالموضحة، وما دونها)؟
قال المناظر: لا تعقل العاقلة ما دون الموضحة؛ لأن رسول الله ﷺ لم يقض فيما دون الموضحة بشيء.
رد الشافعي بإلزام دقيق: إذا لم يقض فيها بشيء، فهل أهدرها؟ لا، بل الجراح مضمونة. وإذا كانت العاقلة تغرم الأكثر (الدية الكاملة)، فما المانع أن تغرم الأقل؟ فالقاعدة المطردة أن جميع ما كان جناية خطأ فعلى العاقلة، ولو كان درهماً يسيراً، طالما أن السبب هو "الخطأ". 
أما ما لزم الإنسان في غير جناية الخطأ (كالزكاة، والمهر، والكفارة، وجناية العمد، وما أفسدته مواشيه)، فكل هذا في ماله الخاص ولا تتحمله العاقلة أبدًا.

\section{مناظرات تطبيقية في القياس}

\subsection{دية العبد المقتول خطأً}
اختلفوا في جناية الحر على العبد خطأً؛ فقال الأحناف وأهل المدينة: "جراح العبد في ثمنه كجراح الحر في ديته" (أي تُحسب كنسبة مئوية). وخالفهم الشافعي، قائلاً إن جراح العبد تُقدر بما نقص من ثمنه (كأي سلعة تُتلف).
طالب المخالفُ الشافعيَّ بالدليل، ثم استدل المخالف بقول سعيد بن المسيب.
فرد عليه الشافعي بإنصاف وذكاء: أنت في أصولك لا ترى أقوال التابعين حجة، فكيف تحتج بها هنا؟ ثم أقام الشافعي القياس الصحيح مبيناً أن العبد يُفارق الحر في أن دية الحر مؤقتة شرعاً، بينما العبد ديته ثمنه، فأشبه السلع والبهائم المضمونة بقيمتها. وساق إلزاماً عقلياً وشرعياً للمخالف حتى أثبت تناقضه في قياسه.

\subsection{جواز استقراض الحيوان ورده بأحسن منه}
أنكر المخالف أن تُشترى الإبل بصفة في الذمة إلى أجل (كأن يستقرض بعيراً). فرد الشافعي بحديث أبي رافع رضي الله عنه: «أن النبي ﷺ استسلف من رجل بعيراً، فجاءته إبل فصرفه أن يقضيه إياه، فلم يجد إلا جملاً خياراً رباعياً، فقال ﷺ: أعطه إياه، فإن خيار الناس أحسنهم قضاءً». فبطل إنكارهم بالسنة الصريحة الجلية.

\section{الرخص والاستثناءات لا يُقاس عليها}
من درر القواعد الأصولية التي أرساها الشافعي رحمه الله: أن ما استُثني من أصل محرم للحاجة (الرخصة)، لا يجوز أن يُجعل أصلاً يُقاس عليه غيره.

\textbf{المثال الأول: المسح على الخفين:}
أصل فرض الله في الوضوء غسل القدمين، فلما مسح النبي ﷺ على الخفين، أخذناه اتباعاً ورخصةً لمن لبسهما على طهارة. ولم يجز لنا أن نقيس على الخفين غيرهما؛ فلا يجوز المسح على العمامة أو البرقع أو القفازين قياساً على الخفين، بل نبقي الأصل (الغسل) فيما عدا المستثنى بالنص.

\textbf{المثال الثاني: بيع العرايا:}
نهى النبي ﷺ عن بيع الثمر بالتمر (المزابنة) لوجود التفاضل والجهالة وكونه ربا. ولكنه رخص لبعض الفقراء الذين لا يملكون نقوداً أن يشتروا الرطب على النخل بخرصه من التمر ليأكلوه (في حدود خمسة أوسق)، وهذه هي "العرايا".
فالنهي عام ومستقر، والرخصة في العرايا خاصة، فلا يجوز أن نضرب الحديثين ببعضهما، ولا أن نقيس بيوع غرر أخرى على رخصة العرايا.

\section{أحكام تعبدية وأخرى معقولة المعنى}

\subsection{دية الجنين (الغرة)}
قضى النبي ﷺ في الجنين إذا أسقطته أمه ميتةً بجناية بغرة (عبد أو أمة، وتُقوّم بخمس من الإبل).
قال الشافعي: هذا حكم فارق حكم النفوس الأحياء؛ فلو سقط حياً ثم مات، ففيه دية كاملة (مائة إبل للذكر، وخمسون للأنثى). ولما حكم النبي ﷺ بالغرة للجنين الميت ولم يسأل أذكر هو أم أنثى؟ دل على أنه حكم "تعبدي"، سنه النبي ﷺ لحكمة يعلمها الله، فنسلم له، ولا نقيس عليه غيره. 
أما الأحكام المعقولة المعنى (التي ندرك علتها)، فنطيع الله فيها ونقيس عليها ما اشترك معها في ذات العلة.

\subsection{قاعدة الخراج بالضمان وحديث المصراة}
قضى النبي ﷺ أن «الخراج بالضمان»؛ فمن اشترى عبداً فاستغله وربح منه، ثم وجد فيه عيباً فرده، فإن المشتري يأخذ الربح ولا يرده للبائع؛ لأنه كان ضامناً للعبد لو مات. فقاس الشافعي على ذلك ثمر النخل وولد الجارية ولبن الماشية المشتراة.

لكنه استثنى "المصراة" (الشاة التي يُربط ضرعها ليجتمع اللبن فيغتر المشتري)، فإن النبي ﷺ أمر بردها ورد "صاع من تمر" معها كبدل للبن. فهنا نقف عند النص ولا نقيسه، ونعمل بالقياس العام في غير المصراة.

\section{أنواع الاختلاف (المحرم والسائغ)}
قال السائل: إني أجد أهل العلم قديماً وحديثاً مختلفين، فهل يسعهم هذا الاختلاف؟
قال الشافعي: الاختلاف من وجهين:
\begin{enumerate}
    \item \textbf{اختلاف محرم:} وهو ما أقام الله به الحجة في كتابه، أو على لسان نبيه ﷺ نصاً بيناً لا يحتمل التأويل؛ فهذا لا يحل الاختلاف فيه لمن علمه. (مثل من يعارض النصوص القطعية بالآراء العقلانية المحدثة).
    \item \textbf{اختلاف سائغ:} ما كان يحتمل التأويل ويُدرك قياساً، فذهب المتأول أو القايس إلى معنى يحتمله الخبر. فهذا لا أضيق فيه على المخالف كما أضيق في المنصوص، وهو من باب الاجتهاد المأجور صاحبه.
\end{enumerate}
وقد ذم الله التفرق بعد مجيء البينات ﴿وَمَا تَفَرَّقَ الَّذِينَ أُوتُوا الْكِتَابَ إِلَّا مِن بَعْدِ مَا جَاءَتْهُمُ الْبَيِّنَةُ﴾.

\section{معنى القرء في اعتداد المطلقة}
اختلف الصحابة رضي الله عنهم في تفسير قول الله تعالى: ﴿وَالْمُطَلَّقَاتُ يَتَرَبَّصْنَ بِأَنفُسِهِنَّ ثَلَاثَةَ قُرُوءٍ﴾. 
فقالت عائشة وزيد بن ثابت وابن عمر: الأقراء هي "الأطهار". 
وقال أبو بكر وعمر وعلي وابن عباس: الأقراء هي "الحيض".

وبيّن الشافعي أن اللفظ لغةً يحتمل الأمرين، لأن "القرء" يعني "الوقت" ويطلق عليهما معاً. لكن الشافعي رجح أن المراد به (الطهر)، مستدلاً باللغة؛ فمصدر (قرأ) بمعنى (جمع وحبس)، والرحم في حال الطهر يجمع الدم ويحبسه، وفي حال الحيض يرخيه ويرسله، فالطهر أولى باسم القرء.
كما استدل بالسنة؛ حيث أمر النبي ﷺ ابن عمر أن يطلق امرأته طاهراً من غير جماع، وقال: «فتلك العدة التي أمر الله أن يُطلق لها النساء».
وهذا مثال على الاختلاف السائغ الذي لا يُعنف فيه المخالف لتوفر أدلة محتملة للطرفين.
\chapter{عدة الحامل، وأحكام الإيلاء، وميراث الجد (ختام الكتاب)}

\section{عدة الحامل المتوفى عنها زوجها}
الحمد لله رب العالمين، والصلاة والسلام على نبينا محمد وعلى آله وصحبه أجمعين. أما بعد؛ فهذا هو الدرس الختامي من دروس شرح كتاب «الرسالة» للإمام أبي عبد الله الشافعي رحمه الله تعالى.

بدأ الشافعي رحمه الله بذكر مسألة اختلف فيها الصحابة قديماً، ليوضح كيف حسمت السنة هذا الخلاف. قال الله عز وجل: ﴿وَالْمُطَلَّقَاتُ يَتَرَبَّصْنَ بِأَنفُسِهِنَّ ثَلَاثَةَ قُرُوءٍ﴾، وقال: ﴿وَاللَّائِي يَئِسْنَ مِنَ الْمَحِيضِ مِن نِّسَائِكُمْ إِنِ ارْتَبْتُمْ فَعِدَّتُهُنَّ ثَلَاثَةُ أَشْهُرٍ وَاللَّائِي لَمْ يَحِضْنَ ۚ وَأُولَاتُ الْأَحْمَالِ أَجَلُهُنَّ أَن يَضَعْنَ حَمْلَهُنَّ﴾، وقال: ﴿وَالَّذِينَ يُتَوَفَّوْنَ مِنكُمْ وَيَذَرُونَ أَزْوَاجًا يَتَرَبَّصْنَ بِأَنفُسِهِنَّ أَرْبَعَةَ أَشْهُرٍ وَعَشْرًا﴾.

اختلف بعض أصحاب النبي ﷺ في الحامل المتوفى عنها زوجها؛ هل تعتد بأبعد الأجلين (أي تنتظر حتى تضع حملها وتمر عليها أربعة أشهر وعشراً)؟ أم تنقضي عدتها بوضع الحمل فقط؟ 
ظاهر القرآن يحتمل المعنيين، فدلت سنة رسول الله ﷺ على أن وضع الحمل هو آخر العدة مطلقاً. فقد روى الزهري أن سُبيعة الأسلمية رضي الله عنها وضعت حملها بعد وفاة زوجها بليالٍ قليلة، فمر بها أبو السنابل بن بعكك رضي الله عنه ورآها قد تصنعت للخطاب، فأنكر عليها وقال: إنها تعتد أربعة أشهر وعشرا. فذكرت ذلك لرسول الله ﷺ، فقال: «كَذَبَ أبو السنابل، قد حللتِ فتزوجي». 
(ومعنى "كذب" في لغة أهل الحجاز: أخطأ، أي أخبر بخلاف الواقع). فدلت السنة على أن وضع الحمل ينهي العدة فوراً.

\section{تقديم السنة على أقوال الرجال}
يستنبط الشافعي قاعدة عظيمة من هذه المسألة: "ما دلت عليه السنة فلا حجة في أحد خالف قوله السنة، ولو كان من أصحاب رسول الله ﷺ". فلا يحل لمسلم أن يقدم قول أحد من الناس على كلام الله وكلام رسوله ﷺ، أو على إجماع سلف الأمة. 

\section{أحكام الإيلاء (اليمين على ترك الجماع)}
ثم انتقل الشافعي لمسألة ليس فيها نص سنة مباشر، وإنما يُستدل عليها من القرآن واستنباطاً، وهي مسألة "المُولي" (الذي يحلف ألا يجامع امرأته).
قال تعالى: ﴿لِّلَّذِينَ يُؤْلُونَ مِن نِّسَائِهِمْ تَرَبُّصُ أَرْبَعَةِ أَشْهُرٍ ۖ فَإِن فَاءُوا فَإِنَّ اللَّهَ غَفُورٌ رَّحِيمٌ * وَإِنْ عَزَمُوا الطَّلَاقَ فَإِنَّ اللَّهَ سَمِيعٌ عَلِيمٌ﴾.

اختلف الصحابة؛ هل بمجرد انقضاء الأشهر الأربعة تطلق المرأة تلقائياً؟ أم يوقف الزوج ويُخيّر؟
رجح الشافعي أنه يوقف، فلا يلزمه الطلاق تلقائياً، بل إذا مضت الأشهر الأربعة قيل له: إما أن تفيء (تجامع) وإما أن تطلق. لأن الله خيره بين الفيئة والطلاق بعد مضي المدة، فلو كانت تطلق بمجرد انقضاء المدة لما كان هناك معنى للتخيير بعدها. والفيئة تكون بالجماع لمن قدر عليه، وهو يمين جعل الله له مدة لئلا تُضار الزوجة، فإن مضت أُجبر الزوج على أحد الأمرين.

\section{مسائل في المواريث (رد الفضل، وميراث الجد)}

\subsection{رد فضل المواريث}
اختلفوا فيما لو مات إنسان وترك أختاً، فلها النصف. فماذا يُفعل بالنصف الآخر؟
ذهب زيد بن ثابت رضي الله عنه والشافعي إلى أن النصف الباقي يذهب لجماعة المسلمين (بيت المال)، ولا يُرد عليها؛ لأن الله فرض لها النصف منفردة أو مع أخواتها، فلو رددنا عليها النصف الباقي لكنا ورثناها غير ما فرض الله لها. وذهب آخرون إلى الرد على ذوي الأرحام لقوله تعالى: ﴿وَأُولُو الْأَرْحَامِ بَعْضُهُمْ أَوْلَىٰ بِبَعْضٍ فِي كِتَابِ اللَّهِ﴾.

\subsection{ميراث الجد مع الإخوة}
اختلف الصحابة في الجد مع الإخوة؛ فذهب أبو بكر وابن عباس رضي الله عنهم إلى أن الجد كالأب يحجب الإخوة تماماً ويأخذ الميراث. وذهب زيد بن ثابت وعمر وعلي وابن مسعود رضي الله عنهم إلى أن الجد يُقاسم الإخوة ولا يحجبهم (بشروط وتفصيلات معروفة في الفرائض).
ورجح الشافعي مذهب المشاركة؛ وناظر من قال بالحجب بأن الجد يدلي بقرابة غيره (أبو أب الميت) والأخ يدلي بقرابة غيره (ابن أب الميت)، وكلاهما يدلي بقرابة الأب، فليس أحدهما أولى بالحجب من الآخر، فكان التشريك أوفر بالقياس وأقرب للعدل.

\section{ختام كتاب الرسالة}
ختم الإمام الشافعي كتابه بالتأكيد على أن العلم مراتب، وأن الاجتهاد يكون بالدليل، وأن التوفيق بيد الله.
وقد ختم الربيع بن سليمان (تلميذ الشافعي وراوية كتابه) نسخته قائلاً: "أجاز الربيع بن سليمان صاحب الشافعي نسخ كتاب الرسالة... في ذي القعدة سنة خمس وستين ومائتين وكتب الربيع بخطه".

ونحمد الله العلي القدير أن يسر لنا قراءة ومدارسة هذا الكتاب العظيم للإمام الشافعي طيب الله ثراه، ونسأله بأسمائه الحسنى وصفاته العلى أن يرزقنا العلم النافع والعمل الصالح، وأن يقسم لنا من خشيته ما يحول به بيننا وبين معاصيه. وصلى الله وسلم وبارك على نبينا محمد وعلى آله وصحبه أجمعين.

\appendix
\chapter{الملاحق: الفوائد المستخلصة من الدروس}

\section{الفوائد المنهجية والتربوية}
\begin{itemize}
    \item \textbf{الاستمرار في المدارسة:} ضرب الإمام المزني أروع الأمثلة في الدأب على العلم، حيث مكث يقرأ كتاب "الرسالة" خمسين سنة، وفي كل مرة يستفيد شيئاً جديداً.
    \item \textbf{فقه القلوب وشكر النعم:} لا يملك العبد أن يشكر الله إلا بتوفيق الله، وهذا التوفيق في ذاته نعمة تستوجب شكراً جديداً، مما يبقي قلب المؤمن معلقاً بالله متذللاً له.
    \item \textbf{آداب طالب العلم:} لخصها الشافعي في أربعة أركان: بذل الجهد، الصبر على المكاره والعوارض، الإخلاص التام لله، والتضرع وطلب العون والفهم من الله وحده.
    \item \textbf{الاعتزاز بالهوية:} وجوب اعتزاز المسلم بسمته ولغته العربية، فإن المتابعة في الهدي الظاهر (كاللغة واللباس) تورث الانتماء والاهتداء في الباطن، والذوبان في لغات وعادات المشركين عنوان للهزيمة النفسية.
    \item \textbf{مقاصد القرآن الكريم:} القرآن الكريم هو كتاب هداية وإرشاد للتي هي أقوم، وليس كتاباً متخصصاً في العلوم الدنيوية البحتة كالفيزياء والكيمياء، بل أُمِرنا فيها بالرجوع لأهل الاختصاص.
    \item \textbf{النجاة في السنة:} الفوز بالجنة منوط بالاعتصام بسنة النبي صلى الله عليه وسلم والتمسك بها والدعوة إليها والحذر التام من مخالفتها.
    \item \textbf{أهمية اللغة العربية:} فهم مراد الله ورسوله يتوقف على فهم لغة العرب، والجهل بها يؤدي إلى الظن بوجود تعارض وتناقض بين النصوص الشرعية.
    \item \textbf{طاعة النبي صلى الله عليه وسلم طاعة لله:} لا يمكن ادعاء اتباع القرآن مع إنكار السنة، فهذا تناقض ويكذب القرآن نفسه الذي أمر بطاعة الرسول طاعة مطلقة.
    \item \textbf{خطورة القول بلا علم:} أصل الشافعي قاعدة ذهبية بأن الواجب على العلماء ألا يقولوا إلا بدليل وعلم راسخ، وأن الكثير ممن خاضوا في العلم والدين لو سكتوا لكان أسلم لهم وللأمة.
    \item \textbf{نصيحة للعوام:} الواجب على العامي أن يلزم كبار العلماء الموثوقين، وأن يصدق مع الله في طلب الحق، فمن صدق واستفتى قلبه بصدق دله الله على الحق وأعانه عليه.
    \item \textbf{شروط الاجتهاد الجماعي:} يتطلب الوصول إلى الحق علماء أتقياء أوفياء لدينهم، متجردين عن العصبية والتقليد واتباع الهوى والشهوات، غايتهم نصرة السنة.
    \item \textbf{خطورة التكلف:} من تكلف الفتيا وتكلم بجهل وخالف مسلك السلف، فإن وافق الحق مصادفة لم يُحمد لفساد مسلكه، وإن أخطأ لم يُعذر. أما من سلك منهج السلف فإنه مأجور على اجتهاده وإن أخطأ.
    \item \textbf{أهمية النصيحة:} بذل النصيحة لعامة المسلمين وخاصتهم فرض ودين، ومن ترك إيضاح الحق والنصح للمسلمين خشية الناس فقد سفه نفسه وترك حظه من الخير.
    \item \textbf{عدم العصمة وذم التقليد:} أخطأ الشافعي بسبق لسان في الاستدلال بآية ﴿فَآمِنُوا بِاللَّهِ وَرَسُولِهِ﴾ في سورة النساء والصواب ﴿وَرُسُلِهِ﴾، وتناقلها النساخ لقرون دون تنبه ثقةً في الإمام، وهو تطبيق عملي لقول الشافعي نفسه "وبالتقليد أغفل من أغفل منهم"، فالعصمة للأنبياء فقط.
    \item \textbf{حدود طاعة أولي الأمر:} طاعة أولي الأمر ليست مطلقة بل مستثناة ومقيدة بالمعروف وبطاعة الله ورسوله، وعليهم إقامة شرع الله في الأرض وحفظ أعراض المسلمين ودمائهم.
    \item \textbf{رحمة الشريعة في الحدود:} الشريعة تحتاط للدماء والأعراض احتياطاً بالغاً، ففي قصة العسيف لم يرجم النبي ﷺ المرأة بمجرد اعتراف الشريك، بل أرسل من يسألها، فلو أنكرت لدرئ عنها الحد.
    \item \textbf{الحكم بالظاهر لا بالبواطن:} القاضي والحاكم ملزم بالحكم بالظاهر (البينة أو الإقرار)، ولا يجوز له الحكم بعلمه الشخصي أو بالقرائن الباطنة، كما أمضى النبي ﷺ حكم اللعان ودرأ الحد رغم ظهور علامات الشبه في الولد.
    \item \textbf{تحكيم السنة لا العواطف:} الميزان الذي توزن به الأقوال والأفعال هو كتاب الله وسنة رسوله ﷺ ومنهج السلف، ولا يجوز التلاعب بالدين باسم العواطف، أو التجارب، أو العقل المصادم للشرع.
    \item \textbf{العذر بالجهل للعلماء:} العالم إذا خالف السنة الصحيحة فإنما يخالفها لعدم بلوغها إياه، وهو معذور مأجور لا يُطعن فيه، لكن لا يجوز للمسلم المقلد أن يترك السنة التي بلغته تعصباً لقول العالم.
    \item \textbf{فقه جمع الطرق:} لا يمكن فهم السنة فهماً صحيحاً وتفسير الاختلاف الظاهري فيها إلا بجمع طرق الحديث وألفاظه لمعرفة أسباب الورود والسياق، وعدم الاكتفاء بظاهر لفظ واحد.
    \item \textbf{حفظ السنة وتكاملها بين الصحابة:} تتفاوت مستويات حفظ الصحابة للسياق؛ فمنهم من نقل النهي المجرد، ومنهم من نقل الرخصة، ومنهم من نقل الناسخ والمنسوخ وعلة الحكم (كحديث عائشة في لحوم الأضاحي)، مما يبرز حكمة تعدد الروايات.
    \item \textbf{الإنصاف عند الاختلاف السائغ:} يُعذر المخالف ولا يُعنف إذا استند إلى دليل شرعي صحيح (كما في اختلاف صيغ التشهد)، بخلاف من يصادم النصوص بالآراء المحدثة.
    \item \textbf{عدم الاغترار بالكثرة:} الكثرة لا اعتبار بها إذا خالفت النص الشرعي، ولا يُنال نصر الله والتمكين بمداهنة المخالفين وترك الثوابت.
    \item \textbf{التأدب مع الشيوخ عند الرد:} إذا ظهر الدليل بخلاف قول الشيخ (كرد الشافعي على سفيان بن عيينة)، وجب اتباع الدليل، ويكون الرد بنقاش علمي رصين وتأصيل تاريخي دون تنقيص أو إساءة أدب.
    \item \textbf{أقسام العلم الشرعي:} ينقسم العلم إلى فرض عين لا يعذر مكلف بجهله (كأصول الفرائض والمحرمات)، وفرض كفاية تتكفل به طائفة من الأمة (كالتفقه الدقيق وتفاريع الأحكام).
    \item \textbf{ثبات المنهج ونصرة الدين:} نصرة الدين والتمكين له لا تنال بالمعاصي ولا بالوسائل المخالفة للشرع، فالغاية المشروعة تفتقر إلى وسيلة مشروعة.
    \item \textbf{الاحتياط للدين:} الخطأ في الرواية عن النبي ﷺ يترتب عليه تحليل الحرام وتحريم الحلال في الأمة، لذا كان احتياط المحدثين وتشددهم في شروط الرواة أعظم من احتياط القضاة في قبول الشهود.
    \item \textbf{نضارة الوجه لمبلغ السنة:} وعد النبي ﷺ بنضارة الوجه وحسنه لمن يسمع سنته ويحفظها ويبلغها للناس، مما يستوجب الحرص على نشر السنن.
    \item \textbf{الحفظ والفقه:} ليس كل راوٍ وحافظ للحديث فقيهاً، وربما نقل الفقه لمن هو أفقه منه، والكمال في الجمع بين الحفظ والفقه في الدين.
    \item \textbf{الرجوع للحق واتباع الدليل:} من أعظم خصال العالم الرجوع عن قوله أو قضاءه (كما فعل عمر وابن عباس وابن عمر) بمجرد بلوغه الدليل الصحيح عن رسول الله ﷺ، دون تعصب لرأي مسبق.
    \item \textbf{المعنى الصحيح للزوم الجماعة:} الجماعة المأمور بلزومها ليست لزوم الأبدان، بل هي لزوم منهج السلف وما أجمعوا عليه من عقيدة وأحكام، ولا عبرة باجتماع الأبدان إن اختلفت العقائد.
    \item \textbf{التكليف بالظاهر لا الغيب:} كلف الله العباد بالحكم على الظاهر في الشهادات والاجتهاد، ووكل السرائر والبواطن لنفسه سبحانه، ولا يكلف الله نفساً علم الغيب.
    \item \textbf{الإنصاف وقبول الحق:} يجب على المجتهد وطالب العلم أن ينصف من نفسه، ولا يمتنع من الاستماع للمخالف، فإما أن يزداد تثبيتاً على الحق، وإما أن يتبين له خطؤه فيرجع للصواب.
    \item \textbf{احترام التخصص الشرعي:} لا يجوز لمن لم يجمع آلات الاجتهاد (من علم الكتاب والسنة واللغة والقياس) أن يفتي أو يقيس، كما لا يُسأل الخياط عن تقييم البناء وأسعار السوق.
    \item \textbf{التحذير من إدعاء الإجماع:} لا يصح ادعاء الإجماع المطلق إلا فيما يُعلم من الدين بالضرورة ولا يُعلم له مخالف، والتهور في ادعاء الإجماعات في المسائل الخلافية لتمرير الآراء مسلك مذموم.
    \item \textbf{أدب المناظرة:} من فقه المناظر إلزامه لخصمه بالحجة من داخل أقواله نفسها، وبيان التناقض المنهجي لمن يحتج بأقوال لا يراها حجة في أصوله لمجرد نصرة مذهبه.
\end{itemize}

\section{قواعد المناظرة والرد على المخالف}
\begin{itemize}
    \item \textbf{الرد على محرفي النصوص:} بيان أن القرآن الكريم عربي خالص، وأن دعوى وجود العجمة فيه باطلة، والرد على من يحاول تطويع النصوص الشرعية لإرضاء أعداء الأمة.
    \item \textbf{تأصيل الغاية والوسيلة:} الغاية الشرعية (التمكين للدين) لا تُنال بوسائل غير شرعية ولا بمعصية الله، ولا عبرة بكثرة المخالفين إذا خالفوا منهج السلف.
    \item \textbf{الرد على القرآنيين:} ادعاء الاكتفاء بالقرآن وترك السنة هو في الحقيقة تكذيب للقرآن؛ لأن القرآن أمر باتباع النبي صلى الله عليه وسلم وطاعته.
    \item \textbf{المطالبة بالدليل الشرعي:} لا يحق لأحد غير النبي صلى الله عليه وسلم أن يُلزم الناس بقول إلا بدليل شرعي، والتقليد الأعمى مذموم.
    \item \textbf{آداب الاختلاف:} التفرقة بين اختلاف الأفهام الذي يسع الأمة (ألا يستقيم أن نكون إخواناً وإن لم نتفق في مسألة؟) وبين الاختلاف المصادم للنصوص والذي يُعد هدماً للشريعة (كاختلاف التضاد في الثوابت).
    \item \textbf{ذم التقليد الأعمى:} التقليد بغير بصر يعمي البصيرة، ولا يقلد تعصباً إلا غبي أو ذو هوى، وطالب العلم مطالَب بمعرفة الدليل ومأخذ العالم قبل اتباعه.
    \item \textbf{دحض شبهة الاكتفاء بالقرآن:} من ادعى الاكتفاء بالقرآن وترك السنة فقد كذب القرآن الذي أمره بالرد للرسول ﷺ عند التنازع في قوله ﴿فَرُدُّوهُ إِلَى اللَّهِ وَالرَّسُولِ﴾، وأن عدم التسليم لقضائه ينفي أصل الإيمان.
    \item \textbf{خطر التعصب المذهبي:} رد الأحاديث الصحيحة بدعوى أنها منسوخة أو مؤولة لنصرة المذهب المقلَّد هو مسلك خطير يؤدي إلى إضاعة السنن وتسهيل اختراق الشريعة بالقوانين الوضعية.
    \item \textbf{الرد العملي على القرآنيين:} إثبات أن الصلاة والصيام والحج والبيوع ذُكرت مجملة في القرآن، ولولا السنة المطهرة لما عرف المسلمون تفاصيل دينهم؛ فمن أنكر السنة هدم الدين كله.
    \item \textbf{الرد على دعوى تعارض السنة مع القرآن:} الادعاء بأن إباحة القرآن للزواج ممن لم يُذكرن في آيات المحرمات يتعارض مع تحريم السنة للجمع بين المرأة وعمتها هو ادعاء باطل؛ فالسنة مقيِّدة ومخصصة وليست معارضة.
    \item \textbf{الرد على العقلانيين:} دعوى السلفية مع تبني العقلانية وتقديم العقل المحدود على النقل هي طعن في المنهج السلفي ومسلك يؤدي إلى العلمنة، فالعقل الصريح لا يخالف النقل الصحيح أبداً.
    \item \textbf{تطبيق الشريعة على أهل الكتاب:} تخضع الأقليات غير المسلمة (كاليهود والنصارى) في ديار الإسلام للنظام العام وقانون العقوبات الإسلامي، كما رجم النبي ﷺ اليهوديين الزانيين لكونه ولي أمر المدينة.
    \item \textbf{أدب الرد على الشيوخ:} يجب على طالب العلم التزام الأدب عند الرد على شيوخه ومخالفتهم (كما فعل الشافعي في رده على شيخه سفيان بن عيينة)، ويكون الرد بالحجج العلمية والتاريخية الدقيقة.
    \item \textbf{مسلك الجمع مقدم على الترجيح:} لا يصار إلى ادعاء التعارض والاختلاف بين الأحاديث إلا بعد تعذر الجمع، فالجمع بين بيع العرايا والمزابنة، وبين السلم ونهي بيع الغائب، هو مسلك الفقهاء الراسخين.
    \item \textbf{الرد على استقلال العقل بالتشريع:} لا يُقاس النص على مجرد العقول، فالأصل في التشريع هو النقل الصحيح المتصل (الرواية)، والقياس الفقهي فرع عن النص لا مصادم له.
    \item \textbf{الرد على منكري حجية الآحاد:} الأحاديث والوقائع المتواترة في إرسال النبي ﷺ للآحاد (كمعاذ لليمن، وعلي لمنى)، وقبول الصحابة لخبر الواحد في نقض اليقين، تدحض شبهة منكري خبر الآحاد.
    \item \textbf{الرد على المتعصبة بالآحاد:} رد السنة المتصلة الصحيحة مع الاحتجاج بالمراسيل المنقطعة متى وافقت المذهب هو من أبطل الباطل وعين التناقض.
    \item \textbf{إبطال الاستحسان:} الاستحسان المجرد عن الأدلة الشرعية والقياس الصحيح هو محض اتباع للهوى والتشهي، ولا يجوز الفتيا به في دين الله.
    \item \textbf{عدم القياس على الرخص:} الرخص الشرعية المستثناة من قواعد عامة لحاجة (كالمسح على الخفين والمصراة والعرايا) يُوقف فيها عند النص ولا تُتخذ أصلاً يُقاس عليه لمخالفتها القياس الجلي.
    \item \textbf{ضوابط الخلاف السائغ:} يُقبل الخلاف ويُعذر صاحبه في المسائل الاجتهادية المبنية على التأويل المحتمل، ويُذم ويُحرم الخلاف إذا صدم نصاً بيناً لا يقبل التأويل.
    \item \textbf{تقديم النص على الرجال:} لا تُقدم أقوال الرجال مهما علت مراتبهم (ولو كانوا من الصحابة) على السنة النبوية إذا صحت واتضحت دلالتها.
\end{itemize}

\section{الفوائد العقدية والإيمانية}
\begin{itemize}
    \item \textbf{الصفات الإلهية:} لا يبلغ الواصفون كُنه عظمة الله، والواجب الإيمان بما وصف الله به نفسه وبما وصفه به رسوله صلى الله عليه وسلم دون تكييف أو تمثيل.
    \item \textbf{كمال الدين وحفظه:} بخلاف الكتب السابقة التي حُرفت وضاعت أسانيدها (بينهم وبين أنبيائهم مئات السنين من الانقطاع)، تفردت الأمة الإسلامية بحفظ الإسناد المتصل إلى رب العالمين.
    \item \textbf{حكمة الابتلاء في الاجتهاد:} الله ابتلى طاعة العباد في الفروض المنصوصة، وابتلاهم أيضاً في بذل الجهد (الاجتهاد) للوصول إلى الحق في المسائل التي ليس فيها نص قاطع كاستقبال القبلة عند البعد عنها.
    \item \textbf{التحليل والتحريم حق لله:} لا يجوز كائناً من كان أن يُحل أو يُحرم إلا من جهة العلم الشرعي المتمثل في الوحيين، والافتئات على الله في التشريع كفر وظلم.
    \item \textbf{عصمة الأمة:} العالم الفرد قد ينسى أو يجهل (كما خفي على عمر حديث الاستئذان)، لكن مجموع الأمة معصوم عن الخطأ والضلال، ومنهج السلف معصوم بالكلية.
    \item \textbf{اقتران الإيمانين:} الإيمان بالله لا يصح ولا يقع عليه اسم الإيمان أبداً حتى يقترن بالإيمان برسوله محمد صلى الله عليه وسلم، وتصديقه في كل ما أخبر.
    \item \textbf{التسليم المطلق للشرع:} لا يثبت إيمان العبد حتى يسلم تسليماً مطلقاً لحكم رسول الله ﷺ ولا يجد في نفسه حرجاً مما قضى، وإن خالف العقل أو الهوى.
    \item \textbf{التدرج حكمة ورحمة:} التدرج الإلهي في تشريع الأحكام ونسخها (كالخمر والقتال والصلاة) هو من رحمته سبحانه وتوسعة على العباد لتستأنس النفوس للتشريع.
    \item \textbf{انتظار الوحي وكمال التسليم:} كان النبي ﷺ يتوقف في القضايا النازلة (كقضية رمي الزوجة بالزنا) ولا يحكم فيها من تلقاء نفسه حتى يأتيه الوحي من الله، كمالاً لافتقاره وعبوديته.
    \item \textbf{الولاء والبراء في الميراث:} لا يتوارث أهل ملتين، فلا يرث الكافر قريبه المسلم، ولا المسلم يرث الكافر، تمييزاً لصف الإيمان وإعلاءً لرابطة الدين على رابطة النسب.
    \item \textbf{كمال التسليم لله ورسوله:} لا يجوز معارضة السنة بالعقول والأهواء، ولا يُقال لحكم رسول الله ﷺ "لِمَ" أو "كيف" على وجه الاعتراض، والتسليم المطلق هو المحك الحقيقي للإيمان.
    \item \textbf{سيادة الشريعة على الجميع:} أحكام الشريعة الإسلامية نافذة على جميع رعايا الدولة المسلمة، بمن فيهم أهل الذمة، كما فعل النبي ﷺ في إقامة حد الرجم على اليهوديين الزانيين.
    \item \textbf{الغاية لا تبرر الوسيلة الباطلة:} إظهار الدين ونصرته يكون بقضاء الله وقدره، ويستوجب الأخذ بالأسباب الشرعية الصافية، واللجوء لوسائل محرمة أو كفرية لن يؤدي إلى التمكين بل للخذلان.
    \item \textbf{حجية الإجماع:} لزوم جماعة المسلمين أصل عقدي، والإجماع المتيقن حجة لازمة لا يجوز مخالفتها، وادعاء الإجماع المعاصر على ما يخالف إجماع السلف (كالمبادئ الوضعية) باطل ومردود.
    \item \textbf{التسليم للأحكام التعبدية:} الأحكام غير معقولة المعنى التي لا تُدرك علتها (كتعويض الجنين الميت بغرة) يُتعبد لله بتسليمها، بخلاف المعقولة التي يُقاس عليها.
\end{itemize}

\section{الفوائد التفسيرية واللغوية}
\begin{itemize}
    \item \textbf{الفرق بين الحمد والشكر:} الحمد أعم سبباً (يكون على النعمة وغيرها) وأخص متعلقاً (بالقلب واللسان فقط). والشكر أخص سبباً (على النعمة فقط) وأعم متعلقاً (بالقلب واللسان والجوارح).
    \item \textbf{معنى أزلفت:} أزلفت تأتي بمعنى (قربت) في حق الجنة، لكنها في سياق الاستغفار في خطبة الشافعي (لما أزلفت وأخرت) تأتي بمعنى ما قدّمتُ من أعمال، بدليل مقابلتها بكلمة (أخرت).
    \item \textbf{كفاية النص القرآني:} قول مجاهد في تفسير ﴿وَإِنَّهُ لَذِكْرٌ لَّكَ وَلِقَوْمِكَ﴾ بأنهم العرب قريش؛ عقّب عليه الشافعي بأن المعنى بيّنٌ ومستغنىً فيه بالتنزيل عن التفسير، دلالة على فصاحة القرآن وبيانه الذاتي.
    \item \textbf{معنى المجمل والبيان:} المجمل ما احتمل معنيين أو أكثر من غير ترجيح، والبيان هو الإيضاح والإظهار لما فيه خفاء.
    \item \textbf{معنى الشطر:} الشطر يأتي بمعنى الجهة في لغة العرب كقوله تعالى: ﴿فَوَلُّوا وُجُوهَكُمْ شَطْرَهُ﴾.
    \item \textbf{معنى العدل:} العدل ليس مجرد الصدق، بل هو من له ملكة تحمله على ملازمة المروءة وفعل الطاعات واجتناب الفسق وخوارم المروءة.
    \item \textbf{تعريف التقوى:} أحسن ما قيل فيها قول طلق بن حبيب: أن تعمل بطاعة الله على نور من الله ترجو ثوابه، وتترك المعصية على نور من الله تخاف عقابه.
    \item \textbf{سعة لغة العرب:} لغة العرب هي أوسع الألسنة وأكثرها ألفاظاً، ولا يحيط بها فرد إحاطة تامة إلا نبي، ولكنها محفوظة ومستوعبة في مجموع الأمة.
    \item \textbf{خصائص لسان العرب:} من سعة لغة العرب إطلاق الأسماء الكثيرة على المسمى الواحد (كالأسد والليث)، وإطلاق الاسم الواحد على معانٍ متعددة المشتركة لفظاً (كالعين للبصر والماء والذهب).
    \item \textbf{تفسير (الذين قال لهم الناس):} تكرر لفظ (الناس) في الآية ثلاث مرات، وفي كل مرة يُراد به فئة مخصوصة تختلف عن الأخرى (المُخْبِرون، والمُخْبَرون، والجَمْعُ الكافر المنصرف من أحد).
    \item \textbf{معنى القرية:} تُطلق القرية في القرآن ويراد بها الأبنية تارة، والسكان تارة أخرى، والسياق هو الذي يُعين المعنى المقصود كقوله: ﴿إِذْ يَعْدُونَ فِي السَّبْتِ﴾.
    \item \textbf{تفسير الحكمة:} إذا ذُكرت «الحكمة» في القرآن الكريم مقترنة بـ «الكتاب»، فالمراد بها سنة النبي صلى الله عليه وسلم بإجماع المحققين من أهل العلم.
    \item \textbf{تفسير ذوي القربى:} هم بنو هاشم وبنو المطلب خاصة دون غيرهم من بني عبد مناف، لقوله ﷺ: "إنما بنو هاشم وبنو المطلب شيء واحد في الجاهلية والإسلام".
    \item \textbf{إفراد فعل الطاعة مع الرسول:} تكرار فعل ﴿وَأَطِيعُوا﴾ مع الرسول في الآية يفيد وجوب الاستقلال بطاعته ﷺ فيما شرعه حتى لو لم ينص عليه القرآن، بخلاف أولي الأمر الذي لم يسبق بفعل أمر مستقل.
    \item \textbf{تعريف النسخ:} النسخ لغة هو الإزالة والنقل، واصطلاحاً هو رفع حكم شرعي متقدم بدليل شرعي متراخٍ عنه.
    \item \textbf{معنى الإحصان للإماء:} الإحصان المذكور في قوله تعالى ﴿فَإِذَا أُحْصِنَّ﴾ في حق الإماء المراد به (الإسلام) وليس الزواج، فالأمة المسلمة تجلد نصف المحصنة.
    \item \textbf{معاني الإحصان:} يُطلق الإحصان في النصوص الشرعية ويراد به أحد معانٍ أربعة: الإسلام، الزواج، الحرية، والعفة. وكل مانع يمنع من المحرم فهو إحصان.
    \item \textbf{الفرق بين الوَضوء والوُضوء:} الوَضوء (بفتح الواو) هو الماء الذي يُستخدم للطهارة، والوُضوء (بضم الواو) هو الفعل نفسه أي عملية غسل الأعضاء.
    \item \textbf{معنى الكلالة:} هو من مات وليس له أصل وارث (أب أو جد) ولا فرع وارث (ابن أو بنت).
    \item \textbf{تفسير الماعون:} فُسر الماعون في قوله تعالى ﴿وَيَمْنَعُونَ الْمَاعُونَ﴾ بأنه الزكاة المفروضة، وهو قول عدد من السلف كابن عباس.
    \item \textbf{معنى الخَرْص:} هو التقدير المعتمد على الخبرة لما على الأشجار من الثمار (كالنخيل والعنب) لتحديد مقدار الزكاة الواجبة فيها.
    \item \textbf{الفرق بين الدجر والعلس والقطنية:} أسماء لبعض المزروعات التي تقتات وتدخر فتجب فيها الزكاة (كالعدس والحمص).
    \item \textbf{معنى الدافة والودك والعسيف:} الدافة: القوم السائرون الذين قدموا ضيوفاً على المدينة. الودك: دسم اللحم وشحمه. العسيف: الأجير العامل.
    \item \textbf{معنى الإسفار والغلس:} الإسفار هو انكشاف الظلمة وظهور ضوء النهار، أو التيقن من طلوع الفجر الصادق. والغلس هو بقايا الظلمة المختلطة بتباشير الصباح الأولى.
    \item \textbf{معنى البيات (الغارة):} يُقصد به الهجوم على الأعداء ليلاً على غفلة، حيث لا يتمايز المقاتل عن غيره كالنساء والولدان.
    \item \textbf{معنى المزابنة والعرايا:} المزابنة بيع الثمر بالتمر كيلاً وهو محرم لما فيه من الغرر والجهالة، والعرايا رخصة مستثناة بخرصها لفقير لا يملك النقد، بحدود خمسة أوسق.
    \item \textbf{معنى اشتمال الصماء والتعريس:} اشتمال الصماء أن يلتف بثوب واحد يمنع إخراج يديه ويفضي لكشف عورته. والتعريس هو النزول للراحة والنوم في السفر.
    \item \textbf{معنى الفضيخ والمهراس والأريكة:} الفضيخ شراب يتخذ من البسر المطبوخ أو المشقوق. المهراس حجر منقور يدق فيه أو يتوضأ منه. والأريكة هي السرير أو ما يتكئ عليه الإنسان براحة.
    \item \textbf{معنى العاقلة والغرة والمخابرة:} العاقلة: عصبة الرجل وقرابته الذين يتحملون دفع الدية. الغرة: عبد أو أمة، وهو مقدار دية الجنين. المخابرة: المزارعة على حصة من المحصول.
    \item \textbf{معنى الخراج بالضمان والموضحة:} الخراج هو ما يغلّه العبد من ربح وعمل يستحقه المشتري لضمانه له. والموضحة هي الشجة التي تشق اللحم وتوضح العظم (ديتها 5 من الإبل).
    \item \textbf{تفسير معنى القرء:} يُطلق القرء لغة على الوقت، ورجّح الشافعي أن المراد به (الطهر) في العدة، لأن أصله بمعنى الحبس والجمع، والرحم في حالة الطهر يحبس الدم بخلاف الحيض الذي يُرخيه.
    \item \textbf{معنى المصراة:} الشاة التي يُحبس لبنها في ضرعها أياماً لإيهام المشتري بكثرة حليبها، وحكمها حق الرد مع صاع تمر بدلاً عن اللبن المحلوب.
    \item \textbf{معنى "كذب" في لغة الحجاز:} تُطلق في لغة أهل الحجاز بمعنى "أخطأ" أو "أخبر بخلاف الواقع دون تعمد للكذب"، كما قال ﷺ: "كذب أبو السنابل".
\end{itemize}

\section{الفوائد الأصولية والحديثية}
\begin{itemize}
    \item \textbf{أهمية كتاب الرسالة:} هو أول ما صُنف في أصول الفقه، وتميز بجمع التقعيد الأصولي والتأصيل اللغوي مع ضرب الأمثلة التطبيقية الفقهية.
    \item \textbf{حجية الحديث وعلومه:} لا يُرد الحديث بحجة عدم وجوده في القرآن، فالسنة شارحة ومبينة ومستقلة بالتشريع، ويجب فهمها في ضوء منهج السلف.
    \item \textbf{وجوه البيان عند الشافعي:} تتنوع بين نص قرآني لا يحتاج لبيان، ومجمل قرآني بينته السنة، وفرض قرآني قيدت السنة بعض أحكامه، وسنة مستقلة بتشريع لم ينص عليه القرآن.
    \item \textbf{تعريف القياس الصحيح:} هو إعمال للنصوص الشرعية والاجتهاد في تطبيقها كالاجتهاد في تحديد جهة القبلة، أو تحديد المثل في جزاء الصيد، وهو يختلف عن القياس الفاسد المصادم للنص.
    \item \textbf{أقسام القياس عند الشافعي:} قياس المعنى (العلة)، وهو أن يحرّم الله شيئاً لمعنى معقول فإذا وُجد في غيره أُلحق به. وقياس الشبه، وهو أن يتردد فرع بين أصلين فيلحق بأقربهما شبهاً.
    \item \textbf{شروط الاجتهاد:} يشترط فيمن يتصدر للاجتهاد والفتيا الإلمام الواسع بلغة العرب دلالةً وتطبيقاً، والعلم بمعاني الأوامر والنواهي، ومعرفة الناسخ والمنسوخ من النصوص.
    \item \textbf{الإحاطة بالسنة:} لا يحيط عالم واحد بالسنة كاملة، ولكنها لا تغيب عن مجموع العلماء. ومصادر الإحاطة بالسنة تشمل الكتب الستة، والمسانيد (كأحمد وأبي يعلى والبزار)، والمعاجم، والسنن.
    \item \textbf{تعريف العام والخاص:} العام هو اللفظ المستغرق لجميع ما وضع له بحسب وضع واحد دفعة، والتخصيص هو قصر العام على بعض مسمياته وإخراج بعض ما كان داخلاً تحته.
    \item \textbf{أقسام العام في القرآن:} ينقسم إلى عام باقٍ على عمومه (كخلق الله لكل شيء)، وعام يراد به الخصوص (كقوله: الذين قال لهم الناس)، وعام يجمع بين العموم والخصوص في آية واحدة.
    \item \textbf{قرائن تحديد المعنى:} دلالات الألفاظ تختلف باختلاف السياق والتركيب والإضافة، فالسياق هو المحدِّد لمراد المتكلم عند احتمال اللفظ لأكثر من معنى.
    \item \textbf{تخصيص السنة للقرآن:} السنة تقيد مطلق القرآن وتخصص عامه؛ كاستثناء القاتل والمخالف في الدين من عموم الإرث، وتخصيص حد السرقة بربع دينار فصاعداً، وحد الزنا بالبكر، وجواز المسح على الخفين من عموم آية غسل القدمين.
    \item \textbf{التفريق بين أنواع القياس:} القياس الصحيح هو الذي يُعمل النص ولا يصادمه (كقياس تحديد القبلة)، والقياس الفاسد هو المصادم للنص (كقياس إبليس في الامتناع عن السجود).
    \item \textbf{مذهب الشافعي في النسخ:} القرآن يُنسخ بالقرآن، والسنة بالسنة. ولا تنسخ السنة القرآن أبداً بل تبينه، والنسخ لا يكون إلا إلى بدل.
    \item \textbf{حكم التنصيف في الحدود:} التنصيف يقع في الحدود التي تقبل العدد كالجلد (يُنصف للإماء)، أما الحدود المتلفة للنفس كالرجم فلا تقبل التنصيف ولذلك لا تُرجم الأمة.
    \item \textbf{الفرق بين المغمى عليه والسكران:} المغمى عليه فاقد العقل بغير اختياره فلا قضاء عليه كالحائض، أما السكران فأدخل المانع باختياره فيلزمه القضاء تبكيتاً وعقوبة له.
    \item \textbf{حقيقة الإجماع مع النسخ:} الإجماع لا يَنْسَخ ولا يُنْسَخ لانقطاع الوحي بوفاة النبي ﷺ، وإنما يُستدل بالإجماع على وجود نص أوجب النسخ.
    \item \textbf{تقييد الوصية بالسنة:} خصصت السنة عموم آيات الوصية في القرآن بأمرين: إخراج الوارث منها (لا وصية لوارث)، وتقييد مقدارها بالثلث فما دونه لقوله ﷺ (الثلث والثلث كثير).
    \item \textbf{قبول أخبار الآحاد المدعومة بالعمل:} الحديث المنقطع أو الذي فيه ضعف يسير إذا اعتضد وتلقاه أهل العلم والمغازي بالقبول والتطبيق (كحديث لا وصية لوارث) يُعمل به ويُحتج به لتخصيص القرآن.
    \item \textbf{أدنى الفرض وما زاد نافلة:} فعل النبي ﷺ للوضوء مرة واحدة دل على أنه الحد الأدنى المجزئ للفرض، وفعله له ثلاثاً دل على أن الزيادة نافلة ومستحبة لطلب الفضل.
    \item \textbf{تخصيص السنة للبيوع:} عموم إحلال البيع في القرآن ﴿وَأَحَلَّ اللَّهُ الْبَيْعَ﴾ خُصص بالسنة التي حرمت بيوعاً محددة (كالمبادلة بالذهب مع دفع الفرق) لعلة الربا أو الغرر وإن تراضى العاقدان.
    \item \textbf{النسخ في صلاة الخوف:} تأخير الصلاة بسبب الخوف (كما حدث يوم الخندق) نُسخ بنزول أحكام صلاة الخوف، فلا يجوز تأخير الصلاة عن وقتها بحال، وتُصلى على أي هيئة أمكنت.
    \item \textbf{تخصيص السنة لعموم المأكولات والمحرمات:} السنة تستقل بالتشريع كما خصصت تحريم (كل ذي ناب من السباع) زيادة على ما في القرآن، وخصصت تحريم (الجمع بين المرأة وعمتها/خالتها).
    \item \textbf{تخصيص عدة الوفاة للحامل:} السنة بينت أن الحامل المتوفى عنها زوجها عدتها بوضع الحمل مطلقاً، حتى لو وضعت بعد الوفاة بساعات، ولا تجمع بين العدتين (وضع الحمل والأشهر الأربعة).
    \item \textbf{أسباب الاختلاف الظاهري للأحاديث:} ترجع إلى أسباب منها: سماع الراوي للجواب دون السؤال، أو حفظه للمنسوخ دون الناسخ، أو اختلاف الحالين اللذين سن فيهما النبي ﷺ كحالة العزيمة وحالة الرخصة.
    \item \textbf{نسخ السنة للسنة:} كما أن القرآن ينسخ بعضه بعضاً، فإن السنة ينسخ بعضها بعضاً كنسخ النهي عن زيارة القبور بالإباحة، ولا تضيع سنة ناسخة عن مجموع الأمة أبداً.
    \item \textbf{فقه بيوع العرايا:} هي استثناء ورخصة من بيوع الغرر والمزابنة، رخص فيها النبي ﷺ لحاجة الفقراء أن يشتروا الرطب بالتمر في حدود خمسة أوسق، ولا يقاس عليها للتحريم الأصلي.
    \item \textbf{تصرف النبي ﷺ بالإمامة ومراعاة المصلحة:} بعض أوامر النبي ﷺ ونواهيه (كالنهي عن ادخار الأضاحي لأجل الدافة) لم تكن تشريعاً عاماً، بل تصرفاً سياسياً مصلحياً بوصفه إماماً للمسلمين يزول بزوال علته.
    \item \textbf{التنصيف يقع على القابل للتجزئة:} الحدود المتلفة للنفس (كالرجم) لا تقبل التنصيف، لذا لا تُرجم الأمة المذنبة، بل يُنصّف في حقها الحد القابل للعدد وهو (الجلد) فتُجلد خمسين جلدة.
    \item \textbf{اختلاف التنوع في العبادات:} ما ورد باختلاف في الصيغ (كأنواع التشهد والقراءات القرآنية) يُحمل على التوسعة والتنوع، وكلها مجزئة وصحيحة ولا يُضرب بعضها ببعض.
    \item \textbf{قواعد الترجيح عند التعارض:} إذا تعذر الجمع يُصار إلى الترجيح بناءً على: كثرة الرواة وأقدميتهم في الصحبة، وموافقة ظاهر القرآن، وموافقة القواعد والأصول، وعمل الخلفاء الراشدين.
    \item \textbf{التوفيق بحمل الحديث على اختلاف الحال:} الجمع بين النهي المطلق لاستقبال القبلة للغائط والبول مع فعل النبي ﷺ يُحمل على اختلاف الحال؛ فالنهي للصحراء والفضاء، والرخصة للبنيان لضيقها ووجود الستر.
    \item \textbf{ترجيح أحاديث الربا:} قدم الشافعي والجمهور أحاديث تحريم ربا الفضل على حديث أسامة (إنما الربا في النسيئة) لكثرة الرواة وكونهم أحفظ، مع إمكانية توجيه الأخير ببيع الأصناف المختلفة.
    \item \textbf{فقه الجمع في قتل النساء والولدان:} النهي عن قتلهم محمول على التعمد والقصد حال تمايزهم. والإباحة المحمولة في حديث "هم منهم" تكون في حال اختلاطهم ليلاً في "البيات" وعدم إمكانية التمييز.
    \item \textbf{استنباط الوجوب الاختياري (المستحب):} فهم الشافعي أن غسل الجمعة مستحب وليس شرطاً لصحة الصلاة، بدليل إقرار عمر لعثمان حين دخل واكتفى بالوضوء ولم يأمره بالخروج ليعيد الغسل.
    \item \textbf{تقييد نهي الخطبة على خطبة الأخ:} خصص بحديث فاطمة بنت قيس أن النهي لا يقع بمجرد التقدم، بل يقع بعد (الركون والقبول) من المرأة للرجل الأول، فما لم تقبل جاز لغيره التقدم.
    \item \textbf{تخصيص أوقات النهي:} النهي عن الصلاة في الأوقات المكروهة يخص (النوافل المطلقة)، وتُستثنى منها الفرائض المقضية والصلوات ذوات الأسباب (كركعتي الطواف) لورود مخصصات لها في السنة.
    \item \textbf{دلالة النهي وتنوعه:} الأصل في نهي الشارع أنه للتحريم ويقتضي الفساد (كنكاح الشغار والمتعة)، وقد يصرف إلى الكراهة أو الإرشاد بقرائن (كالنهي عن الأكل من وسط الصحفة، أو التعريس بالطريق).
    \item \textbf{حقيقة فرض الكفاية:} هو ما قُصد به إيجاد الفعل بقطع النظر عن الفاعل، فإذا قام به من يكفي سقط الإثم عن الباقين، كصلاة الجنازة والجهاد ورد السلام والتفقه.
    \item \textbf{شروط الحديث الصحيح:} سبقت عبارات الشافعي علماء المصطلح في التأصيل لشروط الحديث الصحيح (العدالة، الضبط، اتصال السند، عدم الشذوذ، عدم العلة/التدليس).
    \item \textbf{الفرق بين الشهادة والرواية:} تفترق الرواية عن الشهادة في جواز قبول الواحد، والمرأة المنفردة، وقبول العنعنة (لمن سلم من التدليس)، لأن الرواية تشريع عام يتطلب الحفظ والضبط أكثر من مجرد الإخبار بمشاهدة لواقعة.
    \item \textbf{حكم عنعنة المدلس:} من عُرف بالتدليس (رواية عمن لقيه ما لم يسمع منه) لا يُقبل حديثه إن قال "عن"، حتى يصرّح بالتحديث أو السماع، والتدليس يُعد عورة في النقل وليس كذباً محضاً.
    \item \textbf{التحديث عن بني إسرائيل:} إباحة التحديث عن بني إسرائيل (وحدثوا عن بني إسرائيل ولا حرج) هي رخصة للنقل للاتعاظ والاعتبار، وليست رخصة لتعمد نقل الكذب، بخلاف الرواية عن النبي ﷺ التي يُشترط فيها صحة الإسناد.
    \item \textbf{حجية خبر الآحاد في العقائد والأحكام:} إرسال النبي ﷺ للآحاد (كمعاذ) لتبليغ التوحيد والفرائض يثبت أن خبر الآحاد حجة ملزمة في العقيدة والشريعة على حد سواء وتنقض به الأحكام المتيقنة.
    \item \textbf{شروط قبول الحديث المرسل:} لا يُقبل المرسل إلا من كبار التابعين وبشروط: أن يُعتضد بمسند، أو مرسل من طريق آخر، أو قول صحابي، أو فتوى الجمهور، مع كون المرسِل لا يروي عن الضعفاء.
    \item \textbf{توجيه التثبت في الرواية:} طلب عمر لشاهد ثانٍ مع أبي موسى في حديث الاستئذان كان احتياطاً لسد الباب على الكاذبين، وليس رداً لحجية خبر الآحاد.
    \item \textbf{القياس والاجتهاد:} هما اسمان لمعنى واحد عند الشافعي، وهو بذل الجهد للوصول إلى الحق وإعمال النصوص عبر القرائن والدلائل فيما ليس فيه نص قاطع بعينه.
    \item \textbf{أجر الاجتهاد على الظاهر:} المجتهد إن وافق الحق عند الله باطناً نال أجرين، وإن وافق ظاهر الأدلة وأخطأ الحق الباطن فله أجر واحد على اجتهاده، بشرط أن يكون من أهل الاجتهاد.
    \item \textbf{شروط المجتهد:} حصرها الشافعي في الإحاطة بعلوم الكتاب (ناسخه ومنسوخه، وعامه وخاصه)، والسنن، وأقوال السلف إجماعاً واختلافاً، ولسان العرب، وصحة العقل والتمييز.
    \item \textbf{قياس الأولى (الأقوى):} أن يُحرم الشرع القليل من الشيء، فيُعلم يقيناً أن الكثير منه أشد تحريماً، أو يُحمد على اليسير من الطاعة فيكون الكثير أولى.
    \item \textbf{قاعدة الخراج بالضمان:} استنبط الشافعي منها أن كل ما ينتج عن المبيع في فترة ضمان المشتري (كغلة العبد وثمر الشجر) فهو حق له، ولا يرده إذا رد الأصل بعيب.
    \item \textbf{تحمل العاقلة للدية:} العاقلة (عصبة الجاني) تتحمل دية الخطأ كاملة، والراجح عند الشافعي أنها تتحمل أيضاً ما دون الثلث قياساً على تحملها للأكثر.
    \item \textbf{قصر ضمان العاقلة:} العاقلة لا تتحمل إلا جناية الخطأ فقط. أما جناية العمد، والكفارات، والزكاة، ومسؤولية حفظ الأموال، فكلها تقع في مال المكلف الخاص.
    \item \textbf{تفريد دية الجنين:} دية الجنين الميت قدرها ثابت (الغرة) تعبدياً، ولا يفرّق فيها بين الذكر والأنثى بخلاف الجنين لو نزل حياً ثم مات فإن ديته تختلف باختلاف جنسه.
    \item \textbf{تخصيص عدة الوفاة بوضع الحمل:} حديث سبيعة الأسلمية مخصص لعموم آية التربص أربعة أشهر وعشرا، فالحامل تنقضي عدتها بمجرد الوضع ولو بعد ساعات من الوفاة.
    \item \textbf{أحكام الإيلاء:} المُولي (الحالف على ترك جماع زوجته) لا يقع طلاقه بمجرد انقضاء الأشهر الأربعة تلقائياً، بل يُوقفه القاضي ويُخيره بين الرجوع أو إيقاع الطلاق لرفع الضرر.
\end{itemize}

\backmatter

\renewcommand{\bibname}{قائمة المصادر والمراجع}
\begin{thebibliography}{99}
    \bibitem{risala_shakir} الشافعي، محمد بن إدريس. \textit{الرسالة}. تحقيق الشيخ أحمد محمد شاكر. القاهرة: مطبعة مصطفى البابي الحلبي، 1940م.
    \bibitem{umm_rifat} الشافعي، محمد بن إدريس. \textit{الأم}. تحقيق د. رفعت فوزي عبد المطلب. المنصورة: دار الوفاء، 2001م.
    \bibitem{tabari} الطبري، محمد بن جرير. \textit{جامع البيان عن تأويل آي القرآن}. تحقيق محمود شاكر.
    \bibitem{khattabi} الخطابي، حمد بن محمد. \textit{معالم السنن}. مطبعة محمد علي صبيح، 1932م.
\end{thebibliography}

\printindex[quran]
\printindex[hadith]
\printindex[general]

\end{document}